\documentclass[11pt,openany]{book}
\usepackage[T1]{fontenc}
\usepackage[utf8]{inputenc}
\usepackage[expansion=false]{microtype}
\usepackage{amsmath,amssymb,amsthm,mathtools,stmaryrd}
\usepackage[margin=1.15in]{geometry}
\usepackage{booktabs}
\usepackage{enumitem}
\usepackage[hidelinks]{hyperref}

\setlength{\emergencystretch}{3em}
\newtheorem{theorem}{Theorem}[chapter]
\newtheorem{lemma}[theorem]{Lemma}
\newtheorem{proposition}[theorem]{Proposition}
\newtheorem{definition}[theorem]{Definition}
\newtheorem{corollary}[theorem]{Corollary}
\newtheorem{hypothesis}[theorem]{Hypothesis}
\theoremstyle{remark}
\newtheorem{remark}[theorem]{Remark}
\newtheorem{example}[theorem]{Example}

\newcommand{\Pop}{\mathsf{Pop}}
\newcommand{\Refuse}{\mathsf{Refuse}}
\newcommand{\Bind}{\mathsf{Bind}}
\newcommand{\Transform}{\mathsf{Transform}}
\newcommand{\Verify}{\mathsf{Verify}}
\newcommand{\Collapse}{\mathsf{Collapse}}
\newcommand{\Compensates}{\mathsf{Compensates}}
\newcommand{\Commit}{\operatorname{Commit}}
\newcommand{\dom}{\operatorname{dom}}
\newcommand{\fn}{\operatorname{fn}}
\newcommand{\record}{\operatorname{record}}
\newcommand{\replay}{\operatorname{replay}}
\newcommand{\valid}{\operatorname{valid}}
\newcommand{\subj}{\operatorname{subj}}
\newcommand{\vf}{\mathsf{vf}}
\newcommand{\plan}[1]{\par\medskip\noindent\textbf{Chapter plan.} #1\par}
\newcommand{\spine}[1]{\par\smallskip\noindent\textbf{Spine results.} #1\par}

\title{\textbf{The Calculus of Disposition}\\[0.6em]
\large Processes, Types, and Witnessed Computation}
\author{Flyxion\\Independent Researcher}
\date{September 2026\\[1em]\small Working draft: Preface and eleven chapters drafted,\\ remaining chapters in plan form}

\begin{document}
\frontmatter
\maketitle
\tableofcontents

%===================================================================
\chapter{Preface}
%===================================================================

Every formal calculus of computation decides what a term is, what may
happen to it, when two results count as the same, and what an observer is
allowed to distinguish. These decisions are often introduced separately.
They are not independent. A calculus that treats computation as rewriting
will naturally ask which rewritten terms have the same normal form. A
calculus that treats computation as communication will instead ask which
interactions an environment can observe. A calculus that retains a history
must ask a further question: which differences may be ignored without
asserting that recorded events never occurred?

The calculus developed in this book begins from the last question. Its
basic subject is not a value alone, nor a process considered apart from its
past, but a commitment moving through a witnessed history. A commitment
may be bound, transformed, verified, refused, or collapsed into an
observation. Each such action changes what continuations remain admissible.
None is permitted to revise the past that made it possible.

This orientation yields a system with three formal layers and one durable
semantic object. A lambda layer computes values, predicates,
transformations, reasons, and witnesses. A process layer organizes fresh
commitments, interaction, choice, concurrency, and disposition. A type
algebra separates reusable knowledge from linear rights to advance a
commitment. An append-only ledger records the events from which the current
disposition state can be reconstructed. The layers are kept distinct on
purpose. Lambda terms compute values; processes consume and produce
commitments, capabilities, evidence, and observations. A transformation has
a functional type such as $f:A\multimap B$, while a process applies that
transformation within the history of one particular commitment.

The principal theorem is therefore not merely that well-typed terms do not
go wrong. It is that a commitment reaches a terminal disposition only by an
authorized, causally complete, and durably witnessed history. We call this
property \emph{typed disposition safety}. The proof is cumulative:
\begin{center}
\begin{tabular}{c}
parsing and $\alpha$-invariance\\
$\Downarrow$\\
lifecycle, freshness, and history preservation\\
$\Downarrow$\\
authority, binding, and witness ancestry\\
$\Downarrow$\\
collapse soundness and terminal exclusivity\\
$\Downarrow$\\
replay and compensation\\
$\Downarrow$\\
\textbf{typed disposition safety}
\end{tabular}
\end{center}
The final proof is short because the book is arranged so that nearly all of
its content has already been proved locally. In the present calculus most
of the untyped spine turns out to rest on a single invariant, the
\emph{ledger path lemma} of Chapter~\ref{ch:path}: the records of each
commitment, read in order, trace a legal path through its lifecycle ending
at its current state. Binding integrity, witness ancestry, collapse
soundness, terminal exclusivity, and refusal preservation are then
corollaries rather than separate theorems. The type algebra of Part~IV
strengthens the spine in its final form: commitment uniqueness from
linearity, lifecycle monotonicity from typestate preservation, and the
impossibility of a premature collapse from the input type of
$\mathsf{collapse}$.

\section*{Two layers, one vocabulary}

A reader who knows the canonical Spherepop kernel will notice that the
disposition calculus uses its operator names with more specific meanings.
In the kernel, $\Pop$ commits to one option among those available;
$\Refuse$ records that an option is inadmissible without removing it from
the history; $\Bind$ couples two elements without identifying them; and
$\Collapse_c$ observes a history under a rule $c$. In the disposition
calculus, $\mathsf{pop}$ introduces a fresh commitment; $\mathsf{refuse}$
is one of two terminal dispositions, available from every live state;
$\mathsf{bind}$ associates a commitment with a value or capability;
$\mathsf{collapse}$ produces an observation from a verified commitment
under a declared rule $k$; and $\mathsf{transform}$ and $\mathsf{verify}$
appear as lifecycle stages.

The book's claim is that the disposition calculus is a discipline
\emph{over} the kernel rather than a separate calculus, and it is a claim
to be proved, not assumed. Chapter~\ref{ch:layers} gives an explicit
translation of disposition records into kernel histories and proves it in
both directions at the level of ledgers: every disposition step becomes an
admissible kernel step, and every kernel step of translated shape whose
guard holds comes from a disposition step. Under that result the four
kernel operators remain the causal basis; transformation, verification,
and discharge are guarded compositions of them; the kernel itself enforces
freshness and finality; and everything else the lifecycle requires lives in
the guard. What the result does not establish is that the kernel performs
the computations the guards check. That stronger, transparent reading is
stated as a conjecture.

The two refusals agree on the property the whole corpus cares about: a
refusal is recorded rather than erased. They differ in scope. Kernel
refusal concerns an option and the history continues around it.
Disposition refusal concerns a commitment's entire future, and a later
reconsideration must create a distinct successor commitment rather than
reopen the refused one.

\section*{Three comparative claims}

Alongside the construction, the book tests three claims about the
calculus's position, each of which could turn out false. Each is stated
with what would settle it.

The first is that the calculus's nearest formal relative is the theory of
event structures with conflict rather than the $\pi$-calculus or the
lambda calculus. Every disposition world has a labelled event structure
underneath it, together with annotations $\eta$ carrying principals,
reasons, witnesses, certificates, versions, and rules
(Chapter~\ref{ch:worlds}). The claim is settled positively if the
admissibility conditions can be expressed as structural constraints on
such a structure over a finite alphabet, and negatively, with $\eta$ as the
identified residue, if some condition depends irreducibly on evidence
predicates.

The second concerns refusal. In its defensible form it says that recorded
refusal is not represented by the standard failures model of CSP without
adding an internal event or instrumentation. A CSP refusal set is a
property of behaviour inferred from outside a process; a recorded refusal
is authorized, reasoned, and witnessed data inside it. The stronger claim,
that no counterpart exists in any form, would need a separation result,
and until one is proved the book asserts only the weaker statement. It is
settled by an encoding of recorded refusal into instrumented CSP together
with a proof that uninstrumented failures equivalence identifies
histories that differ only in their recorded refusals.

The third is that making collapse rule-indexed changes the theory of
equivalence rather than relocating it. The observable label of a collapse
is the pair $\langle k,o\rangle$, so equal payloads under different rules
are distinct. The claim is settled by defining equivalence of processes
relative to a rule, $P\approx_k Q$, and determining when it is a
congruence, which classical equivalences arise as rules, and whether
simultaneous observation of one history under several rules yields a
result that van Glabbeek's linear-time/branching-time spectrum does not
already state. If it yields none, the spectrum contains the substance, and
Part~VII says so.

One early piece of evidence bears on the first claim. Chapter~\ref{ch:replay}
proves that replay is invariant under exchanging adjacent records of
distinct commitments, so that a ledger determines its disposition state
only up to a Mazurkiewicz commutation class. The causal order that the
ancestry theorems need is therefore the dependence order of a trace, not
the raw order of the ledger, and that is exactly the order an event
structure makes primary.

\section*{What counts as a result}

Positive comparative claims are stated relative to Gorla's encodability
criteria and the standard of full abstraction; negative claims follow the
model of Palamidessi's separation results. Spine theorems are stated with
their hypotheses visible. The liveness result in particular holds only
under fairness, responsive external dependencies, terminating
transformations, and decidable verification, and is stated with all four.
Three grades of assertion are kept apart: results proved here; standard
results cited and used without re-derivation, marked where load-bearing;
and conjectures or programmatic claims, collected in
Appendix~\ref{app:open}.

The book also uses three levels of validation, and they are not
interchangeable:
\[
\text{exhaustive finite checks}\;\subset\;
\text{untyped semantic proofs}\;\subset\;
\text{typed universal guarantees}.
\]
The laboratory of Chapter~\ref{ch:worlds} enumerates every admissible
ledger up to a bound and finds countermodels and specification errors; it
proves nothing about larger worlds. The untyped proofs of Parts~III, V,
and VI establish the behaviour of every operational derivation. The typed
metatheory of Part~IV establishes that well-formed programs cannot even
express several classes of duplicated, abandoned, or misindexed
continuation. Each level catches things the next cannot state, and none
substitutes for the one above it.

\section*{Status of this edition}

Drafted in full: Chapters~\ref{ch:past} (Computation with a Past),
\ref{ch:worlds} (Disposition Worlds), \ref{ch:layers} (The Kernel and the
Disposition Layer), \ref{ch:syntax} (Syntax and Binding), \ref{ch:lts}
(Labelled Disposition Semantics), \ref{ch:types} (The Type Algebra),
\ref{ch:session} (Preservation, Progress, and Protocol Fidelity),
\ref{ch:path} through \ref{ch:compensation} (Part~V), \ref{ch:replay}
(Replay), \ref{ch:progress} (Eventual Disposition), and
\ref{ch:typedsafety} (Typed Disposition Safety), \ref{ch:parsing}
(Concrete Syntax and Unique Parsing), and \ref{ch:authority} (Authority as
a Modality). The remaining chapters, chiefly the comparative ones, appear
as plans stating what each must establish.

%===================================================================
\mainmatter
\part{Orientation}
%===================================================================

\chapter{Computation with a Past}
\label{ch:past}

\section{Four commitments of a calculus}

A calculus is sometimes presented as syntax followed by reduction rules.
That description omits two choices which are equally fundamental. Besides
its terms and steps, a calculus fixes an equivalence and an observational
boundary. The untyped lambda calculus, for example, begins with variables,
abstractions, and applications; advances by beta reduction; compares terms
through a selected convertibility or observational relation; and exposes a
result through a normal form, a head form, or interaction with a context.
Process calculi alter every component. Their terms denote persisting
agents, their transitions carry labels, their equivalences compare
behaviours, and their observers distinguish actions rather than returned
values alone.

The coupling matters. If an observer sees only final values, two histories
may be identified even when one passed through a refusal and the other did
not. If refusal is itself observable, that identification becomes invalid.
If an append-only record survives the computation, equivalence cannot mean
literal equality of histories unless no abstraction is permitted. It must
instead mean equality under a declared observation rule.

\begin{definition}[Disposition]
A \emph{disposition} is an attributable and recordable determination of
the admissible future of a commitment. A disposition may advance the
commitment, refuse its continuation, or produce a terminal observation.
\end{definition}

This definition separates refusal from failure. A machine failure may
prevent the production of any disposition. A refusal is already a result:
it identifies a commitment, an authorized source, a reason or predicate,
and a record which persists after the refusal has become historical.

\section{State and history}

The running configuration is written
\[
  C=\langle P,S,A,O,H\rangle .
\]
Here $P$ is a process, $S$ is mutable computational state, $A$ maps
commitment identities to their current disposition states, $O$ contains
observations, and $H$ is an append-only history. The components are kept
separate because they answer different questions. The store says what is
currently available to computation. The commitment map says which
lifecycle rights are live. The observation set says what has been made
public under a collapse rule. The history says how those facts arose.

Two executions may therefore reach extensionally equal stores while
remaining semantically different:
\[
  S_1=S_2 \quad\text{but}\quad H_1\ne H_2.
\]
This inequality is not administrative residue. If $H_1$ contains a refusal
and $H_2$ does not, a later observer may legitimately distinguish them.

\begin{definition}[History extension]
For histories $H$ and $H'$, write $H\sqsubseteq H'$ when $H$ is a prefix of
$H'$. A transition system is \emph{historically monotone} when
\[
  C\longrightarrow C' \quad\Longrightarrow\quad H\sqsubseteq H'.
\]
\end{definition}

Historical monotonicity does not prohibit correction. It prohibits
correction by erasure. A later event may compensate for an earlier event,
but the compensation is intelligible only because both remain available.

\section{The lifecycle}

A commitment has an identity $c$, content of type $A$, and a lifecycle
state. At the untyped operational level we use
\[
\begin{aligned}
d ::= {}& \mathsf{open}(q)
\mid \mathsf{bound}(q,b)
\mid \mathsf{transformed}(b,v)\\
&\mid \mathsf{verified}(v,w)
\mid \mathsf{verificationFailed}(v,\varphi)\\
&\mid \mathsf{refused}(r)
\mid \mathsf{collapsed}(o).
\end{aligned}
\]
The live states are $\mathsf{open}$, $\mathsf{bound}$,
$\mathsf{transformed}$, $\mathsf{verified}$, and
$\mathsf{verificationFailed}$. The terminal states are $\mathsf{refused}$
and $\mathsf{collapsed}$.

Verification has two outcomes and two successor states. A positive result
leads to $\mathsf{verified}$, carrying a witness. A negative result leads to
$\mathsf{verificationFailed}$, carrying an authenticated certificate of
the failure. The name is chosen with care: nothing has failed to execute.
Verification has returned a certified negative result. The two branches are
symmetric: successful verification does not collapse the commitment, and
failed verification does not refuse it. Verification classifies evidence;
an authority disposes of the continuation.

Refusal is permitted from every live state. This is essential. A reviewer
may reject a binding, a transformation may be found inadmissible, or
verification may have failed. If refusal were available only from
$\mathsf{open}$, these outcomes would become unrecorded dead ends rather
than dispositions.

Let $\preceq_D$ be the least reflexive and transitive relation containing
the constructive edges
\[
\mathsf{open}\prec_D\mathsf{bound}
\prec_D\mathsf{transformed}
\prec_D\mathsf{verified}
\prec_D\mathsf{collapsed},
\]
\[
\mathsf{transformed}\prec_D\mathsf{verificationFailed},
\]
and a refusal edge from each live state to $\mathsf{refused}$. From
$\mathsf{verificationFailed}$ the refusal edge is the only one. The two
terminal states are incomparable. Consequently, no ordinary lifecycle
transition can turn a refusal into a collapse or a collapse into a refusal.

The lifecycle is small enough to tabulate. For each state $s$, let
$\mathsf{Next}(s)$ be the set of actions that may follow it:
\[
\begin{aligned}
\mathsf{Next}(\mathsf{open})&=\{\mathsf{bind},\mathsf{refuse}\},&
\mathsf{Next}(\mathsf{bound})&=\{\mathsf{transform},\mathsf{refuse}\},\\
\mathsf{Next}(\mathsf{transformed})&=\{\mathsf{verify},\mathsf{refuse}\},&
\mathsf{Next}(\mathsf{verified})&=\{\mathsf{collapse},\mathsf{refuse}\},\\
\mathsf{Next}(\mathsf{verificationFailed})&=\{\mathsf{refuse}\},&
\mathsf{Next}(\mathsf{refused})&=\mathsf{Next}(\mathsf{collapsed})=\varnothing,
\end{aligned}
\]
This \emph{refusal completion table} is the single source from which the
rest of the book's descriptions of the lifecycle are derived: the
operational rules of Chapter~\ref{ch:lts}, the lifecycle step function of
Chapter~\ref{ch:path}, which is the table with the data carried by each
state added, and the session type of Chapter~\ref{ch:session}, which is the table read
as a protocol. Every row for a live state contains $\mathsf{refuse}$; every
terminal row is empty. Those two facts are, in miniature, the statements
that refusal is always available and that terminal dispositions are
final. The $\mathsf{verify}$ entry has two outcomes, to $\mathsf{verified}$
and to $\mathsf{verificationFailed}$. A policy may extend the failed row,
for instance with a fresh verification attempt or replacement evidence,
but nothing is added to it implicitly.

\section{Records and the step function}
\label{sec:records}

Every disposition, and every discharge of an external obligation, leaves
a record in the history. The records are
\[
\begin{gathered}
\Pop(c,q),\qquad \Compensates(a,c,q,d),\qquad \Bind(a,c,b),\qquad
\Transform(c,f,b,v),\\
\Verify(a,c,v,w),\qquad \mathsf{VerificationFailed}(a,c,v,\varphi),\\
\Refuse(a,c,r),\qquad
\Collapse(a,c,k,w,o),\qquad
\mathsf{Discharged}(a,c,\omega,\varphi),\qquad
\mathsf{TimedOut}(a,c,\omega,t,\varphi),
\end{gathered}
\]
where $a$ names a principal, $q$ a request, $b$ and $v$ values, $f$ a
transformation, $w$ a witness, $k$ an observation rule, and $o$ an
observation.

A verifier $g$ applied to a commitment and value returns one of two
results: $\mathsf{ok}(w)$, a witness, or $\mathsf{fail}(\varphi)$, a
\emph{failure certificate}. Certificates have a fixed shape,
\[
\varphi=\langle\kappa,g,c,h(v),\mathsf{outcome},\sigma\rangle,
\]
naming the verifier's semantic version $\kappa$, the issuer $g$, the
commitment $c$, a hash $h(v)$ of the value or claim the certificate is
about, the outcome it attests, and a signature $\sigma$ over all of
these. The same shape serves for external evidence of discharge, below,
with the obligation in place of the value. A negative verification is
recorded by its own record, $\mathsf{VerificationFailed}(a,c,v,\varphi)$,
carrying the certificate. Refusal reasons $r$ are always written by the
refusing principal; there is no reserved reason, because a failed
verification is not a refusal.

The \emph{subject} of a record is the commitment whose state it changes.
For a compensation record it is the new successor $c$; the predecessor $d$
is referred to but not changed. In every other record it is the
commitment named first after the principal, or first overall for $\Pop$
and $\Transform$. For a history $H$, the \emph{projection} $H|_c$ is the
subsequence of records whose subject is $c$, in ledger order.

Two predicates on evidence are fixed: $\valid(w,c,v)$, which says $w$ is a
correct witness for $c$ having value $v$, and
\[
\mathsf{checkCert}(\varphi,c,x,\mathsf{outcome}),
\]
which says that $\varphi$ is well formed, that its commitment field is $c$,
its claim field is $h(x)$, its outcome field is the one named, and its
signature verifies for the issuer and version it names. The check reads
the certificate; it never reruns the verifier. Rerunning may be
impossible, expensive, nondeterministic, or altered by a later version,
and replay must remain deterministic when the original verifier is gone.

\begin{definition}[Lifecycle step]\label{def:delta}
The partial function $\delta$ takes a state, or $\bot$ for ``not yet
introduced'', and a record, and returns the resulting state:
\[
\begin{aligned}
\delta(\bot,\Pop(c,q))&=\mathsf{open}(q),\\
\delta(\bot,\Compensates(a,c,q,d))&=\mathsf{open}(q),\\
\delta(\mathsf{open}(q),\Bind(a,c,b))&=\mathsf{bound}(q,b),\\
\delta(\mathsf{bound}(q,b),\Transform(c,f,b,v))&=\mathsf{transformed}(b,v),\\
\delta(\mathsf{transformed}(b,v),\Verify(a,c,v,w))&=\mathsf{verified}(v,w)
  &&\text{if }\valid(w,c,v),\\
\delta(d,\Refuse(a,c,r))&=\mathsf{refused}(r)
  &&\text{if }\mathsf{live}(d),\\
\delta(\mathsf{transformed}(b,v),\mathsf{VF}(a,c,v,\varphi))
  &=\mathsf{vFailed}(v,\varphi)
  &&\text{if }\mathsf{checkCert}(\varphi,c,v,\mathsf{fail}),\\
\delta(d,\mathsf{Discharged}(a,c,\omega,\varphi))&=d
  &&\text{if }\mathsf{live}(d)\text{ and}\\
  &&&\quad\mathsf{checkCert}(\varphi,c,\omega,\mathsf{discharged}),\\
\delta(d,\mathsf{TimedOut}(a,c,\omega,t,\varphi))&=d
  &&\text{if }\mathsf{live}(d)\text{ and}\\
  &&&\quad\mathsf{checkCert}(\varphi,c,\langle\omega,t\rangle,\mathsf{expired}),\\
\delta(\mathsf{verified}(v,w),\Collapse(a,c,k,w,o))&=\mathsf{collapsed}(o)
  &&\text{if }o=\mathsf{observe}_k(v),
\end{aligned}
\]
where $\mathsf{VF}$ abbreviates the record $\mathsf{VerificationFailed}$
and $\mathsf{vFailed}$ the state $\mathsf{verificationFailed}$, and $\delta$
is undefined in every other case. Repeated variables must match: the
$b$ of a transformation record must be the $b$ held in the bound state,
the $v$ of a verification record the $v$ produced by transformation, and
the $w$ of a collapse record the witness held in the verified state. Write
$\delta^{*}$ for the left fold of $\delta$ over a sequence from $\bot$.
\end{definition}

The step function is the refusal completion table with the data carried
by each state checked against the data in each record. Discharge and
timeout records are the two clauses that leave the lifecycle state
unchanged. A discharge records that principal $a$ discharged an external
obligation $\omega$ that $c$ was waiting on, with authenticated evidence
$\varphi$. A timeout records, with evidence from a clock authority $a$,
that $\omega$ was \emph{not} discharged within the declared interval $t$.
Neither ends the commitment; both are consumed and checked by the fold, not
skipped. Formally $\delta$ also carries, beside each state, the set of
obligations already settled for $c$, by discharge or by timeout, and
rejects a second settlement of the same obligation; since that set never
affects any other clause, it is suppressed from the notation. Note the
verification-failure clause: it is admissible only from
$\mathsf{transformed}$, only with an authentic certificate for that
commitment and value, and it leads to a live state, not a terminal one.
Refusal is admissible from every live state, $\mathsf{verificationFailed}$
included, and only by explicit act.

\begin{hypothesis}[Stability]\label{hyp:stable}
The predicates $\valid$ and $\mathsf{checkCert}$ and the maps
$\mathsf{observe}_k$ are deterministic and fixed for the duration of an
execution, relative to the verifier versions recorded in witnesses and
certificates.
\end{hypothesis}

Stability is what lets a verification carried out at one moment count as
evidence at a later moment, and what lets replay re-check a record long
after it was appended. Versions matter because a verifier upgraded
mid-execution changes what $\valid$ means; recording $\kappa$ keeps the
question ``valid under which verifier?'' answerable.

\begin{hypothesis}[Verifier soundness]\label{hyp:verifier}
If $g(c,v)\Downarrow\mathsf{ok}(w)$ then $\valid(w,c,v)$, and if
$g(c,v)\Downarrow\mathsf{fail}(\varphi)$ then
$\mathsf{checkCert}(\varphi,c,v,\mathsf{fail})$. Moreover
$\mathsf{checkCert}$ holds only for certificates actually issued by the
issuer named in $\varphi$, at the version named, on the commitment and
claim named.
\end{hypothesis}

The first half is a requirement on verifiers. The second is an
authentication assumption, met in practice by signing certificates with a
key held by the verifier.



\section{A first example}

Consider a request $q$ that introduces a fresh commitment $c$. A value
$b$ is bound to it by principal $a$, a function $f$ transforms $b$ into
$v$, a verifier $a'$ produces witness $w$, and $a''$ collapses the result
under observation rule $k$. The successful history is
\[
\begin{aligned}
H={}&\langle
  \Pop(c,q),\;
  \Bind(a,c,b),\;
  \Transform(c,f,b,v),\\
 &\hspace{1.2em}\Verify(a',c,v,w),\;
  \Collapse(a'',c,k,w,o)
\rangle .
\end{aligned}
\]
The identity $c$ threads through the record. It is not replaced by $b$ or
$v$. The witness is indexed by $c$ and $v$, preventing evidence produced
for another commitment from being replayed here.

If verification instead rejects $v$, the history continues
\[
  \mathsf{VerificationFailed}(a',c,v,\varphi)\,;\;\Refuse(h,c,r),
\]
where $\varphi$ is the verifier's authenticated failure certificate and
$h$ is a principal with authority to refuse. The verifier changed the
evidentiary state of the commitment; the refusal is a separate act by
someone entitled to dispose of it. This is not a truncated successful
history. It is a complete refused history, and it retains the evidence of
failure beside the decision that followed it. Any later
reconsideration must introduce a successor commitment $c'$ and record that
$c'$ compensates for the practical matter associated with $c$.

\chapter{Disposition Worlds}
\label{ch:worlds}

A formal semantics becomes usable when a reader can build small models of
it by hand, ask whether statements are true in them, and find the smallest
model in which a plausible statement fails. This chapter introduces the
finite models used throughout the book, together with three instruments
for working with them: countermodels, history twins, and causal shuffles.
They are elementary enough to use before any metatheory is available, and
later chapters show that they were instantiating the central theorems all
along. A computation is a finite causal world; a ledger is one
serialization of that world; and a collapse rule declares what an observer
retains.

\section{Worlds and configurations}

\begin{definition}[Disposition world]
A \emph{disposition world} is a tuple
\[
W=\langle E,\prec,\#,\lambda,\mathcal P,\mathcal K\rangle
\]
where $E$ is a finite set of events; $\prec$ is a strict partial order on
$E$ (causal precedence); $\#$ is a symmetric, irreflexive relation on $E$
(incompatibility) that is inherited upward, so that $e\,\#\,e'\prec e''$
implies $e\,\#\,e''$; $\lambda$ labels each event with a ledger record;
$\mathcal P$ is an authority policy, a set of triples $(a,\kappa,c)$
stating that principal $a$ may perform acts of kind $\kappa$ on commitment
$c$; and $\mathcal K$ is a finite set of observation rules. The
commitments of $W$ are the subjects of its labels.
\end{definition}

A \emph{configuration} of $W$ is a set $x\subseteq E$ that is downward
closed under $\prec$ and contains no two incompatible events. It is one
possible state of the computation: the events that have happened. A
\emph{linearization} of $x$ is a listing of its events consistent with
$\prec$, read as a ledger through $\lambda$. A \emph{pointed world}
$(W,x)$ fixes a current configuration.

\begin{definition}[Admissible world]\label{def:admissible-world}
$W$ is \emph{admissible} when
\begin{enumerate}[label=(W\arabic*)]
\item any two events whose records are dependent, in the sense of
Definition~\ref{def:dependence}, are either $\prec$-comparable or
incompatible;
\item every linearization of every configuration has a defined replay;
\item for every event labelled with an attributed record, the required
authority of Chapter~\ref{ch:lts} is in $\mathcal P$.
\end{enumerate}
\end{definition}

Condition (W1) says the world orders everything that replay cares about.
Condition (W2) says every possible state of the world is a legal
disposition state. Condition (W3) says nothing in the world was done
without entitlement.

Worlds of this kind are prime event structures, in the sense of Nielsen,
Plotkin, and Winskel, with events labelled by records and an authority
policy attached. The book's first comparative claim is that this is not an
accident, and this chapter is where the resemblance first becomes
concrete.

A world with empty $\#$ has a single maximal configuration and describes
one execution. Every execution produces one:

\begin{proposition}[Executions are worlds]\label{prop:exec-world}
Let $H$ be the history of an execution from the empty configuration. The
world whose events are the records of $H$, ordered by the dependence order
$\prec_H$ of Definition~\ref{def:dependence}, with empty incompatibility
and the execution's policy, is admissible.
\end{proposition}

\begin{proof}
(W1) holds because dependent records are $\prec_H$-ordered by definition.
For (W2), let $x$ be a downset of $\prec_H$. Some linear extension of
$\prec_H$ lists $x$ first; it is trace equivalent to $H$, so by
Proposition~\ref{prop:commute} its replay is defined, and replay of a
prefix of a defined fold is defined. Any other linearization of $x$ is
trace equivalent to that prefix (Theorem~\ref{thm:shuffle}). (W3) is
Theorem~\ref{thm:auth}.
\end{proof}

\subsection*{What an ordinary event structure lacks}

A disposition world has a plain labelled event structure underneath it.
Write it as
\[
\mathcal E_D=\langle E,\le,\#,\lambda,\eta\rangle,
\]
where $\langle E,\le,\#,\lambda\rangle$ is a labelled prime event
structure whose labels $\lambda(e)$ record only the kind of act and its
subject, and $\eta$ carries the disposition annotations: principal,
reason, witness, certificate, verifier version, and observation rule. The
underlying structure determines configurations and linearizations, and
with them everything in the Causal Shuffle theorem below. What it does not
determine is admissibility: conditions (W2) and (W3) consult $\eta$ through
$\valid$, $\mathsf{checkCert}$, $\mathsf{observe}_k$, and the policy.

This makes the first comparative claim precise. It is settled in the
positive if (W2) and (W3) can be expressed as structural constraints on a
labelled event structure over a finite alphabet, in which case the
calculus is an event-structure discipline with rich labels. It is settled
in the negative if some admissibility condition depends irreducibly on
evidence predicates over unbounded data, in which case $\eta$ is the
residue the claim predicts.

\section{A small language}

Formulas are built from atomic statements about a pointed world, the
Boolean connectives, quantification over the finitely many commitments,
principals, and values of $W$, and one modality.
\begin{itemize}[leftmargin=1.5em]
\item $\mathsf{Live}(c)$, $\mathsf{Refused}(c)$, $\mathsf{Collapsed}(c)$,
and the other state predicates hold at $(W,x)$ when replay of a
linearization of $x$ puts $c$ in the named state.
\item $\mathsf{Occurred}(\rho)$ holds when some event of $x$ is labelled
$\rho$.
\item $\mathsf{Authorized}(a,\kappa,c)$ holds when $(a,\kappa,c)\in\mathcal P$.
\item $\mathsf{Valid}(w,c,v)$ is the validity predicate.
\item $\Diamond\varphi$ holds at $(W,x)$ when $\varphi$ holds at $(W,y)$
for some configuration $y\supseteq x$. $\Box$ is its dual.
\end{itemize}

The state predicates are well defined only because every linearization of
$x$ gives the same replay. That is the Causal Shuffle theorem below, so
the language rests on it from its first clause.

\begin{proposition}[Modal terminal exclusivity]\label{prop:modal-terminal}
In every admissible pointed world,
\[
(W,x)\models\mathsf{Refused}(c)\rightarrow\neg\Diamond\mathsf{Collapsed}(c).
\]
\end{proposition}

\begin{proof}
Suppose $c$ is refused at $x$ and $y\supseteq x$. The refusal event of $c$
is in $y$. In any linearization of $y$, the projection onto $c$ contains a
refusal record, and since $\delta$ is undefined on every record whose
incoming state is terminal, that refusal is the last record of the
projection and the fold yields $\mathsf{refused}$. So $c$ is not collapsed
at $y$.
\end{proof}

\section{Countermodels}

The most useful operation on worlds is the search for the smallest one
that falsifies a proposed statement:
\[
\mathsf{countermodel}(\varphi,n)=\text{a minimal admissible }W
\text{ with at most }n\text{ events and }W\not\models\varphi.
\]
Minimality is lexicographic: fewest commitments, then fewest events, then
fewest principals, then least causal depth. Small countermodels teach more
than arbitrary failing cases, because every event in a minimal
countermodel is doing work.

Two examples, found by the laboratory of Section~\ref{sec:lab} under a
policy in which $a_1$ may bind and $a_2$ may verify:
\begin{itemize}[leftmargin=1.5em]
\item The conjecture ``if $c$ is verified, the principal who bound $c$ is
authorized to verify it'' fails in a four-event world,
$\Pop(c)\,;\Bind(a_1,c)\,;\Transform(c)\,;\Verify(a_2,c)$. The conjecture
was underspecified: authority attaches to acts, not to commitments, and
different acts on one commitment may belong to different principals.
\item The conjecture ``a refused commitment was never verified'' fails in
a five-event world in which a reviewer refuses after a successful
verification. Refusal is available from $\mathsf{verified}$ as from every
live state, and a valid witness does not oblige anyone to accept.
\item The conjecture $\mathsf{Live}(c)\rightarrow\Diamond\mathsf{Collapsed}(c)$,
that every live commitment can still collapse, fails in a four-event
world, $\Pop(c)\,;\Bind(a_1,c)\,;\Transform(c)\,;\mathsf{VF}(a_2,c)$. A
commitment whose verification failed is live, but no continuation leads to
collapse; the only way out is refusal. This countermodel appeared when the
laboratory was rerun after the $\mathsf{verificationFailed}$ state was
introduced, and it is correct: it is the formal content of the decision
that a negative verification does not dispose, and that nothing but an
authorized refusal disposes after it.
\end{itemize}
Conversely, no countermodel with at most six events exists for
Proposition~\ref{prop:modal-terminal}, nor for the weaker reachability
statement $\mathsf{Live}(c)\rightarrow\Diamond\mathsf{Terminal}(c)$ under
that policy.

\begin{remark}[Horizons]
A bounded search for $\Diamond\varphi$ must look far enough ahead. A
search that extends worlds only to the global bound will report
$\mathsf{Pop}(c_1)\,;\Bind(a_1,c_1)\,;\Pop(c_2)$ as a countermodel to
$\mathsf{Live}(c)\rightarrow\Diamond\mathsf{Terminal}(c)$ at bound six,
simply because the remaining budget is too small for $c_1$ to finish. The
laboratory therefore extends each world by the longest lifecycle path,
four events, beyond its own length. The general moral is that a bounded
countermodel to a modal statement is evidence only relative to its
horizon.

The same care applies to what the positive result means. The modality
$\Diamond$ is existential, so the absence of a countermodel to
$\mathsf{Live}(c)\rightarrow\Diamond\mathsf{Terminal}(c)$ is a
\emph{reachability} result: from every live state some continuation of at
most four records reaches a terminal state, and when any principal holds
refusal authority, one record suffices. It is not the liveness statement
that every execution eventually disposes of $c$. Unrelated events may
postpone $c$ indefinitely and external obligations may never be answered;
eventual disposition is a conditional temporal theorem, stated with its
hypotheses in Chapter~\ref{ch:progress}.
\end{remark}

\section{History twins and collapse lenses}

A \emph{lens} is an observation rule $k$ applied to a pointed world,
yielding some value $k(W,x)$. Two pointed worlds are equivalent under $k$,
written $(W,x)\approx_k(W',x')$, when their lens values agree.

\begin{definition}[History twins]
Pointed worlds are \emph{twins} for lenses $(k_1,k_2)$ when they are
equivalent under $k_1$ and not under $k_2$.
\end{definition}

The laboratory finds three canonical twins.
\begin{enumerate}[leftmargin=1.5em]
\item \emph{Same output, different refusals.} A single commitment that is
bound, transformed, verified, and collapsed to $o$; and a world in which the
same happens while a second commitment is opened and refused. A lens
reporting the observed values sees $\{o\}$ in both. A lens reporting
refusals distinguishes them.
\item \emph{Same output, different authority.} The same five-event history
twice, once collapsed by $a_1$ and once by $a_2$. The value lens cannot
tell them apart; a lens reporting who performed the decisive act can.
\item \emph{Same state, different refusal.} The worlds
\[
\Pop(c)\,;\Refuse(a_1,c,r)
\quad\text{and}\quad
\Pop(c)\,;\Bind(a_1,c)\,;\Transform(c)\,;\mathsf{VF}(a_2,c)\,;\Refuse(a_1,c,r).
\]
Both
leave $c$ refused. One refusal was a decision before any work was done;
the other was a decision taken after a certified verification failure,
and the ledger keeps the failure beside the decision.
This is the operational content of indexing the refusal type by the
lifecycle stage, $\mathsf{Refusal}\langle c,A,s\rangle$, in
Chapter~\ref{ch:types}.
\end{enumerate}
Two worlds are never simply equivalent. They are equivalent under a
declared lens.

\begin{definition}[Lens refinement]
Lens $k_1$ \emph{refines} $k_2$, written $k_1\preceq k_2$, when
$k_2=h\circ k_1$ for some function $h$: everything $k_2$ reports can be
computed from what $k_1$ reports.
\end{definition}

\begin{proposition}\label{prop:lens}
Refinement is a preorder, and if $k_1\preceq k_2$ then
$(W,x)\approx_{k_1}(W',x')$ implies $(W,x)\approx_{k_2}(W',x')$.
\end{proposition}

\begin{proof}
Reflexivity takes $h$ to be the identity and transitivity composes the
two functions. If $k_1$ values agree, applying $h$ to both gives equal
$k_2$ values.
\end{proof}

Lenses identified up to mutual refinement correspond to the equivalence
relations they induce on pointed worlds, and equivalence relations on a
set form a complete lattice under inclusion. The collapse rules of the
calculus therefore form a structured spectrum rather than an unrelated
collection. Whether that spectrum coincides with van Glabbeek's for the
classical equivalences, when those are recast as lenses, is the question
of Chapter~\ref{ch:spectrum}.

\section{Causal shuffles}

The third instrument enumerates every ledger that could have recorded one
world: all linearizations of a configuration.

\begin{theorem}[Causal shuffle]\label{thm:shuffle}
In a world satisfying (W1), all linearizations of a configuration have
the same replay: either all are undefined, or all are defined and equal.
In particular this holds in every admissible world.
\end{theorem}

\begin{proof}
Any two linear extensions of a finite partial order are connected by a
sequence of exchanges of adjacent elements that are incomparable in the
order: move the element that comes first in the target order to the front
by successive adjacent exchanges, each of which passes an element that the
target places later and which therefore cannot precede it in the order,
and recurse on the rest. Two incomparable events in a configuration are not
incompatible, so by (W1) their records are independent. By
Proposition~\ref{prop:commute}, exchanging adjacent independent records
does not change replay, and does not change whether it is defined.
The proof uses only (W1), which is why it can be invoked in establishing
(W2) for Proposition~\ref{prop:exec-world}.
\end{proof}

The theorem identifies exactly when the sequence numbers a sequencer
assigns are incidental. Every order among independent records is one
serialization of the same world. Order is semantically necessary only
between dependent records, and (W1) requires the world to state it. For
the finite class $\mathcal H_{\le 7}$ of Section~\ref{sec:lab}, the
laboratory confirms the theorem instance by instance: each of the
$372{,}961$ linearizations of the $15{,}167$ ledgers in the class replays
to the same state as the ledger it came from.

\section{The witness mixer}

Let $c_1$ and $c_2$ both be $\mathsf{transformed}(b,v)$, with witnesses
$w_1$ produced for $(c_1,v)$ and $w_2$ produced for $(c_2,v)$. Try verifying
$c_1$ with $w_2$. A check that asked only whether the witness vouches for
the value $v$ would accept, since the values coincide. The step function
refuses, because it asks for $\valid(w_2,c_1,v)$.

That refusal depends on the verifier, not only on the calculus. Indexing
forces the question to be asked about the right commitment, but the
predicate must actually distinguish commitments for the answer to differ.

\begin{hypothesis}[Index sensitivity]\label{hyp:index}
$\valid(w,c,v)\wedge\valid(w,c',v)$ implies $c=c'$.
\end{hypothesis}

Under this hypothesis evidence is non-transferable: a witness valid for
one commitment is valid for no other, and the witness ancestry of
Corollary~\ref{cor:witness} cannot be satisfied by borrowed evidence. The
hypothesis is a requirement on verifiers, typically met by signing or
hashing the commitment identity into the witness, and it is stated here
because nothing in the operational rules can supply it.

\section{Perturbations}

The laboratory's third mode takes an admissible world and changes exactly
one thing. On the five-event successful history, every perturbation but
one is caught, and each is caught by a specific clause. A sixth row comes
from a third base, a commitment refused after a certified verification
failure, whose certificate is replaced by a forged one:
\begin{center}
\small
\begin{tabular}{ll}
\toprule
Perturbation & Detected by\\
\midrule
unauthorized binder & authority (W3)\\
foreign witness in verification & $\delta$: $\valid(w,c,v)$ fails\\
collapse citing a different witness & $\delta$: witness mismatch\\
exchange of two dependent records & $\delta$: wrong incoming state\\
reuse of a commitment identity & $\delta$: identity not fresh\\
forged failure certificate & $\delta$: $\mathsf{checkCert}$ fails\\
deletion of an interior record & $\delta$: wrong incoming state, \emph{usually}\\
deletion of the final record & \emph{not detected}\\
\bottomrule
\end{tabular}
\end{center}
The last row is not a flaw in the checker. A ledger with its final record
removed is a prefix of an admissible ledger, and every prefix of an
admissible ledger is admissible, since the world it describes is simply an
earlier configuration. Replay certifies that what is present could have
happened. It cannot certify that nothing is missing from the end.

This is an impossibility, not an oversight, and Chapter~\ref{ch:replay}
states it as the Prefix Indistinguishability theorem together with the
remedy: a signed checkpoint retained outside the ledger. Completeness
itself requires a witness outside the history whose completeness is in
question.

The qualification in the interior row is a finding of the second base.
Deleting the transformation from
\[
\Pop(c)\,;\Bind(a_1,c)\,;\Transform(c)\,;\Refuse(a_1,c,r)
\]
leaves $\Pop(c)\,;\Bind(a_1,c)\,;\Refuse(a_1,c,r)$, which is admissible,
because explicit refusal is legal from the bound state as well. The same
generosity that makes refusal available at every live stage makes some
interior deletions invisible to replay. Any deletion whose remainder is
still a legal path escapes detection, and the same anchor is the remedy:
a hash chain detects every modification, not only truncation.

\section{Disposition control}
\label{sec:control}

The instruments so far evaluate statements in given worlds. A different
question asks which decisions determine a disposition: what must be fixed
so that a commitment is guaranteed to end one way.

\begin{definition}[Interventions, forcing, influence]
An \emph{intervention} $K$ is a set of constraints on decision points,
each removing some records from the enabled set wherever they would occur:
for instance ``verifications of $c$ succeed'' or ``principal $a$ never
refuses $c$''. A \emph{completion} under $K$ is a maximal admissible
ledger all of whose steps respect $K$. Then
\begin{itemize}[leftmargin=1.5em]
\item $\mathsf{force}(K,d)$ holds when \emph{every} completion under $K$
leaves $c$ in disposition $d$;
\item $\mathsf{influence}_\mu(K,d)=\Pr_\mu\bigl[\text{a completion under
$K$ leaves $c$ in }d\bigr]$, for an explicitly declared distribution $\mu$
over schedules, initial states, or worlds.
\end{itemize}
\end{definition}

The two are different kinds of claim. Forcing is a universal statement
over completions and is proved or refuted by exhaustion. Influence is a
probability and depends on $\mu$; a high influence is not a guarantee.
Without a declared $\mu$, a figure of this kind is not a probability at
all: it can only be a coverage figure or a reachability count, and should
be reported as one. The
distinction is taken from recent work on control of Boolean network
models of biological systems, where controlling a feedback vertex set
together with the source nodes carries a structural guarantee of reaching
a target attractor, while smaller control sets found through network
modularity are evaluated by \emph{percent coverage}: the proportion of the
target outcome that fixing the control set is shown to reproduce. In that
study's biological models, one variant reached about 87\% average
coverage with control sets slightly larger than the minimum feedback
vertex set, and another, choosing one high-ranked feedback vertex per
module, reached about 77\% with sets about 63\% of that size
\cite{nasrollahi}. Coverage measures partial achievement of a target, and
influence as defined here is a frequency over schedules, so the two
quantities differ. The lesson carries over all the same: a measure of
partial success, however high, is not a guarantee.

For a single commitment under the laboratory's policy, with $\mu$ the
uniform scheduler that picks among enabled records with equal
probability:
\begin{itemize}[leftmargin=1.5em]
\item With no intervention, the influence toward collapse is $1/72$ and
toward refusal $71/72$. The asymmetry reflects only that refusal options
outnumber advancing ones at every stage under a uniform scheduler, not
anything about real systems.
\item The unique minimal forcing set for \emph{collapse} has three
elements: verifications succeed, and neither principal refuses.
\item The unique minimal forcing set for \emph{refusal} has one element:
verifications fail.
\item The best two-element set that does not force collapse, ``neither
principal refuses'', reaches influence $1/2$: every remaining completion
turns on the verifier.
\end{itemize}
Forcing refusal is cheap and forcing collapse is expensive, because
refusal is available everywhere and collapse only at the end of one path.
That asymmetry is structural. It is the formal shadow of the book's
insistence that refusal be available from every live state.

\begin{proposition}[No cycles, one re-entry]\label{prop:reentry}
The lifecycle is acyclic, and each commitment's projection contains at
most one opening, four lifecycle steps, and one discharge per obligation.
The only rule enabled by a terminal state is compensation, and it creates
a new commitment rather than re-entering the old one. Hence in this
calculus an infinite execution is never a cycle in any commitment's
lifecycle; it is an unbounded supply of new commitments, arriving through
replication or through chains of compensation.
\end{proposition}

\begin{proof}
Acyclicity and the bound on lifecycle steps are properties of $\delta$.
Discharges are bounded by the obligation set. Compensation is the only
rule with a terminality premise, and its subject is fresh. An infinite
execution therefore introduces infinitely many commitments, and
introductions come only from pop, whose instances beyond the finitely many
in the program text come from replication, and from compensation.
\end{proof}

In the Boolean-network setting, a feedback vertex set matters because
cycles are what sustain undesired attractors. Here there are no cycles to
break in the basic calculus, and a feedback-set analysis would be
ornamental. What plays the corresponding role is the set of re-entry
points: replication, and the authority to compensate. Removing
compensation authority leaves replication as the only source of unbounded
work. If the calculus is extended with reopening, retries, delegation
cycles, or mutually dependent verification, genuine cycles appear, and a
theorem of the feedback-set kind becomes available: an intervention that
meets every live dependency cycle removes cyclic obstructions to
disposition. It would still not guarantee eventual disposition, which
needs fairness and responsive dependencies as well
(Chapter~\ref{ch:progress}).

\section{The laboratory}
\label{sec:lab}

The examples in this chapter were produced by a small executable
reference, \texttt{disposition\allowbreak\_lab.py}, which implements the records, the
step function, and replay as the book defines them, and enumerates every
admissible ledger up to a bound by running a maximally nondeterministic
process under a fixed policy. Its results are bounded propositions about a
named finite class, not theorems, and they are versioned with the grammar
they were computed for.

Let $\mathcal H_{\le 7}$ be the class of admissible ledgers of at most
seven records over two commitment identities, two principals, one value,
one transformation, one observation rule, one external obligation per
commitment with one declared deadline, and the policy stated in the
source. For version~4 of the laboratory, whose grammar is the book's
current record grammar, including the $\mathsf{verificationFailed}$ state,
authenticated certificates, and discharge and timeout records,
$|\mathcal H_{\le 7}|=15{,}167$. The laboratory reports each check as a
named property with its class, search depth, normalization, expected
result, and, when it fails, a minimized counterexample, so that it serves
as an executable index of the proof spine. Every property below holds on
every $H\in\mathcal H_{\le 7}$:
\begin{center}
\small
\begin{tabular}{ll}
\toprule
Property & Statement checked\\
\midrule
P-ReplaySound & $\delta^{*}$ of each projection equals replay's map (the path lemma)\\
P-TerminalExclusive & at most one terminal record per commitment, and it is last\\
P-CertificateAuthentic & every verification failure checks against its state and a verifier\\
P-FailedVerificationLive & a verification failure leaves its commitment live\\
P-NoVerifierDisposition & no verification record yields a terminal state\\
P-TraceInvariant & all $372{,}961$ linearizations replay to the state of their ledger\\
P-WorldAdmissible & every configuration has a defined replay\\
\bottomrule
\end{tabular}
\end{center}
In addition, for each of the $15{,}332$ pointed worlds in which some
commitment is refused, no admissible extension within the horizon
collapses it. The companion module \texttt{disposition\allowbreak\_parser.py}
checks the parser properties of Chapter~\ref{ch:parsing}. Two properties
named in the proof spine, P-RightsAgree and P-DelegationConservative, need
the process and typed layers, which the laboratory does not yet model;
they are proved in Chapters~\ref{ch:session} and~\ref{ch:authority} and
not checked here.

Earlier versions checked earlier grammars, and their figures apply to
those grammars only: version~1 ($1{,}230$ ledgers), version~2 with
certificates and discharge ($6{,}305$), and version~3 with timeout and
automatic refusal on failed verification ($14{,}608$). The countermodels,
twins, and control figures are unchanged across versions, with one
deliberate exception: the countermodel to
$\mathsf{Live}(c)\rightarrow\Diamond\mathsf{Collapsed}(c)$, which exists
only once a failed verification leaves its commitment live. The
laboratory's value is not in the counts. It is in finding the smallest
world where an idea goes wrong, which is how the horizon remark, the
interior-deletion finding, the forged-certificate row, and that
countermodel were found.

\section*{Exercises}
\begin{enumerate}[leftmargin=1.5em]
\item Construct the smallest admissible world in which a commitment is
compensated. How many events does it need, and why can none be removed?
\item Find twins for the lenses ``set of commitments in a terminal state''
and ``set of observations''.
\item Give a world violating (W1) in which two linearizations of the same
configuration replay differently.
\item Define a lens that distinguishes twin~3 above but not twin~1, and
place it in the refinement order relative to the value lens and the
refusal lens.
\item Show that without Hypothesis~\ref{hyp:index} there is an admissible
world in which $c_1$ is verified with a witness that was produced for
$c_2$.
\end{enumerate}

\chapter{The Kernel and the Disposition Layer}
\label{ch:layers}

The disposition calculus uses the names of the Spherepop kernel's
operators. This chapter decides whether that is a formalization of the
kernel or a related calculus built above it, by giving an explicit
translation
\[
\llbracket-\rrbracket:\text{disposition records}\longrightarrow
\text{kernel histories}
\]
and proving two things about it: that every disposition step translates
to an admissible kernel step (forward simulation), and that every kernel
step of translated shape whose guard holds comes from a disposition step
(backward adequacy). The records and the step function it uses were
defined in Section~\ref{sec:records}; the operational rules the guards
refer to are stated in Chapter~\ref{ch:lts}, and the corollary relating
the two uses consistency, established in Chapter~\ref{ch:replay}. Until
both directions are proved, the relation between the layers is a
conjecture, and it should be read that way.

\section{The kernel, restated}

A kernel history is an append-only sequence of events. Four causal
operators act on histories.
\begin{itemize}[leftmargin=1.5em]
\item $\Pop(x)$ commits to the option $x$ among those currently available,
removing it from what remains possible.
\item $\Refuse_r(X)$ documents that the options in $X$ are inadmissible,
with reason $r$, without removing them from the record of what was
considered.
\item $\Bind(H_1,H_2)$ couples two histories so that each constrains the
other's continuation, without identifying them.
\item $\Collapse_k(H)$ observes a history under a rule $k$, projecting it
onto the quotient $k$ induces.
\end{itemize}
Three structural rules build the terms these operators act on: sequence
$H_1;H_2$, tensor $H_1\otimes H_2$, which forms a joint context from two
histories so that an operator can condition one on the other, and branch,
which divides one history into two continuations sharing a prefix. A
separate annotation channel carries metadata that participates in no
causal reduction.

One feature of the kernel matters more than any other here. Applying
$\Collapse_k$ computes an observation. It does not by itself make that
observation a commitment: a history can be observed under many rules, and
none of those observations constrains what happens next unless something
commits to it. The corpus's operating-system specification states this as
``collapse is rendering'' and keeps rendered views out of the
authoritative log. The translation below adopts that reading. Under the
alternative reading, in which every kernel collapse is itself logged, one
step of the translation is absorbed and nothing else changes; the paragraph
on collapse in Section~\ref{sec:translation} says which.

\section{Kernel histories and admissibility}

To state a translation precisely, the kernel needs a semantics of its own,
and it should be the most permissive one consistent with the four
operators. Fix a universe $\mathcal X$ of options. A \emph{kernel history}
$K$ is a finite sequence of events
\[
\Pop(x),\qquad \Refuse_r(X),\qquad \Bind(x,y),
\]
for options $x,y\in\mathcal X$, sets $X\subseteq\mathcal X$, and reasons
$r$. Write $\mathsf{Popped}(K)$ for the options popped in $K$ and
$\mathsf{Refused}(K)$ for the union of the refused sets. An event is
\emph{kernel admissible} after $K$ when:
\begin{itemize}[leftmargin=1.5em]
\item $\Pop(x)$: $x\notin\mathsf{Popped}(K)\cup\mathsf{Refused}(K)$, so an
option can be committed to at most once, and never after it has been
documented as inadmissible;
\item $\Refuse_r(X)$: always, since documenting inadmissibility removes
nothing;
\item $\Bind(x,y)$: $x,y\in\mathsf{Popped}(K)$, so only things that exist
can be coupled.
\end{itemize}
A kernel history is admissible when each event is admissible after its
prefix. The kernel's derived choice, $\mathsf{Choice}(x,X)=\Pop(x)\parallel
\Refuse(X\setminus\{x\})$, commits to one option and documents the rest as
inadmissible. Kernel $\Collapse_k$ is a function on histories, a
rendering, and not an event.

This semantics enforces very little, and that is deliberate. It says
nothing about the order in which options may be popped, nothing about
what data they carry, and nothing about who may pop them. Whatever the
disposition calculus requires beyond freshness and the permanence of
refusal must therefore be supplied by the translation, and the question
of this chapter is exactly how much that is.

\section{Continuation spaces and guards}

The options of the translation are the records themselves. A record
popped in the kernel is a commitment to that record having occurred.

\begin{definition}[Continuation space]
For a ledger $H$ and commitment $c$, the \emph{admissible continuation
space} is
\[
\Omega_c(H)=\{\tau\mid\text{every record of }\tau\text{ has subject }c
\text{ and }\delta^{*}(H|_c\cdot\tau)\text{ is defined}\},
\]
the set of futures that $c$'s own lifecycle still admits. Its one-step
part, the records that could come next, is
$\Omega^1_c(H)=\{\rho\mid\langle\rho\rangle\in\Omega_c(H)\}$. Write
$\mathsf{Opt}_c(H)$ for the set of all records with subject $c$ not yet
popped.
\end{definition}

For a live commitment $\Omega_c(H)$ contains the continuations its state,
the data it carries, and the evidence predicates allow. For a terminal
commitment it contains only the empty continuation.

\begin{definition}[Guard]
The \emph{guard} of a record $\rho$ at ledger $H$, written $G_\rho(H)$, is
the conjunction of the premises of the operational rule that appends
$\rho$, read against the state that replay of $H$ reconstructs: the
lifecycle and data conditions ($\rho\in\Omega^1_{\subj(\rho)}(H)$), the
authority premise, the evaluation or verification premise where there is
one ($f\,b\Downarrow v$, $g(c,v)\Downarrow\mathsf{ok}(w)$ or
$\mathsf{fail}(\varphi)$), and for compensation the terminality of the
predecessor.
\end{definition}

The guard is a statement in the metalanguage. It is not a kernel operator,
and the translation does not emit an event for it. A step whose guard
holds is translated; a step whose guard fails is simply not a
disposition step. If an implementation wishes to record a failed attempt,
it may do so as a kernel refusal of the proposed option or on the
annotation channel, but nothing in what follows depends on that choice.

\section{The translation}
\label{sec:translation}

\begin{definition}[Record translation]\label{def:translation}
For a record $\rho$ with subject $c$ whose guard $G_\rho(H)$ holds, let
$\ell_c$ be the last record of $H|_c$, and define $\llbracket\rho\rrbracket_H$ by
\[
\begin{aligned}
\llbracket\Pop(c,q)\rrbracket_H&=\Pop(\rho),\\
\llbracket\Compensates(a,c,q,d)\rrbracket_H&=\Pop(\rho);\Bind(\rho,\ell_d)
  \quad\text{over }H_c\otimes H_d,\\
\llbracket\rho\rrbracket_H&=\Pop(\rho);\Bind(\rho,\ell_c)
  \quad\text{for every other non-terminal }\rho,\\
\llbracket\Refuse(a,c,\varrho)\rrbracket_H&=
  \Pop(\rho)\parallel\Refuse_{\langle a,\varrho\rangle}(\mathsf{Opt}_c(H)\setminus\{\rho\});\;
  \Bind(\rho,\ell_c),\\
\llbracket\Collapse(a,c,k,w,o)\rrbracket_H&=
  \Pop(\rho)\parallel\Refuse(\mathsf{Opt}_c(H)\setminus\{\rho\});\;
  \Bind(\rho,\ell_c),
  \quad o=\Collapse_k(v).
\end{aligned}
\]
The other non-terminal records are the advancing records ($\Bind$,
$\Transform$, $\Verify$) and the discharge and timeout records.
The translation of a ledger is $\llbracket\varepsilon\rrbracket=\varepsilon$
and $\llbracket H\cdot\rho\rrbracket=\llbracket H\rrbracket;
\llbracket\rho\rrbracket_H$, defined when every guard holds at its
prefix. Conversely, the \emph{decoding} of a kernel history lists its
popped records in order.
\end{definition}

The kernel history component $H_c$ of a commitment is the chain of its
popped records linked by $\Bind$: each record of $c$ is coupled to the one
before it, without being identified with it. The clauses fall into four
groups.

\paragraph{Openings.} Introducing a commitment pops its opening record.
Freshness needs nothing extra: a popped option can never be popped again.
Compensation also binds the new opening to the predecessor's last record,
which needs the tensor rule to form the joint context of two histories,
the same requirement that waiting on a child process and mounting a
namespace place on the operating-system specification. $\Bind$ is the
right operator because it couples without identifying, the property
developed at length for the kernel's binding in \cite{flyxion-binding}: the successor
depends on the predecessor and is a different commitment.

\paragraph{Advancing, discharge, and timeout.} Binding, transformation,
verification, discharge, and timeout all have the same shape, $\Bind\circ\Pop$ onto
the commitment's history. They differ in payload and in guard. Nothing
about transformation or verification requires a kernel operator the other
two do not; their distinctive content is in the conditions under which
they are admissible.

\paragraph{Terminal steps.} Refusal and collapse are both instances of the
kernel's derived choice: one option is popped and every other option of
$c$ is documented as inadmissible. What distinguishes them is what is
popped. A refusal pops an attributed, reasoned decision, and the refused
set carries that principal and reason. A collapse pops a committed
observation $o=\Collapse_k(v)$: kernel collapse renders an observation,
and the pop makes it a commitment. Disposition refusal is therefore kernel
refusal of the \emph{whole remaining option space} of $c$, which is what
distinguishes it in scope from kernel refusal of a single option; and
disposition collapse is $\Pop\circ\Collapse_k$, a committed observation,
rather than kernel collapse alone.

\paragraph{Where the discipline lives.} Because the terminal translations
refuse every remaining option of $c$, the kernel's own admissibility
enforces finality: after either terminal step, no further record of $c$
can be popped. Freshness is likewise enforced by the kernel. Everything
else, the order of the lifecycle, the matching of data between records,
the validity of evidence, and authority, is enforced only by the guard.
That is the precise content of calling the disposition calculus a
discipline over the kernel.

\section{Forward simulation and backward adequacy}

\begin{lemma}[Records occur once]\label{lem:once}
If $\replay(H)$ is defined, no record occurs twice in $H$.
\end{lemma}

\begin{proof}
A repeated record has the same subject as its first occurrence, so both
lie in one projection. Every clause of $\delta$ other than discharge and
timeout moves strictly along the acyclic lifecycle or out of $\bot$, so
the second occurrence would meet a state from which its clause is
undefined. A repeated discharge or timeout of the same obligation is
rejected by the set of settled obligations carried beside the state.
\end{proof}

\begin{theorem}[Forward simulation]\label{thm:fwd}
If $\replay(H)$ is defined and $G_\rho(H)$ holds, then
$\llbracket H\rrbracket;\llbracket\rho\rrbracket_H$ is a kernel-admissible
history whose decoding is $H\cdot\rho$.
\end{theorem}

\begin{proof}
By induction on $H$ the prefix $\llbracket H\rrbracket$ is admissible and
decodes to $H$. For the new events: $\rho$ has not been popped, by
Lemma~\ref{lem:once} applied to $H\cdot\rho$, whose replay is defined
because the guard includes $\rho\in\Omega^1_c(H)$. It has not been
refused, because the only refusals of options with subject $c$ are those
emitted by a terminal step of $c$, and a terminal $c$ has an empty
one-step space, contradicting the guard. So $\Pop(\rho)$ is admissible.
Each $\Bind$ couples $\rho$ to a record already popped. Each $\Refuse$ is
admissible unconditionally. Decoding lists the popped records in order,
which is $H$ followed by $\rho$.
\end{proof}

\begin{corollary}[Operational forward simulation]
If $C$ is consistent and $C\xrightarrow{\alpha}C'$ appends $\rho$, then
$\llbracket H\rrbracket\xRightarrow{\llbracket\rho\rrbracket_H}\llbracket H'\rrbracket$
in the kernel.
\end{corollary}

\begin{proof}
The rule that fired has the guard's premises: its lifecycle and data
premises give $\rho\in\Omega^1_c(H)$ because replay of $H$ reconstructs
$A$ (consistency), and its authority premise is Theorem~\ref{thm:auth}.
Apply Theorem~\ref{thm:fwd}.
\end{proof}

The converse direction is where the phrase ``discipline over the kernel''
earns its meaning. Without it, the translation might embed disposition
steps into a more permissive system without the permissiveness showing.

Call a kernel extension $\beta$ of $\llbracket H\rrbracket$
\emph{typed-guarded} if it has the shape $\llbracket\rho\rrbracket_H$ for
some record $\rho$ and $G_\rho(H)$ holds. Let $U$ be the \emph{universal
process} over a finite set of names and values: the replication of the
sum of every prefix, each followed by $\mathbf 0$. It can attempt every
disposition step, and it performs exactly those whose premises hold.

\begin{theorem}[Backward adequacy]\label{thm:bwd}
Let $C=\langle U,S,A,O,H\rangle$ be consistent, and let
$\llbracket H\rrbracket;\beta$ be kernel admissible with $\beta$
typed-guarded, of shape $\llbracket\rho\rrbracket_H$. Then there is a
transition $C\xrightarrow{\alpha}C'$ appending $\rho$, and
$\llbracket H\rrbracket;\beta=\llbracket H'\rrbracket$.
\end{theorem}

\begin{proof}
The guard is the conjunction of the premises of the rule for $\rho$, read
against replay of $H$, which by consistency is the runtime state. The
universal process offers the corresponding prefix, which by the guard is
enabled; firing it appends $\rho$. The equation is the definition of the
translation.
\end{proof}

Together the two theorems say that, under the opaque reading and at the
level of ledgers, translated disposition histories are exactly the kernel
histories whose steps pass their guards. The disposition calculus adds no
causal primitive and removes no kernel history that its guards admit. It
is in this sense a \emph{conservative guarded presentation} of the kernel.

Three limits should be kept in view. First, the theorems are relative to
the universal process; a particular program offers fewer steps, as every
program does, and that is not a failure of adequacy. Second, they hold
under the opaque reading, in which evaluation and verification results are
inputs checked by the guard. They do not show that the kernel performs
those computations. Third, the guard lives in the metalanguage: the kernel
does not compute it.

\section{Opaque and transparent readings}

Under the \emph{opaque} reading, the results of evaluation and
verification enter as recorded data whose correctness is part of the
guard. That is the reading of the theorems above, and it supports every
operational and ledger result of the book without any claim about what
the kernel can compute.

Under the \emph{transparent} reading, the value $v$ of a transformation
is itself produced inside the kernel by its encoding of application,
$\Collapse_{\mathrm{subst}}(\Bind(f,b))$, and similarly for verifiers. Only
this reading would support the claim that Spherepop itself performs the
computation. It requires a theorem the book does not yet have: that the
kernel's application encoding agrees with evaluation in the lambda layer,
$\Collapse_{\mathrm{subst}}(\Bind(f,b))=v$ exactly when $f\,b\Downarrow v$.
It is stated as a conjecture and listed in Appendix~\ref{app:open}.

\section{What the translation preserves}

\begin{proposition}[Compositionality]
$\llbracket H\cdot H'\rrbracket=\llbracket H\rrbracket;\llbracket H'\rrbracket_H$,
where the second translation is taken record by record at the successive
prefixes.
\end{proposition}

\begin{proof}
Induction on $H'$ using the defining clause for histories.
\end{proof}

\begin{proposition}[Locality]\label{prop:locality}
The translation of a record pops and binds only options with its own
subject, except that a compensation record also binds to the last record
of its predecessor.
\end{proposition}

\begin{proof}
Inspection of Definition~\ref{def:translation}.
\end{proof}

From these the translation preserves commitment identity (distinct
commitments occupy disjoint chains of options, and no clause identifies
two of them); causal order (records of one commitment are chained by
$\Bind$ in ledger order, while independent records touch disjoint options
and may be interleaved, so the dependence order of the ledger is carried to
the kernel and nothing is added); refusal records (each disposition
refusal is a kernel refusal in an append-only history); and declared
observations (each $(c,k,o)$ in $O$ corresponds to exactly one popped
collapse record carrying $k$). Finally, the commitment map itself has a
kernel reading: $\delta$ is an observation rule on histories, and
$A=\Collapse_\delta(\llbracket H\rrbracket)$ after decoding. The
disposition state is rendered from the history, not stored beside it.

\section{Decisions fixed by this chapter}

\begin{enumerate}[leftmargin=1.5em]
\item \emph{No new causal primitive.} Under the opaque reading and at the
level of ledgers, the disposition calculus is a conservative guarded
presentation of the kernel (Theorems~\ref{thm:fwd} and~\ref{thm:bwd}).
Transformation, verification, and discharge are guarded $\Bind\circ\Pop$.
Whether the kernel also performs the computations, the transparent
reading, is conjectural.
\item \emph{Refusal is available from every live state} and translates to
a choice that documents the whole remaining option space of $c$ as
inadmissible.
\item \emph{Collapse is always indexed by a rule $k$}, recorded in the
ledger, and means a committed observation, with observable label
$\langle k,o\rangle$.
\item \emph{Verification records both outcomes and disposes of neither.}
Success appends a $\Verify$ record and leads to $\mathsf{verified}$;
failure appends a $\mathsf{VerificationFailed}$ record carrying an
authenticated certificate and leads to the live state
$\mathsf{verificationFailed}$, from which only an authorized refusal
exits.
\item \emph{Plain pop requires no authority.} Introducing an identity
commits no principal to anything; the first attributed act on a
commitment is its binding, discharge, or refusal.
\item \emph{Replay reconstructs the commitment map and observation set},
not the store. The store is reconstructible exactly when every
state-affecting input is recorded.
\item \emph{Discharge and timeout are kernel-derived.} Each is a guarded
$\Bind\circ\Pop$ of authenticated evidence onto a waiting commitment's
history, not a fifth primitive and not an annotation. Neither ends a
commitment.
\end{enumerate}

\chapter{Encodings and Their Criteria}
\label{ch:encodings}
\plan{Gorla's five criteria (compositionality, name invariance,
operational correspondence, divergence reflection, success sensitivity)
and full abstraction, with examples meeting some and failing others. Three
possible outcomes of any comparison: equivalence, strict inclusion,
incomparability. Why Turing completeness says almost nothing about
expressiveness in the sense needed here.}

%===================================================================
\part{Terms and Names}
%===================================================================

\chapter{Syntax and Binding}
\label{ch:syntax}

\section{Two syntactic categories}

The calculus separates expressions from processes. Expressions compute;
processes interact and alter disposition. The expression grammar is
\[
\begin{aligned}
e ::= {}& x \mid \lambda x{:}A.e \mid e\,e
\mid \langle e,e\rangle \mid \pi_1 e \mid \pi_2 e\\
&\mid \mathsf{inl}\;e \mid \mathsf{inr}\;e
\mid \mathsf{case}\;e\;\mathsf{of}\cdots .
\end{aligned}
\]
The process grammar is
\[
\begin{aligned}
P,Q ::= {}& \mathbf 0
\mid \mathsf{pop}\;q\;\mathsf{as}\;c.P
\mid \mathsf{pop}_a\;q\;\mathsf{as}\;c\;\mathsf{after}\;d.P\\
&\mid \mathsf{refuse}_a\;c\;\mathsf{because}\;r.P
\mid \mathsf{bind}_a\;c\;\mathsf{to}\;b.P
\mid \mathsf{transform}\;c\;\mathsf{with}\;f.P\\
&\mid \mathsf{verify}_a\;c\;\mathsf{using}\;g\;\mathsf{then}\;P\;
\mathsf{else}\;Q
\mid \mathsf{collapse}_{a,k}\;c.P\\
&\mid \mathsf{await}_a\;c\;\omega.P
\mid \mathsf{await}_a\;c\;\omega.P\;\mathsf{within}_{b}\;t\;\mathsf{else}\;Q\\
&\mid P\mid Q \mid (\nu c)P \mid P+Q \mid !P.
\end{aligned}
\]
Subscripts name the acting principal $a$ and, for collapse, the
observation rule $k$. Pop and transformation are unattributed: the first
only introduces an identity, and the second is a lambda-layer computation.
The compensating pop $\mathsf{pop}_a\;q\;\mathsf{as}\;c\;\mathsf{after}\;d$
introduces a successor $c$ to a terminal commitment $d$ and is attributed,
since reconsidering a disposition is itself an act someone must be
entitled to perform.

Verification carries two continuations because it may decide either way.
A single-continuation form would leave a negative decision without a
successor state, and the commitment would sit in $\mathsf{transformed}$
with nothing recorded, which is exactly the unrecorded dead end the
lifecycle was designed to exclude.

Collapse names no witness. It uses the witness recorded in the
commitment's verified state, so a program cannot cite the wrong one, and a
well-typed program cannot stall on a mismatch between the witness it holds
and the one the state holds. The prefix $\mathsf{await}_a\;c\;\omega.P$
suspends the process until principal $a$ discharges an external
obligation $\omega$ on which commitment $c$ depends, such as a committee
decision or a response from an outside service. It is the calculus's
positive form of waiting: a process that is waiting says so, names the
commitment, and names who it is waiting for, rather than simply failing to
move. Waiting itself writes nothing to the ledger; only a successful
discharge does. The deadline form
$\mathsf{await}_a\;c\;\omega.P\;\mathsf{within}_b\;t\;\mathsf{else}\;Q$
adds an explicit branch: either $a$ discharges $\omega$ and the process
continues as $P$, or a clock authority $b$ certifies that the interval $t$
passed without discharge and the process continues as $Q$. A timeout does
not convert silence into refusal. It records that the obligation went
unanswered, and it is up to an authorized rule in $Q$ to decide whether
the commitment then waits further, is compensated, or is refused.

Replication creates new commitments on every unfolding. Because
$\mathsf{pop}$ binds its identity, the law $!P\equiv P\mid\,!P$ together
with $\alpha$-conversion yields a fresh copy of each generative binder in
each copy of $P$, and the pop rule's freshness premise forces distinct
identities when those copies fire. No two instances ever share a
commitment.

A negative verification does not refuse the commitment. The failure
branch records the verifier's certificate and continues as $Q$ with the
commitment live in $\mathsf{verificationFailed}$; if the commitment is to
end, $Q$ must contain a refusal by a principal entitled to refuse. A
verifier therefore changes the evidentiary state of a commitment without
ever disposing of it.

The binder in $\mathsf{pop}\;q\;\mathsf{as}\;c.P$ scopes over $P$, as does
the binder of the compensating pop and restriction $(\nu c)P$. Bound
variables in lambda abstractions scope over their bodies. Commitment
identities and expression variables form disjoint sorts, and they remain
sorted even if a concrete implementation represents them with the same
strings.

\begin{definition}[Alpha-equivalence]
Alpha-equivalence, written $M\equiv_\alpha N$, is the least congruence
identifying terms which differ only by a consistent capture-avoiding
renaming of bound variables or bound commitment identities.
\end{definition}

\section{Free names and substitution}

The free commitment names of representative processes are defined by
\[
\begin{aligned}
\fn(\mathbf 0)&=\varnothing,\\
\fn(P\mid Q)&=\fn(P)\cup\fn(Q),\\
\fn((\nu c)P)&=\fn(P)\setminus\{c\},\\
\fn(\mathsf{pop}\;q\;\mathsf{as}\;c.P)&=\fn(q)\cup
  (\fn(P)\setminus\{c\}),\\
\fn(\mathsf{pop}_a\;q\;\mathsf{as}\;c\;\mathsf{after}\;d.P)&=
  \{d\}\cup\fn(q)\cup(\fn(P)\setminus\{c\}),\\
\fn(\mathsf{bind}_a\;c\;\mathsf{to}\;b.P)&=
  \{c\}\cup\fn(b)\cup\fn(P),\\
\fn(\mathsf{verify}_a\;c\;\mathsf{using}\;g\;\mathsf{then}\;P\;
\mathsf{else}\;Q)&=\{c\}\cup\fn(g)\cup\fn(P)\cup\fn(Q),\\
\fn(\mathsf{collapse}_{a,k}\;c.P)&=\{c\}\cup\fn(P),\qquad
\fn(\mathsf{await}_a\;c\;\omega.P)=\{c,\omega\}\cup\fn(P),\\
\fn(\mathsf{await}_a\;c\;\omega.P\;\mathsf{within}_b\;t\;\mathsf{else}\;Q)
&=\{c,\omega\}\cup\fn(P)\cup\fn(Q).
\end{aligned}
\]
Expression substitution $e[v/x]$ and name substitution $P[d/c]$ are
capture avoiding. Whenever a substitution would capture a free name, a
bound name is first replaced by a fresh one.

\begin{lemma}[Freshening]
For every finite set of names $N$ and term $M$, there exists
$M'\equiv_\alpha M$ whose bound names are disjoint from $N$.
\end{lemma}

\begin{proof}
The syntax tree of $M$ contains finitely many binders. Rename them one at a
time using names outside the union of $N$, the free names of $M$, and the
names already selected. The supply of names is countably infinite, while
the excluded set at each step is finite.
\end{proof}

\section{Structural congruence}

Structural congruence identifies processes whose difference is purely
administrative:
\[
\begin{gathered}
P\mid\mathbf 0\equiv P,
\qquad P\mid Q\equiv Q\mid P,
\qquad (P\mid Q)\mid R\equiv P\mid(Q\mid R),\\
P+Q\equiv Q+P,
\qquad !P\equiv P\mid\,!P,
\qquad (\nu c)(\nu d)P\equiv(\nu d)(\nu c)P,\\
(\nu c)\mathbf 0\equiv\mathbf 0,
\qquad
(\nu c)(P\mid Q)\equiv P\mid(\nu c)Q
\quad\text{if }c\notin\fn(P),
\end{gathered}
\]
together with $\alpha$-equivalence. No structural law may delete a
recorded event or identify distinct commitment identities. Structural
congruence rearranges a process before an event occurs; it does not
rewrite the history afterward.

\begin{theorem}[Equivariance]
For every capture-avoiding permutation of names $\rho$,
\[
C\xrightarrow{\alpha}C'
\quad\Longrightarrow\quad
\rho(C)\xrightarrow{\rho(\alpha)}\rho(C').
\]
\end{theorem}

\begin{proof}
By induction on the transition derivation of
Chapter~\ref{ch:lts}. Each primitive rule compares names only for
equality, freshness, or membership in a finite support. Permutations
preserve all three properties. The contextual cases are immediate from
the induction hypothesis, and the congruence case follows from closure of
structural congruence under capture-avoiding renaming.
\end{proof}

The theorem is stated for permutations rather than arbitrary renamings
because freshness is preserved by bijections and not in general by
many-to-one maps. It licenses every later proof to choose fresh names
without the result depending on a particular identifier generator.

\chapter{Concrete Syntax and Unique Parsing}
\label{ch:parsing}

Semantic disagreements must not hide inside parsing disagreements. If one
source text could be read as two processes, a theorem about processes
would say nothing definite about programs. This chapter fixes a concrete
syntax and proves that it is read one way only. The claim is made at
three levels of strength, and they are kept apart: lexical analysis is
deterministic; every accepted token sequence has exactly one derivation;
and a bounded check confirms that the implementation and printer agree
with that theory. The third does not stand in for the second.

\section{Lexical classes}

Tokens fall into disjoint classes, and each class is recognizable from its
first character:
\begin{center}
\small
\begin{tabular}{lll}
\toprule
Class & Form & Example\\
\midrule
commitment names ($\mathit{CN}$) & \texttt{\#}NAME & \texttt{\#c}\\
principals ($\mathit{PR}$) & \texttt{@}NAME & \texttt{@a}\\
obligations ($\mathit{OB}$) & \texttt{?}NAME & \texttt{?w}\\
observation rules ($\mathit{RL}$) & \texttt{\%}NAME & \texttt{\%alloc}\\
intervals ($\mathit{TM}$) & \texttt{\textasciitilde}NAME & \texttt{\textasciitilde t}\\
reasons ($\mathit{RS}$) & quoted, with backslash escapes & \texttt{"late"}\\
values, functions, requests ($\mathit{VAL}$) & lowercase-initial NAME, not a keyword & \texttt{f}\\
keywords and punctuation & fixed strings & \texttt{verify}, \texttt{|}\\
\bottomrule
\end{tabular}
\end{center}
Here NAME is an ASCII letter or underscore followed by ASCII letters,
digits, and underscores. The classes are chosen to make the sorts of the
abstract syntax visible in the text: a commitment name cannot stand where
a principal is expected, because the two cannot even be the same token.

\begin{lemma}[Unique tokenization]\label{lem:lex}
Every string has at most one tokenization.
\end{lemma}

\begin{proof}
After whitespace, the first character determines the class: the five
sigils, the quotation mark, a lowercase letter, or a punctuation
character, and these are pairwise distinct. Within a class the token is
the longest match, which is unique. A lowercase identifier is a keyword
exactly when it appears in the fixed keyword list. A quoted reason ends at
the first unescaped quotation mark.
\end{proof}

Failure certificates, witnesses, and ledger records have no concrete
syntax at all. They are runtime objects, produced by verifiers and
appended by rules, and no production of the grammar below mentions them.
That is a fact about the grammar's formation, not the outcome of a test:
no source program can contain a certificate, so none can forge one.

\section{The grammar}

The grammar is stratified by precedence, with choice loosest, parallel
composition tighter, and prefixes, restriction, and replication tightest:
\[
\begin{aligned}
P &::= S\;P_1 & P_1 &::= \texttt{+}\;P\;\mid\;\varepsilon\\
S &::= U\;S_1 & S_1 &::= \texttt{|}\;S\;\mid\;\varepsilon
\end{aligned}
\]
\[
\begin{aligned}
U ::= {}& \texttt{0}\mid \texttt{(}\,P\,\texttt{)}
\mid \texttt{nu}\;\mathit{CN}\;\texttt{.}\;U \mid \texttt{!}\;U
\mid \texttt{pop}\;\mathit{POP}\\
&\mid \texttt{bind}\;\texttt{[}\mathit{PR}\texttt{]}\;\mathit{CN}\;\texttt{to}\;\mathit{VAL}\;\texttt{.}\;U
\mid \texttt{transform}\;\mathit{CN}\;\texttt{with}\;\mathit{VAL}\;\texttt{.}\;U\\
&\mid \texttt{refuse}\;\texttt{[}\mathit{PR}\texttt{]}\;\mathit{CN}\;\texttt{because}\;\mathit{RS}\;\texttt{.}\;U
\mid \texttt{collapse}\;\texttt{[}\mathit{PR}\texttt{,}\,\mathit{RL}\texttt{]}\;\mathit{CN}\;\texttt{.}\;U\\
&\mid \texttt{verify}\;\texttt{[}\mathit{PR}\texttt{]}\;\mathit{CN}\;\texttt{using}\;\mathit{VAL}\;\texttt{then}\;U\;\texttt{else}\;U
\mid \texttt{await}\;\texttt{[}\mathit{PR}\texttt{]}\;\mathit{CN}\;\mathit{OB}\;\mathit{AW}\\
\mathit{POP} ::= {}& \mathit{VAL}\;\texttt{as}\;\mathit{CN}\;\texttt{.}\;U
\mid \texttt{[}\mathit{PR}\texttt{]}\;\mathit{VAL}\;\texttt{as}\;\mathit{CN}\;\texttt{after}\;\mathit{CN}\;\texttt{.}\;U\\
\mathit{AW} ::= {}& \texttt{.}\;U
\mid \texttt{within}\;\texttt{[}\mathit{PR}\texttt{]}\;\mathit{TM}\;\texttt{then}\;U\;\texttt{else}\;U
\end{aligned}
\]
Choice and parallel composition associate to the right, by the direction
of their recursion. Parentheses reset precedence.

One scoping decision needs stating because readers from other traditions
may expect otherwise. The continuation after a prefix's dot, and each
branch after \texttt{then} and \texttt{else}, is a unary process $U$, not an
arbitrary process. So
\begin{center}
\texttt{pop q as \#c . 0 | 0}\quad reads as\quad
$(\mathsf{pop}\;q\;\mathsf{as}\;c.\mathbf 0)\mid\mathbf 0$,
\end{center}
and a parallel continuation must be parenthesized,
\texttt{pop q as \#c . (0 | 0)}. The same holds for the branches of
verification and of a deadline. This is the usual convention of CCS, where
prefix binds more tightly than parallel composition, and it is what makes
the two-branch forms unambiguous: every \texttt{then} has exactly one
\texttt{else}, and no \texttt{else} can be claimed by more than one
\texttt{then}. The deadline form is distinguished from the plain
$\mathsf{await}$ by its second token after the obligation,
\texttt{within} rather than the dot, so there is no dangling-else
problem there either.

\section{Unique derivations}

For a nonterminal $X$ and one of its alternatives $\beta$, the
\emph{predict set} is $\mathrm{FIRST}(\beta)$, together with
$\mathrm{FOLLOW}(X)$ if $\beta$ can derive the empty string. A grammar is
LL(1) when, for each nonterminal, the predict sets of its alternatives
are pairwise disjoint. For the grammar above:
\begin{itemize}[leftmargin=1.5em]
\item $P_1$: the alternative $\texttt{+}\,P$ is predicted by \texttt{+};
the empty alternative by $\mathrm{FOLLOW}(P_1)=\mathrm{FOLLOW}(P)=
\{\texttt{)},\$\}$.
\item $S_1$: $\texttt{|}\,S$ is predicted by \texttt{|}; the empty
alternative by $\mathrm{FOLLOW}(S)=\{\texttt{+},\texttt{)},\$\}$.
\item $U$: each of the eleven alternatives begins with a distinct
terminal.
\item $\mathit{POP}$: $\mathit{VAL}$ versus \texttt{[}.
\item $\mathit{AW}$: \texttt{.} versus \texttt{within}.
\end{itemize}
The laboratory computes the same sets mechanically from the grammar and
finds no conflict (property P-ParserLL1).

\begin{theorem}[Unique parsing]\label{thm:parse}
Every token sequence in the language of the grammar has exactly one
derivation. Consequently, by Lemma~\ref{lem:lex}, every accepted source
text denotes exactly one abstract syntax tree.
\end{theorem}

\begin{proof}
Suppose a token sequence has two distinct leftmost derivations. Consider
the first step at which they differ. Both expand the same leftmost
nonterminal $X$, with different alternatives $\beta\ne\beta'$, and the
remaining input begins with the same token $t$. In each derivation, $t$
is either the first token derived from the chosen alternative, or, if the
alternative derives the empty string there, a token that follows $X$. So
$t$ lies in the predict sets of both $\beta$ and $\beta'$, contradicting
the LL(1) property. The conversion from derivation to abstract syntax
tree is a function, so the tree is unique too.
\end{proof}

Structural congruence plays no part in this theorem. The texts
\texttt{0 | pop q as \#c . 0} and \texttt{pop q as \#c . 0 | 0} have
different, unique trees; that those trees denote congruent processes is a
fact about the semantics, not an ambiguity of the syntax.

\section{Adequacy of the printer}

Let $\mathrm{show}(P)$ print a tree with the fewest parentheses its
precedence requires, and let $\mathrm{Concrete}(P)$ be the set of texts
obtained from $\mathrm{show}(P)$ by wrapping any subterms in additional,
redundant parentheses.

\begin{theorem}[Adequacy]\label{thm:adequacy}
$\mathrm{parse}(s)=P$ if and only if $s\in\mathrm{Concrete}(P)$.
\end{theorem}

\begin{proof}
From right to left, by induction on $P$: the minimal printing parses to
$P$ because it inserts parentheses exactly where the precedence
stratification requires them, and a redundant pair around a subterm is
the production $U\to\texttt{(}\,P\,\texttt{)}$, whose tree is the tree of
the enclosed process. From left to right, by induction on the unique
derivation of $s$: the conversion to a tree discards only parentheses and
the empty expansions of $P_1$ and $S_1$, so $s$ differs from
$\mathrm{show}(\mathrm{parse}(s))$ only by parentheses around subterms,
which is membership in $\mathrm{Concrete}(\mathrm{parse}(s))$.
\end{proof}

\section{The executable check}

The reference parser is generated from the same grammar: a table-driven
predictive parser whose table is the one computed for the LL(1) check,
followed by a conversion from parse trees to abstract syntax trees. There
is one grammar in the implementation, not two, so the grammar that is
analyzed is the grammar that is parsed.

Beyond the general property, the module checks bounded properties over
every abstract syntax tree with at most five constructors, $9{,}469$
distinct trees over a small alphabet of atoms. The alphabet can be small
because parsing depends only on token classes, not on which name or value
appears. For every such $P$:
\begin{itemize}[leftmargin=1.5em]
\item P-ParserUnique: $\mathrm{show}(P)$ has exactly one derivation,
counted by a chart that enumerates all derivations rather than following
the table;
\item P-ParserRoundTrip: $\mathrm{parse}(\mathrm{show}(P))=P$, and
$\mathrm{parse}$ of a fully redundantly parenthesized printing of $P$ is
also $P$;
\item P-ParserNormalize: $\mathrm{show}(\mathrm{parse}(\mathrm{show}(P)))
=\mathrm{show}(P)$, and $\mathrm{parse}(\mathrm{show}(\mathrm{parse}(s)))
=\mathrm{parse}(s)$ for the redundantly parenthesized $s$, so redundant
parentheses disappear without changing the tree.
\end{itemize}
Eleven named cases pin down the decisions of this chapter with expected
trees, among them prefix scope, the relative precedence of choice and
parallel composition, right associativity, the unary branches of
verification, the deadline form, and escaped reasons (P-ParserNamed).
Seventeen malformed inputs are rejected, among them an unquoted value in
the reason position, an unterminated reason, a keyword in a value
position, an empty parenthesized process, an extra \texttt{else}, a
malformed rule index, an obligation in a commitment position, and a
deadline lacking its clock principal (P-ParserRejects). As a negative
control, adding an ambiguous alternative to the grammar produces both an
LL(1) conflict and a sentence with two derivations; the checks detect the
ambiguity they are meant to detect.

These checks establish that the implementation and printer agree with the
theory over the enumerated class. The universal statement is
Theorem~\ref{thm:parse}, which rests on the LL(1) property and not on the
enumeration.

\chapter{The Lambda Layer}
\plan{Reduction and big-step evaluation $e\Downarrow v$ for the value
language, which the transformation and verification rules consume as a
premise. Compares the kernel encodings of abstraction and application with
the lambda layer's own, and asks what a history retains that a normal form
discards.}

\chapter{Combinators, Erasure, and Linear Resource}
\plan{Uses combinatory logic to test how much depends on names, and the
discarding combinator to ask whether erasure must be recorded. Introduces
linear logic as the resource discipline Part~IV needs, and compares kernel
$\Pop$'s removal of an option with linear consumption of a hypothesis.}

%===================================================================
\part{Processes and Labelled Transitions}
%===================================================================

\chapter{CCS: Synchronization and Binding}
\plan{A finite CCS fragment encoded into disposition histories and the
converse examined. The central question is whether binding is
synchronization or something weaker, since a CCS handshake consumes both
complementary actions in one step. The worked encoding of
Appendix~\ref{app:worked} belongs here.}

\chapter{CSP and the Two Refusals}
\plan{Sets failures-semantics refusal sets against recorded refusal. The
chapter's target is the weakened claim of the Preface: that recorded
refusal is not represented by the standard failures model without an
added internal event or instrumentation. It gives the instrumented
encoding and then tries to prove the separation, that uninstrumented
failures equivalence identifies histories differing only in recorded
refusals; the history twins of Chapter~\ref{ch:worlds} are the candidate
witnesses. The sequence of a certified verification failure followed by
a separately authorized refusal gives a sharp test case: a refusal whose
grounds are recorded evidence, for which a refusal set has no slot at
all.}

\chapter{The Pi-Calculus: Restriction and Fresh Commitments}
\label{ch:pi}
\plan{Name passing, scope extrusion, and restriction. Develops the remark
of Chapter~\ref{ch:lts} that restriction here scopes naming but not
visibility, and asks whether this makes the disposition calculus's
$(\nu c)$ closer to a capability discipline than to $\pi$-calculus
privacy.}

\chapter{Labelled Disposition Semantics}
\label{ch:lts}

\section{Actions and records}

Observable actions are
\[
\begin{aligned}
\alpha ::= {}& \mathsf{pop}(c,q)
\mid \mathsf{compensate}(a,c,q,d)
\mid \mathsf{bind}(a,c,b)
\mid \mathsf{transform}(c,f,b,v)\\
&\mid \mathsf{verify}(a,c,v,w)
\mid \mathsf{refuse}(a,c,r)
\mid \mathsf{collapse}(a,c,k,w,o).
\end{aligned}
\]
Every observable action has a canonical ledger record $\record(\alpha)$,
the record of Section~\ref{sec:records} with the capitalized constructor
and the same fields. This pairing is part of the semantics rather than an
optional logging policy. Subjects and projections are as defined there.

\section{Primitive rules}

POP introduces a fresh commitment:
\[
\frac{c\notin\dom(A)}{
\langle \mathsf{pop}\;q\;\mathsf{as}\;c.P,S,A,O,H\rangle
\xrightarrow{\mathsf{pop}(c,q)}
\langle P,S,A[c\mapsto\mathsf{open}(q)],O,
H\cdot\Pop(c,q)\rangle}.
\]

A compensating pop introduces a fresh successor to a terminal
commitment:
{\footnotesize
\[
\frac{c\notin\dom(A)\qquad \mathsf{terminal}(A(d))\qquad
      \mathsf{may}(a,\mathsf{compensate},d)}{
\langle \mathsf{pop}_a\;q\;\mathsf{as}\;c\;\mathsf{after}\;d.P,S,A,O,H\rangle
\xrightarrow{\mathsf{compensate}(a,c,q,d)}
\langle P,S,A[c\mapsto\mathsf{open}(q)],O,
H\cdot\Compensates(a,c,q,d)\rangle}.
\]
}

Binding advances an open commitment while preserving its identity:
\[
\frac{A(c)=\mathsf{open}(q)\qquad
      \mathsf{may}(a,\mathsf{bind},c)}{
\langle \mathsf{bind}_a\;c\;\mathsf{to}\;b.P,S,A,O,H\rangle
\xrightarrow{\mathsf{bind}(a,c,b)}
\langle P,S,A[c\mapsto\mathsf{bound}(q,b)],O,
H\cdot\Bind(a,c,b)\rangle}.
\]

Transformation evaluates a lambda-layer function and records both its
input and output:
{\footnotesize
\[
\frac{A(c)=\mathsf{bound}(q,b)\qquad f\,b\Downarrow v}{
\langle \mathsf{transform}\;c\;\mathsf{with}\;f.P,S,A,O,H\rangle
\xrightarrow{\mathsf{transform}(c,f,b,v)}
\langle P,S,A[c\mapsto\mathsf{transformed}(b,v)],O,
H\cdot\Transform(c,f,b,v)\rangle}.
\]
}

Verification runs the verifier and branches on its result. Write
\[
V\equiv\mathsf{verify}_a\;c\;\mathsf{using}\;g\;\mathsf{then}\;P\;
\mathsf{else}\;Q.
\]
The successful branch produces an indexed witness:
{\footnotesize
\[
\frac{A(c)=\mathsf{transformed}(b,v)\qquad
      \mathsf{may}(a,\mathsf{verify},c)\qquad
      g(c,v)\Downarrow\mathsf{ok}(w)}{
\langle V,S,A,O,H\rangle
\xrightarrow{\mathsf{verify}(a,c,v,w)}
\langle P,S,A[c\mapsto\mathsf{verified}(v,w)],O,
H\cdot\Verify(a,c,v,w)\rangle}.
\]
}
The failing branch records the failure certificate and leaves the
commitment live:
{\footnotesize
\[
\frac{A(c)=\mathsf{transformed}(b,v)\qquad
      \mathsf{may}(a,\mathsf{verify},c)\qquad
      g(c,v)\Downarrow\mathsf{fail}(\varphi)}{
\langle V,S,A,O,H\rangle
\xrightarrow{\mathsf{verificationFailed}(a,c,v,\varphi)}
\langle Q,S,A[c\mapsto\mathsf{verificationFailed}(v,\varphi)],O,
H\cdot\mathsf{VerificationFailed}(a,c,v,\varphi)\rangle}.
\]
}

Explicit refusal is available from any live state $d$:
\[
\frac{A(c)=d\qquad\mathsf{live}(d)\qquad
      \mathsf{may}(a,\mathsf{refuse},c)}{
\langle \mathsf{refuse}_a\;c\;\mathsf{because}\;r.P,S,A,O,H\rangle
\xrightarrow{\mathsf{refuse}(a,c,r)}
\langle P,S,A[c\mapsto\mathsf{refused}(r)],O,
H\cdot\Refuse(a,c,r)\rangle}.
\]

Collapse commits to an observation of a verified commitment, using the
witness held in its state:
{\footnotesize
\[
\frac{A(c)=\mathsf{verified}(v,w)\qquad
      \mathsf{may}(a,\mathsf{collapse},c)\qquad
      o=\mathsf{observe}_k(v)}{
\langle \mathsf{collapse}_{a,k}\;c.P,S,A,O,H\rangle
\xrightarrow{\mathsf{collapse}(a,c,k,w,o)}
\left\langle\begin{array}{@{}l@{}}P,S,A[c\mapsto\mathsf{collapsed}(o)],
O\cup\{(c,k,o)\},\\ H\cdot\Collapse(a,c,k,w,o)\end{array}\right\rangle}.
\]
}
The subscript $k$ names the observation rule, and the record carries it.
Collapse is therefore not one universal quotient imposed on every history.
Different applications may expose different distinctions, but every use
declares which rule did the identifying. The observable label of a
collapse is the pair $\langle k,o\rangle$, not the payload $o$ alone, so
equal payloads produced under different rules remain distinct:
$\langle k_1,o\rangle\neq\langle k_2,o\rangle$ when $k_1\neq k_2$. An
observer who wants to ignore the rule may project it away, but that
projection is itself a declared collapse rule, not a default. Observations
are stored as $(c,k,o)$, so two commitments collapsing to equal values also
remain distinguishable in $O$.

External obligations are discharged by the environment, with
authenticated evidence:
{\footnotesize
\[
\frac{\mathsf{live}(A(c))\qquad\mathsf{may}(a,\mathsf{discharge},\omega)\qquad
      \mathsf{checkCert}(\varphi,c,\omega,\mathsf{discharged})}{
\langle \mathsf{await}_a\;c\;\omega.P,S,A,O,H\rangle
\xrightarrow{\mathsf{discharge}(a,c,\omega,\varphi)}
\langle P,S,A,O,H\cdot\mathsf{Discharged}(a,c,\omega,\varphi)\rangle}.
\]
}
For the deadline form, discharge takes the first branch by the same rule,
and a certified timeout takes the second:
{\footnotesize
\[
\frac{\mathsf{live}(A(c))\qquad\mathsf{may}(b,\mathsf{clock},\omega)\qquad
      \mathsf{checkCert}(\varphi,c,\langle\omega,t\rangle,\mathsf{expired})}{
\langle \mathsf{await}_a\;c\;\omega.P\;\mathsf{within}_b\;t\;\mathsf{else}\;Q,
S,A,O,H\rangle
\xrightarrow{\mathsf{timeout}(b,c,\omega,t,\varphi)}
\langle Q,S,A,O,H\cdot\mathsf{TimedOut}(b,c,\omega,t,\varphi)\rangle}.
\]
}
The discharge record must be classified, since the book has just argued
that no fifth causal primitive is needed. It is \emph{kernel-derived}: a
guarded $\Bind\circ\Pop$ of an evidence-carrying proposal to the history of
$c$, the same shape as binding, transformation, and verification
(Chapter~\ref{ch:layers}), differing from them only in leaving the
lifecycle state where it was. It is not a primitive and not a mere
annotation, because it changes what the process may do next. Three
nearby things are \emph{not} ledger events and must stay distinct from
it: a pending obligation, which is process state (the $\mathsf{await}$
prefix is simply still there); a denied attempt to discharge, which
belongs at most on the annotation channel; and nonresponse, which is the
absence of any discharge. None of these ends the commitment. If a pending
obligation is to end it, an authorized principal must refuse it
explicitly. A timeout record is classified the same way as a discharge:
kernel-derived, lifecycle-neutral, and authenticated.

The stability hypothesis of Section~\ref{sec:records} applies to the
validity predicate, failure certificates, and observation maps used in
these rules.

\section{Contextual rules}

The primitive rules act on a prefix at the top of the process. Parallel
components share the store, commitment map, observation set, and history,
so a step of one component is a step of the whole:
\[
\frac{\langle P,S,A,O,H\rangle\xrightarrow{\alpha}
      \langle P',S',A',O',H'\rangle}{
\langle P\mid Q,S,A,O,H\rangle\xrightarrow{\alpha}
\langle P'\mid Q,S',A',O',H'\rangle}
\qquad
\frac{\langle P,S,A,O,H\rangle\xrightarrow{\alpha}
      \langle P',S',A',O',H'\rangle}{
\langle P+Q,S,A,O,H\rangle\xrightarrow{\alpha}
\langle P',S',A',O',H'\rangle}
\]
\[
\frac{\langle P,S,A,O,H\rangle\xrightarrow{\alpha}
      \langle P',S',A',O',H'\rangle}{
\langle (\nu c)P,S,A,O,H\rangle\xrightarrow{\alpha}
\langle (\nu c)P',S',A',O',H'\rangle}
\qquad
\frac{P\equiv P_1\quad
      \langle P_1,\ldots\rangle\xrightarrow{\alpha}\langle P_2,\ldots\rangle
      \quad P_2\equiv P'}{
\langle P,\ldots\rangle\xrightarrow{\alpha}\langle P',\ldots\rangle}
\]
Choice is resolved by whichever summand acts first. The unchosen summand
is neither evaluated nor refused: the rule for $P+Q$ requires a transition
of $P$ only, inspects no premise of $Q$, and records nothing about $Q$.
Internal choice resolution is therefore a different thing from
disposition refusal, and nothing about a discarded branch can later be
read off the ledger.

\begin{remark}[Restriction scopes naming, not visibility]
In the $\pi$-calculus, restriction makes a name private and hides actions
on it from the environment. Here the restriction rule places no side
condition on the label, and the record is appended to a history shared by
every component. Restriction here is lexical scope, not
confidentiality: it limits which parts of a process can \emph{name} a
commitment, and so which parts can act on it, but not who can \emph{see}
what happened to it. The ledger is not scoped. Whether this makes
restriction behave like a capability discipline is a question for
Chapter~\ref{ch:pi}, not a claim of this one.
\end{remark}

\section{Worked derivations}

\paragraph{The successful route.}
Let
\[
\begin{aligned}
P_0={}&\mathsf{pop}\;q\;\mathsf{as}\;c.\,
\mathsf{bind}_a\;c\;\mathsf{to}\;b.\,
\mathsf{transform}\;c\;\mathsf{with}\;f.\\
&\mathsf{verify}_{a'}\;c\;\mathsf{using}\;g\;\mathsf{then}\;
\mathsf{collapse}_{a'',k}\;c.\mathbf 0\;
\mathsf{else}\;\mathbf 0,
\end{aligned}
\]
with $f\,b\Downarrow v$, $g(c,v)\Downarrow\mathsf{ok}(w)$,
$o=\mathsf{observe}_k(v)$, and each principal holding the authority its
step requires. From the empty configuration the five primitive rules fire
in order:
\[
C_0\xrightarrow{\mathsf{pop}(c,q)}C_1
\xrightarrow{\mathsf{bind}(a,c,b)}C_2
\xrightarrow{\mathsf{transform}(c,f,b,v)}C_3
\xrightarrow{\mathsf{verify}(a',c,v,w)}C_4
\xrightarrow{\mathsf{collapse}(a'',c,k,w,o)}C_5,
\]
with
\[
\begin{array}{lll}
\toprule
i & A_i(c) & H_i\\
\midrule
1 & \mathsf{open}(q) & \langle\Pop(c,q)\rangle\\
2 & \mathsf{bound}(q,b) & H_1\cdot\Bind(a,c,b)\\
3 & \mathsf{transformed}(b,v) & H_2\cdot\Transform(c,f,b,v)\\
4 & \mathsf{verified}(v,w) & H_3\cdot\Verify(a',c,v,w)\\
5 & \mathsf{collapsed}(o) & H_4\cdot\Collapse(a'',c,k,w,o)\\
\bottomrule
\end{array}
\]
and $O_5=\{(c,k,o)\}$. Each history is a proper prefix of the next. The
premise of each rule is visible in the previous row: the bind rule needs
the $\mathsf{open}$ state of row~1, the transform rule reads the $b$ of
row~2, and so on.

\paragraph{Refusal at every live stage.}
The four refusal routes are not redundant, because they establish that the
refusal rule covers every live state, the failed state included.
\begin{enumerate}[leftmargin=1.5em]
\item \emph{After pop.} From $C_1$, $\mathsf{refuse}_a\;c\;\mathsf{because}\;r$
fires with $A_1(c)=\mathsf{open}(q)$, giving
$H_1\cdot\Refuse(a,c,r)$.
\item \emph{After binding.} From $C_2$, a reviewer with refusal authority
judges the bound value unsuitable; the explicit rule fires from
$\mathsf{bound}(q,b)$.
\item \emph{After an inadmissible transformation.} From $C_3$, the explicit
rule fires from $\mathsf{transformed}(b,v)$ with a reason stating why $v$ is
unacceptable.
\item \emph{After negative verification.} From $C_3$, if
$g(c,v)\Downarrow\mathsf{fail}(\varphi)$, the failing verification rule
fires and the process continues with its $\mathsf{else}$ branch, giving
$H_3\cdot\mathsf{VerificationFailed}(a',c,v,\varphi)$ with $c$ live in
$\mathsf{verificationFailed}(v,\varphi)$. The $\mathsf{else}$ branch then
contains the refusal, $\mathsf{refuse}_h\;c\;\mathsf{because}\;r$, by a
principal $h$ holding refusal authority, and the explicit refusal rule
fires from the failed state.
\end{enumerate}

\begin{remark}[Denial of an actor is not refusal of a commitment]
One tempting fifth route is ``refusal because the binding was
unauthorized''. There is no such rule, and there should not be:
\[
\neg\mathsf{may}(a,\mathsf{bind},c)\;\not\Rightarrow\;\mathsf{Refused}(c).
\]
If an unauthorized attempt refused $c$, authority failure would become an
ambient cancellation capability: any principal could end any commitment by
attempting something it is not entitled to do. Three outcomes must be kept
apart.
\begin{itemize}[leftmargin=1.5em]
\item $\mathsf{DeniedAttempt}(a,c,\alpha)$: a principal tried an act it
may not perform. Nothing about $c$ changes. If recorded at all, it belongs
on the kernel's non-causal annotation channel, an audit plane beside the
ledger, where it supports security analysis without becoming a causal
event.
\item $\mathsf{Blocked}(c,\omega)$: the commitment cannot advance until an
external obligation $\omega$ is discharged. It is a positive runtime form,
introduced in Chapter~\ref{ch:session}, and it changes nothing about the
lifecycle either.
\item $\mathsf{Refused}(b,c,r)$: an authorized principal $b$ ended the
commitment for reason $r$. Only this outcome changes the lifecycle.
\end{itemize}
In the present calculus an unauthorized attempt simply has no transition.
In the typed calculus it cannot even be written, since every attributed
prefix requires its authority as a typing premise.
\end{remark}

\paragraph{Independent steps commute.}
Let $A(c)=\mathsf{open}(q)$ and $A(c')=\mathsf{open}(q')$ with $c\neq c'$,
and let the process be
$\mathsf{bind}_a\;c\;\mathsf{to}\;b.\mathbf 0\mid
\mathsf{bind}_{a'}\;c'\;\mathsf{to}\;b'.\mathbf 0$. Either component may
act first, by the parallel rule (with commutativity of $\mid$ for the right
component). The two interleavings reach the same process, commitment map,
and observation set, and histories
\[
H\cdot\Bind(a,c,b)\cdot\Bind(a',c',b')
\quad\text{and}\quad
H\cdot\Bind(a',c',b')\cdot\Bind(a,c,b),
\]
which differ by one exchange of independent records. In general:

\begin{proposition}[Diamond for independent steps]\label{prop:diamond}
Let $C=\langle P_1\mid P_2,S,A,O,H\rangle$, and suppose $P_1$ can perform
$\alpha$ and $P_2$ can perform $\beta$ from $C$, with
$\record(\alpha)$ and $\record(\beta)$ independent in the sense of
Definition~\ref{def:dependence}. Then both interleavings are possible and
reach configurations equal in every component except the history, where
they differ by exchanging the two records.
\end{proposition}

\begin{proof}
The premises of each rule read the commitment map only at the subject of
its record, and, for compensation, at its predecessor. By independence,
the other step writes neither key (Theorem~\ref{thm:identity}), so each
premise still holds after the other step, and each update produces the
same value in either order. The process components evolve separately.
Observation sets are joined by union. A pop and a compensating pop each
require their subject to be fresh; independence gives distinct subjects,
so neither introduction uses the other's identity.
\end{proof}

The hypothesis that the two steps come from different parallel components
matters. In $\mathsf{bind}_a\;c\;\mathsf{to}\;b.\,
\mathsf{bind}_{a'}\;c'\;\mathsf{to}\;b'.\mathbf 0$ the records are still
independent, but only one order is possible, because the prefix imposes
it. Independence of records says replay does not care about their order;
it does not say the program allows both.

\paragraph{Competing terminal steps.}
Let $A(c)=\mathsf{verified}(v,w)$ and let the process be
$\mathsf{refuse}_a\;c\;\mathsf{because}\;r.\mathbf 0\mid
\mathsf{collapse}_{a',k}\;c.\mathbf 0$. Either component
may fire first. If the refusal fires, $c$ becomes $\mathsf{refused}(r)$ and
the collapse rule's premise $A(c)=\mathsf{verified}(v,w)$ fails. If the
collapse fires, $c$ becomes $\mathsf{collapsed}(o)$, which is not live, and
the refusal rule's premise fails. Exactly one terminal record is appended.

This settles a point about concurrency. Terminal exclusivity does not rest
on the order in which rules are written; it rests on the commitment map
being shared, so that the first terminal step falsifies the premise of
every competitor. That in turn rests on an assumption the interleaving
semantics makes silently.

\begin{hypothesis}[Atomic commitment map]\label{hyp:atomic}
Each transition reads the commitment map, checks its premises, and
updates the map as one indivisible step.
\end{hypothesis}

Under this hypothesis exclusivity is \emph{operational}. In a distributed
implementation the hypothesis must be earned: without a sequencer, a
compare-and-swap, a consensus object, or some other linearization point,
two replicas could each read a live state, each append a terminal record,
and produce incompatible dispositions that no later replay can reconcile.
Chapter~\ref{ch:sequencer} states what each mechanism provides. There are therefore two grades of the result: operational
exclusivity, which holds for every execution under
Hypothesis~\ref{hyp:atomic}, and static exclusion of competitors, which
holds for well-typed programs. What the untyped calculus does not do is say
anything about the loser, which is left with a prefix it can never perform and no
record of why. The linear discipline of Part~IV turns that situation into
a static error: a well-typed program cannot hold two rights to advance the
same commitment, so it cannot contain both components in the first place.

\section{First preservation results}

\begin{theorem}[Commitment identity]\label{thm:identity}
If $C\xrightarrow{\alpha}C'$, then $\dom(A)\subseteq\dom(A')$, and
\[
\dom(A')\setminus\dom(A)=
\begin{cases}
\{c\} & \text{if }\alpha=\mathsf{pop}(c,q)\text{ or }
         \alpha=\mathsf{compensate}(a,c,q,d),\\
\varnothing & \text{otherwise,}
\end{cases}
\]
with $c\notin\dom(A)$ in the first case. Moreover $A'(e)=A(e)$ for every
$e\neq\subj(\record(\alpha))$.
\end{theorem}

\begin{proof}
By induction on the derivation. Each primitive rule updates $A$ at no
more than one key, the subject of its record; the discharge and timeout
rules write back the value already there. The two introducing rules have
$c\notin\dom(A)$ as a premise and add that key; every other primitive rule
has a premise $A(c)=\ldots$ and so updates an existing key. The contextual
rules pass $A$ and $A'$ through unchanged from their premise.
\end{proof}

Since the domain never shrinks, a commitment identity once introduced
remains in the map for the rest of the execution. The spine's provision
for an archival operation that preserves identity in the ledger is
therefore not needed in the present calculus; it becomes relevant only if
a later extension removes terminal commitments from $A$ for efficiency.

\begin{theorem}[Lifecycle monotonicity]
If $C\to C'$ and $c$ occurs in both commitment maps, then
\[
  A(c)\preceq_D A'(c).
\]
In particular, no ordinary transition reopens a refused or collapsed
commitment.
\end{theorem}

\begin{proof}
By cases on the primitive rule at the root of the derivation, the
contextual rules transmitting the property. The introducing rules concern
a fresh identifier and, by Theorem~\ref{thm:identity}, leave every existing
key unchanged; in particular a compensating pop does not alter its
predecessor. Each advancing rule follows one generating edge of
$\preceq_D$. Both refusing rules follow the refusal edge from a live
state. Neither terminal state satisfies the premise of an advancing or
refusing rule.
\end{proof}

\begin{theorem}[History monotonicity]
If $C\xrightarrow{\alpha}C'$, then
\[
  H'=H\cdot\record(\alpha),
\]
and hence $H\sqsubseteq H'$.
\end{theorem}

\begin{proof}
Inspection of the primitive rules shows that each appends its paired
record. The contextual rules pass the history through unchanged from their
premise, and structural congruence acts on the process component alone. No
rule has a premise or conclusion in which a history suffix is removed or
replaced.
\end{proof}

\begin{corollary}[Terminal exclusivity, map form]
Along any execution, no commitment enters both $\mathsf{refused}$ and
$\mathsf{collapsed}$.
\end{corollary}

\begin{proof}
Whichever terminal state is entered first, lifecycle monotonicity keeps
the commitment in that state at every later configuration, since
$\preceq_D$ relates each terminal state only to itself. By
Theorem~\ref{thm:identity} the identity cannot be removed and
reintroduced in a different state.
\end{proof}

The history form of this corollary, which is what an auditor reading the
ledger needs, is Corollary~\ref{cor:terminal-history}.

\section{Authority}

Define the authority an action requires, $\mathsf{auth}(\alpha)$, by
\[
\begin{aligned}
\mathsf{auth}(\mathsf{bind}(a,c,b))&=\mathsf{bind},&
\mathsf{auth}(\mathsf{verify}(a,c,v,w))&=\mathsf{verify},\\
\mathsf{auth}(\mathsf{collapse}(a,c,k,w,o))&=\mathsf{collapse},&
\mathsf{auth}(\mathsf{compensate}(a,c,q,d))&=\mathsf{compensate},
\end{aligned}
\]
and
\[
\mathsf{auth}(\mathsf{refuse}(a,c,r))=\mathsf{refuse},\qquad
\mathsf{auth}(\mathsf{verificationFailed}(a,c,v,\varphi))=\mathsf{verify},
\]
$\mathsf{auth}(\mathsf{discharge}(a,c,\omega,\varphi))=\mathsf{discharge}$,
and $\mathsf{auth}(\mathsf{timeout}(b,c,\omega,t,\varphi))=\mathsf{clock}$.
The object of the required authority is the action's subject, except for
compensation, where it is the predecessor $d$, and discharge and timeout,
where it is the obligation $\omega$.

\begin{theorem}[Authorized transition]\label{thm:auth}
If $C\xrightarrow{\alpha}C'$ and $\alpha$ is attributed to principal $a$,
then $\mathsf{may}(a,\mathsf{auth}(\alpha),x)$ held in $C$, where $x$ is
the object of the required authority.
\end{theorem}

\begin{proof}
By inversion on the derivation. The contextual rules preserve the label,
so the root primitive rule determines it, and every attributed primitive
rule carries the $\mathsf{may}$ premise for exactly the kind
$\mathsf{auth}$ assigns to its label. In particular a refusal label comes
only from the explicit refusal rule, whose premise is refusal authority,
and a verification-failure label only from the failing verification rule,
whose premise is verification authority.
\end{proof}

Two different guarantees concern verification failures, and each has its
own scope.

\begin{proposition}[Source authenticity]\label{prop:vf}
If $\mathsf{VerificationFailed}(a,c,v,\varphi)$ occurs in a history
produced by execution, then at the configuration where it was appended,
$c$ was $\mathsf{transformed}(b,v)$ for some $b$, principal $a$ held
verification authority over $c$, and the verifier named in the process
returned $\mathsf{fail}(\varphi)$ on $(c,v)$.
\end{proposition}

\begin{proof}
By history monotonicity the record was appended by a single transition,
which can only have been an instance of the failing verification rule,
whose premises are the statement.
\end{proof}

This is a statement about programs. It says nothing about a ledger that
arrives from elsewhere, since someone with write access to storage could
append a record directly, without running any process. For that setting
the guarantee comes from replay.

\begin{proposition}[Ledger acceptance]\label{prop:vf-ledger}
If $\delta^{*}(H|_c)$ is defined and $H|_c$ contains the record
$\mathsf{VerificationFailed}(a,c,v,\varphi)$, then the preceding record
of $H|_c$ left $c$ in a state $\mathsf{transformed}(b,v)$ for which
\[
\mathsf{checkCert}(\varphi,c,v,\mathsf{fail})
\]
holds. Under Hypothesis~\ref{hyp:verifier}, $\varphi$ was issued by the
verifier it names, at the version it names, on that commitment and value.
\end{proposition}

\begin{proof}
The only clause of $\delta$ accepting the record is the
verification-failure clause, whose incoming state and side condition are
the statement. The second sentence is the authentication half of
Hypothesis~\ref{hyp:verifier}.
\end{proof}

What replay cannot check is that the principal $a$ in the record held
verification authority; that requires the policy, which an auditor can
supply alongside the ledger. And the proposition is only as strong as the
certificate scheme: if certificates can be forged, so can ledger-level
verification failures. Neither proposition concerns refusal at all. A
verification failure is evidence; whether it leads to refusal is decided
by a separate, separately authorized act.

In the present calculus the policy judgement $\mathsf{may}$ is fixed, so
the spine's non-amplification corollary holds trivially: no transition
changes what any principal may do. It becomes a substantive result once
delegation is added in Chapter~\ref{ch:authority}, where authority can grow, and the claim
is that it grows only through a delegation record.

\chapter{Choice, Actors, and the Arbiter}
\label{ch:choice}
\plan{The known separation results for mixed choice, and where a choice
resolved by a recording arbiter falls relative to them. The choice rule of
Chapter~\ref{ch:lts} leaves no record of the discarded summand; the
chapter asks whether a variant that records the road not taken is more
expressive, or only more informative. Contrasts the actor model's
unordered arrival with total causal order.}

%===================================================================
\part{Linear and Session Types}
%===================================================================

\chapter{The Type Algebra}
\label{ch:types}

The untyped calculus says what happens when a step is enabled. It says
nothing about programs that attempt steps that can never be enabled: a
collapse of an open commitment, a second binding of the same commitment
from two parallel components, a refusal by a principal without authority.
Such programs simply stop, with nothing recorded. The type algebra
excludes them. It does so by making the right to advance a commitment a
linear resource, indexed by the commitment's identity and lifecycle state,
and by stating exactly how those resources correspond to the runtime
commitment map.

\section{Types, rights, and facts}

The value types are
\[
A,B ::= \mathbf 1\mid A\times B\mid A+B\mid A\to B
\mid A\otimes B\mid A\multimap B\mid\;!A .
\]
On top of them the disposition layer uses three kinds of assumption.
\begin{itemize}[leftmargin=1.5em]
\item \emph{Rights} $\Commit\langle c,A,s\rangle$, for a live state $s$:
one of $\mathsf{Open}$, $\mathsf{Bound}$, $\mathsf{Transformed}$,
$\mathsf{Verified}$, and $\mathsf{VerificationFailed}$. A right is the
permission to advance commitment $c$ from state $s$. Rights are linear: they may be neither duplicated nor
discarded.
\item \emph{Tokens} $\mathsf{Refusal}\langle c,A,s\rangle$ and
$\mathsf{Observed}\langle c,B,k\rangle$, the results of the two terminal
steps. The index $s$ of a refusal token records the stage at which refusal
occurred. Tokens are affine: they may be discarded but not duplicated.
\item \emph{Facts}: the authority policy $\mathsf{may}(a,\kappa,x)$,
failure evidence $\mathsf{FailureEvidence}\langle c,B\rangle$ recording
that verification of $c$ returned a certified negative result, and
terminal facts $\mathsf{Terminal}\langle c,t\rangle$ with
$t\in\{\mathsf{Refused},\mathsf{Collapsed}\}$. Facts are persistent. A
terminal fact is issued by the terminal step of $c$ and may be consulted
any number of times afterward, which is sound because terminality is
monotone (lifecycle monotonicity).
\end{itemize}
The linear context holds the right to advance a commitment, never the
historical fact that the commitment exists. The history and the
commitment map keep $c$ forever (Theorem~\ref{thm:identity}); the context
holds at most one right to move it.

The typing judgement for processes is
\[
\Sigma;\Gamma;\Delta\vdash P ,
\]
where $\Sigma$ is a set of names (commitment identities and obligations),
$\Gamma$ contains value typings and facts, and $\Delta$ contains rights and
tokens. The protocol $\Theta$ offered by $P$ is, in this edition,
determined by $\Delta$: each right $\Commit\langle c,A,s\rangle$ offers
the row of the refusal completion table for $s$, read as a protocol
(Section~\ref{sec:fidelity}). The general judgement
$\Sigma;\Gamma;\Delta\vdash P\triangleright\Theta$ with independently
specified protocols is needed only for interaction among several
principals over shared channels, which this calculus does not yet have.

\section{Operator types}

Read as functions, the disposition operators have these types.
\[
\begin{aligned}
\mathsf{bind}&:\Commit\langle c,A,\mathsf{Open}\rangle\otimes A
\multimap\Commit\langle c,A,\mathsf{Bound}\rangle,\\
\mathsf{transform}&:\Commit\langle c,A,\mathsf{Bound}\rangle
\otimes(A\to B)
\multimap\Commit\langle c,B,\mathsf{Transformed}\rangle,\\
\mathsf{verify}&:\Commit\langle c,B,\mathsf{Transformed}\rangle
\otimes\mathsf{Verifier}\langle c,B\rangle\\
&\qquad\multimap
\bigl(\exists w.\,\Commit\langle c,B,\mathsf{Verified}(w)\rangle
\otimes\mathsf{Evidence}\langle c,B,w\rangle\bigr)\\
&\qquad\oplus\;
\bigl(\exists\varphi.\,\Commit\langle c,B,\mathsf{VerificationFailed}(\varphi)\rangle
\otimes\,!\mathsf{FailureEvidence}\langle c,B,\varphi\rangle\bigr),\\
\mathsf{refuse}_s&:\Commit\langle c,A,s\rangle\otimes\mathsf{Reason}\\
&\qquad\multimap\mathsf{Refusal}\langle c,A,s\rangle
\otimes\,!\mathsf{Terminal}\langle c,\mathsf{Refused}\rangle,\\
\mathsf{collapse}_k&:\Commit\langle c,B,\mathsf{Verified}(w)\rangle\\
&\qquad\multimap\mathsf{Observed}\langle c,B,k\rangle
\otimes\,!\mathsf{Terminal}\langle c,\mathsf{Collapsed}\rangle.
\end{aligned}
\]
Four features deserve comment.

\emph{Pop is generative, not an ordinary existential.} Its external type
may be written $A\to\exists c.\Commit\langle c,A,\mathsf{Open}\rangle$,
but internally it is a name-generation rule analogous to restriction: its
typing premise extends $\Sigma$ with a name not already in it, and the
continuation receives the open right for that name (rule \textsc{T-Pop}
below). An ordinary existential could be eliminated and its witness
compared with other names; a generated name is fresh by construction, and
that is what freshness of commitment identities rests on.

\emph{The success branch of verification carries its witness.} The
verified state is indexed by the witness $w$, and the evidence
$\mathsf{Evidence}\langle c,B,w\rangle$ is indexed by the commitment.
Evidence for one commitment is never evidence for another: the index is
part of the type, and under Hypothesis~\ref{hyp:index} it is part of the
predicate too. In process typing the witness does not need a name, because
$\mathsf{collapse}$ reads it from the runtime state; what the typing
guarantees is that such a witness exists there (Lemma~\ref{lem:canon}).

\emph{Both verification branches keep the right.} The positive branch
produces a verified right with its evidence; the negative branch produces
a right in $\mathsf{VerificationFailed}$ with persistent failure evidence.
Neither produces a refusal token or a terminal fact, because neither is a
disposition. A verifier cannot end a commitment, and the type says so.

\emph{Terminal steps issue persistent facts.} Refusal and collapse consume
the right and produce two things: an affine token recording the result,
and a persistent fact $\mathsf{Terminal}\langle c,t\rangle$. The fact is
what compensation consults. The earlier worry, that compensation needed a
certificate of terminality that no rule issued, is resolved by making the
terminal rules issue it.

\section{Context algebra}

A \emph{linear context} $\Delta$ is a finite multiset of rights and
tokens. It is \emph{well formed} when it contains at most one right for
each commitment, at most one token for each commitment, and never both a
right and a token for the same commitment. Write $\Delta_1\uplus\Delta_2$
for multiset union, defined only when the result is well formed.

\begin{lemma}[Ownership]\label{lem:own}
If $\Delta_1\uplus\Delta_2$ is defined and $\Delta_1$ contains a right for
$c$, then $\Delta_2$ contains neither a right nor a token for $c$.
\end{lemma}

\begin{proof}
Otherwise the union would contain two assumptions for $c$ of kinds that
well-formedness forbids together.
\end{proof}

For an action $\alpha$ with subject $c$, the \emph{context step}
$\Delta\xrightarrow{\alpha}\Delta'$ is defined by: a pop or compensation
adds $\Commit\langle c,A,\mathsf{Open}\rangle$; binding, transformation,
and successful verification replace the right for $c$ by the right for its
successor state; either kind of refusal replaces the right by the token
$\mathsf{Refusal}\langle c,A,s\rangle$; collapse replaces the right by
$\mathsf{Observed}\langle c,B,k\rangle$; a discharge leaves $\Delta$
unchanged.

\begin{lemma}[Recombination]\label{lem:recomb}
If $\Delta_1\uplus\Delta_2$ is defined, $\Delta_1\xrightarrow{\alpha}
\Delta_1'$, and either $\alpha$ introduces a name not mentioned in
$\Delta_2$ or $\Delta_1$ contains the right for the subject of $\alpha$,
then $\Delta_1'\uplus\Delta_2$ is defined and equals
$(\Delta_1\uplus\Delta_2)'$, the context step applied to the union.
\end{lemma}

\begin{proof}
The context step changes only the entry for the subject of $\alpha$. In
the first case that entry is new to both sides. In the second, by
Lemma~\ref{lem:own}, $\Delta_2$ has no entry for it. Either way the
changed entry lives in $\Delta_1'$ alone and well-formedness of the union
is preserved.
\end{proof}

The unrestricted context $\Gamma$ admits weakening, contraction, and
exchange. The linear context admits exchange only, together with
weakening by tokens.

\section{Typing the value layer}

Values are typed by $\Gamma;\Delta\vdash e:A$. The rules for the
multiplicative core are
\[
\begin{gathered}
\frac{}{\Gamma,x{:}A;\varnothing\vdash x:A}
\qquad
\frac{}{\Gamma;x{:}A\vdash x:A}
\qquad
\frac{\Gamma;\Delta,x{:}A\vdash e:B}{\Gamma;\Delta\vdash\lambda x{:}A.e:A\multimap B}
\qquad
\frac{\Gamma,x{:}A;\varnothing\vdash e:B}{\Gamma;\varnothing\vdash\lambda x{:}A.e:A\to B}
\\[1ex]
\frac{\Gamma;\Delta_1\vdash e_1:A\multimap B\qquad\Gamma;\Delta_2\vdash e_2:A}
{\Gamma;\Delta_1\uplus\Delta_2\vdash e_1\,e_2:B}
\qquad
\frac{\Gamma;\Delta\vdash e_1:A\to B\qquad\Gamma;\varnothing\vdash e_2:A}
{\Gamma;\Delta\vdash e_1\,e_2:B}
\end{gathered}
\]
with products, sums, projections, injections, and case analysis typed in
the unrestricted fragment in the standard way. Transformations and
verifiers used by processes are closed unrestricted terms.

\begin{lemma}[Unrestricted substitution]\label{lem:usubst}
If $\Gamma,x{:}A;\Delta\vdash e:B$ and $\Gamma;\varnothing\vdash v:A$,
then $\Gamma;\Delta\vdash e[v/x]:B$.
\end{lemma}

\begin{proof}
Induction on the derivation of $e$. At an unrestricted variable $x$ the
result is the typing of $v$, weakened in $\Gamma$. At a rule that splits
$\Delta$, the variable $x$ lives in $\Gamma$ and so is available on both
sides, and the induction hypothesis applies to each. At a binder, rename
the bound variable apart from $x$ and the free variables of $v$ first.
\end{proof}

\begin{lemma}[Linear substitution]\label{lem:lsubst}
If $\Gamma;\Delta_1,x{:}A\vdash e:B$ and $\Gamma;\Delta_2\vdash v:A$, and
$\Delta_1\uplus\Delta_2$ is defined, then
$\Gamma;\Delta_1\uplus\Delta_2\vdash e[v/x]:B$.
\end{lemma}

\begin{proof}
Induction on the derivation of $e$. At a linear variable, $e=x$ and
$\Delta_1=\varnothing$, so the result is the typing of $v$. The linear
variable $x$ cannot occur at an unrestricted variable rule or in the body
of an unrestricted abstraction, since both require an empty linear
context. At linear application, $\Delta_1,x{:}A$ is split between the two
premises and $x$ lands in exactly one of them; apply the induction
hypothesis to that premise and recombine, the resources of $v$ joining the
side that held $x$. At unrestricted application the argument has an empty
linear context, so $x$ is in the function premise. At linear abstraction,
rename apart and apply the induction hypothesis. In no case are the
resources of $v$ duplicated, because $x$ occurs in exactly one branch of
every split.
\end{proof}

Evaluation $e\Downarrow v$ preserves types for closed terms, by induction
on the evaluation derivation using the two substitution lemmas; closed
values of function type are abstractions, of product type pairs, and of
sum type injections. These canonical-forms facts are standard
\cite{pierce} and are used without further proof.

\section{Typing processes}

The process rules are syntax directed. Write $\Gamma\vdash e:A$ for
$\Gamma;\varnothing\vdash e:A$.
{\small
\[
\begin{gathered}
\textsc{T-Nil}\;
\frac{\Delta\text{ contains only tokens}}{\Sigma;\Gamma;\Delta\vdash\mathbf 0}
\qquad
\textsc{T-Pop}\;
\frac{c\notin\Sigma\quad\Gamma\vdash q:A\quad
\Sigma,c;\Gamma;\Delta,\Commit\langle c,A,\mathsf{Open}\rangle\vdash P}
{\Sigma;\Gamma;\Delta\vdash\mathsf{pop}\;q\;\mathsf{as}\;c.P}
\\[1.2ex]
\textsc{T-Comp}\;
\frac{\begin{array}{c}c\notin\Sigma\quad\mathsf{Terminal}\langle d,t\rangle\in\Gamma\quad
\mathsf{may}(a,\mathsf{compensate},d)\in\Gamma\quad\Gamma\vdash q:A\\
\Sigma,c;\Gamma;\Delta,\Commit\langle c,A,\mathsf{Open}\rangle\vdash P\end{array}}
{\Sigma;\Gamma;\Delta\vdash\mathsf{pop}_a\;q\;\mathsf{as}\;c\;\mathsf{after}\;d.P}
\\[1.2ex]
\textsc{T-Bind}\;
\frac{\mathsf{may}(a,\mathsf{bind},c)\in\Gamma\quad\Gamma\vdash b:A\quad
\Sigma;\Gamma;\Delta,\Commit\langle c,A,\mathsf{Bound}\rangle\vdash P}
{\Sigma;\Gamma;\Delta,\Commit\langle c,A,\mathsf{Open}\rangle\vdash
\mathsf{bind}_a\;c\;\mathsf{to}\;b.P}
\\[1.2ex]
\textsc{T-Trans}\;
\frac{\Gamma\vdash f:A\to B\quad
\Sigma;\Gamma;\Delta,\Commit\langle c,B,\mathsf{Transformed}\rangle\vdash P}
{\Sigma;\Gamma;\Delta,\Commit\langle c,A,\mathsf{Bound}\rangle\vdash
\mathsf{transform}\;c\;\mathsf{with}\;f.P}
\\[1.2ex]
\textsc{T-Ver}\;
\frac{\begin{array}{c}
\mathsf{may}(a,\mathsf{verify},c)\in\Gamma\qquad
\Gamma\vdash g:\mathsf{Verifier}\langle c,B\rangle\\
\Sigma;\Gamma;\Delta,\Commit\langle c,B,\mathsf{Verified}\rangle\vdash P\\
\Sigma;\Gamma,\mathsf{FailureEvidence}\langle c,B\rangle;
\Delta,\Commit\langle c,B,\mathsf{VerificationFailed}\rangle\vdash Q
\end{array}}
{\Sigma;\Gamma;\Delta,\Commit\langle c,B,\mathsf{Transformed}\rangle\vdash
\mathsf{verify}_a\;c\;\mathsf{using}\;g\;\mathsf{then}\;P\;\mathsf{else}\;Q}
\\[1.2ex]
\textsc{T-Ref}\;
\frac{\begin{array}{c}\mathsf{live}(s)\quad\mathsf{may}(a,\mathsf{refuse},c)\in\Gamma\quad
\Gamma\vdash r:\mathsf{Reason}\\
\Sigma;\Gamma,\mathsf{Terminal}\langle c,\mathsf{Refused}\rangle;
\Delta,\mathsf{Refusal}\langle c,A,s\rangle\vdash P\end{array}}
{\Sigma;\Gamma;\Delta,\Commit\langle c,A,s\rangle\vdash
\mathsf{refuse}_a\;c\;\mathsf{because}\;r.P}
\\[1.2ex]
\textsc{T-Col}\;
\frac{\begin{array}{c}\mathsf{may}(a,\mathsf{collapse},c)\in\Gamma\quad
k:\mathsf{Rule}(B)\\
\Sigma;\Gamma,\mathsf{Terminal}\langle c,\mathsf{Collapsed}\rangle;
\Delta,\mathsf{Observed}\langle c,B,k\rangle\vdash P\end{array}}
{\Sigma;\Gamma;\Delta,\Commit\langle c,B,\mathsf{Verified}\rangle\vdash
\mathsf{collapse}_{a,k}\;c.P}
\\[1.2ex]
\textsc{T-Await}\;
\frac{\omega\in\Sigma\quad\mathsf{may}(a,\mathsf{discharge},\omega)\in\Gamma\quad
\Sigma;\Gamma;\Delta,\Commit\langle c,A,s\rangle\vdash P}
{\Sigma;\Gamma;\Delta,\Commit\langle c,A,s\rangle\vdash\mathsf{await}_a\;c\;\omega.P}
\\[1.2ex]
\textsc{T-Deadline}\;
\frac{\begin{array}{c}\omega\in\Sigma\quad
\mathsf{may}(a,\mathsf{discharge},\omega),\ \mathsf{may}(b,\mathsf{clock},\omega)\in\Gamma\\
\Sigma;\Gamma;\Delta,\Commit\langle c,A,s\rangle\vdash P\qquad
\Sigma;\Gamma;\Delta,\Commit\langle c,A,s\rangle\vdash Q\end{array}}
{\Sigma;\Gamma;\Delta,\Commit\langle c,A,s\rangle\vdash
\mathsf{await}_a\;c\;\omega.P\;\mathsf{within}_b\;t\;\mathsf{else}\;Q}
\\[1.2ex]
\textsc{T-Par}\;
\frac{\Sigma;\Gamma;\Delta_1\vdash P\quad\Sigma;\Gamma;\Delta_2\vdash Q}
{\Sigma;\Gamma;\Delta_1\uplus\Delta_2\vdash P\mid Q}
\qquad
\textsc{T-Sum}\;
\frac{\Sigma;\Gamma;\Delta\vdash P\quad\Sigma;\Gamma;\Delta\vdash Q}
{\Sigma;\Gamma;\Delta\vdash P+Q}
\\[1.2ex]
\textsc{T-Res}\;
\frac{c\notin\Sigma\quad\Sigma,c;\Gamma;\Delta\vdash P}
{\Sigma;\Gamma;\Delta\vdash(\nu c)P}
\qquad
\textsc{T-Rep}\;
\frac{\Sigma;\Gamma;\varnothing\vdash P}{\Sigma;\Gamma;\varnothing\vdash\,!P}
\end{gathered}
\]
}
Here $\Gamma\vdash g:\mathsf{Verifier}\langle c,B\rangle$ says $g$ is a
verifier whose witnesses and certificates are issued for $c$, and
$k:\mathsf{Rule}(B)$ says $\mathsf{observe}_k$ is total on $B$.

Several design decisions are visible in the rules. \textsc{T-Nil} forbids
ending a process while it holds a right: a well-typed program cannot
silently abandon a live commitment. \textsc{T-Sum} types both summands
against the same context, so a choice is between alternative uses of the
same right. \textsc{T-Rep} requires an empty linear context, so no right
is ever replicated. The terminal rules add the terminal fact to $\Gamma$
for the continuation, and \textsc{T-Comp} consults it. Every attributed
prefix requires its authority fact, so an unauthorized attempt cannot be
written in a well-typed program at all.

\section{Configurations and the rights--state correspondence}

A typing of the process says nothing, by itself, about the runtime. The
bridge between them is the following invariant. Without it the
commitment map and the linear context could each be internally consistent
while disagreeing with each other.

\begin{definition}[Rights agree]\label{def:rights}
$\mathsf{RightsAgree}(A,\Delta)$ holds when, for every commitment $c$:
\begin{enumerate}[label=(\roman*)]
\item if $A(c)$ is live with state kind $s$, then $\Delta$ contains exactly
one right for $c$, namely $\Commit\langle c,A_c,s\rangle$, and the values
held in $A(c)$ have the types that right indexes;
\item if $A(c)$ is terminal or $c\notin\dom(A)$, then $\Delta$ contains no
right for $c$;
\item a token for $c$ in $\Delta$ occurs only if $A(c)$ is terminal, and a
refusal token only if $A(c)$ is refused, an observation token only if it
is collapsed.
\end{enumerate}
\end{definition}

\begin{definition}[Well-typed configuration]\label{def:wtconf}
$\Sigma;\Gamma;\Delta\vdash\langle P,S,A,O,H\rangle\;\mathsf{ok}$ when
\begin{enumerate}[label=(C\arabic*)]
\item $\Sigma;\Gamma;\Delta\vdash P$;
\item the configuration is consistent: $\replay(H)=\langle A,O\rangle$;
\item $\mathsf{RightsAgree}(A,\Delta)$;
\item every fact in $\Gamma$ is true: the authority facts are the policy,
and $\mathsf{Terminal}\langle c,t\rangle\in\Gamma$ only if $A(c)$ is
terminal of kind $t$;
\item $\dom(A)\subseteq\Sigma$, and the names occurring free in $P$ are
in $\Sigma$.
\end{enumerate}
\end{definition}

\begin{lemma}[Canonical forms for rights]\label{lem:canon}
In a well-typed configuration, if $\Delta$ contains
$\Commit\langle c,B,\mathsf{Verified}\rangle$, then
$A(c)=\mathsf{verified}(v,w)$ with $v$ of type $B$ and $\valid(w,c,v)$. If
$\Delta$ contains $\Commit\langle c,B,\mathsf{VerificationFailed}\rangle$,
then $A(c)=\mathsf{verificationFailed}(v,\varphi)$ with
$\mathsf{checkCert}(\varphi,c,v,\mathsf{fail})$. If $\Delta$ contains
$\Commit\langle c,A,s\rangle$ for any live $s$, then $A(c)$ is in state
$s$ and $H|_c$ is a $\delta$-path from an opening record to $A(c)$.
\end{lemma}

\begin{proof}
By (C3) the state of $c$ is $s$ with correctly typed contents. By (C2) and
the definition of replay, $A(c)=\delta^{*}(H|_c)$, which is a path; for
the verified case the verification clause of $\delta$ requires
$\valid(w,c,v)$, for the failed case the failure clause requires the
certificate check, and by stability both still hold.
\end{proof}

Canonical forms is where the evidence discipline becomes a static
guarantee: possession of a verified right is enough to know that a
correctly indexed witness exists in the state, without the program ever
having named it.

\section{The proof target}

The metatheory of the next chapter establishes preservation, progress, and
protocol fidelity. Their synthesis is the following statement, proved in
Chapter~\ref{ch:typedsafety}.

\begin{theorem}[Typed disposition safety, target form]
Suppose $\Sigma;\Gamma;\Delta\vdash C_0\;\mathsf{ok}$. For every finite
execution $C_0\xRightarrow{\tau}C_n$, the configuration $C_n$ is well
typed; every live commitment has exactly one right to advance it and every
terminal commitment has none; every attributed transition was authorized;
every witness is indexed by the commitment and value it verifies; every
collapse has a complete verified ledger path; every refusal remains in the
history; no commitment has two terminal dispositions; and no well-typed
program contains a competing continuation for a commitment whose right it
does not hold.
\end{theorem}

\chapter{Preservation, Progress, and Protocol Fidelity}
\label{ch:session}

This chapter proves that typing is preserved by execution, that a
well-typed configuration is never silently stuck, and that every execution
follows the lifecycle protocol of each commitment it touches. It closes
with the negative derivations: the programs the type system rejects, and
why.

\section{Structural lemmas}

\begin{lemma}[Inversion]\label{lem:inv}
Each typing rule for a prefix is the only rule whose conclusion has that
prefix at the root. Hence from a derivation of
$\Sigma;\Gamma;\Delta\vdash\pi.P$ the premises of the corresponding rule
can be read off; in particular $\Delta$ contains the right the rule
consumes.
\end{lemma}

\begin{proof}
The rules are syntax directed; only \textsc{T-Nil}, \textsc{T-Par},
\textsc{T-Sum}, \textsc{T-Res}, and \textsc{T-Rep} apply to non-prefix
forms, and each applies to exactly one form.
\end{proof}

\begin{lemma}[Weakening and strengthening]\label{lem:weak}
Typing is preserved by adding true facts to $\Gamma$, by adding tokens to
$\Delta$, and by adding to or removing from $\Sigma$ names that occur
neither free in $P$ nor in $\Delta$.
\end{lemma}

\begin{proof}
Induction on the derivation. Facts are only ever consulted, tokens are
absorbed by \textsc{T-Nil}, and names outside the free names are never
inspected except by freshness premises, which remain satisfiable by
choosing binders apart.
\end{proof}

\begin{lemma}[Equivariance of typing]\label{lem:tequiv}
For every permutation $\rho$ of names,
$\Sigma;\Gamma;\Delta\vdash P$ implies
$\rho\Sigma;\rho\Gamma;\rho\Delta\vdash\rho P$.
\end{lemma}

\begin{proof}
Induction on the derivation; every premise that mentions names compares
them only for equality or membership, which permutations preserve.
\end{proof}

\begin{lemma}[Congruence]\label{lem:congr}
If $P\equiv Q$ then $\Sigma;\Gamma;\Delta\vdash P$ iff
$\Sigma;\Gamma;\Delta\vdash Q$.
\end{lemma}

\begin{proof}
By cases on the axioms, each in both directions. The laws for $\mid$ are
commutativity and associativity of $\uplus$ together with \textsc{T-Nil}
for the unit. Commutativity of $+$ exchanges the two identical-context
premises of \textsc{T-Sum}. For $!P\equiv P\mid\,!P$, \textsc{T-Rep} gives
an empty context, and $\varnothing\uplus\varnothing=\varnothing$. The
restriction laws use Lemma~\ref{lem:weak} for names not free, and
$\alpha$-conversion uses Lemma~\ref{lem:tequiv}.
\end{proof}

\section{Preservation}

\begin{theorem}[Preservation]\label{thm:sr}
If $\Sigma;\Gamma;\Delta\vdash C\;\mathsf{ok}$ and
$C\xrightarrow{\alpha}C'$, then there are $\Sigma'\supseteq\Sigma$ and
$\Gamma'\supseteq\Gamma$ such that
$\Sigma';\Gamma';\Delta'\vdash C'\;\mathsf{ok}$, where
$\Delta\xrightarrow{\alpha}\Delta'$ is the context step.
\end{theorem}

\begin{proof}
By induction on the derivation of the transition, with the process
typing brought into matching shape by Lemma~\ref{lem:congr}.

\emph{Primitive rules.} By Lemma~\ref{lem:inv} the prefix's typing
premises are available.
\begin{itemize}[leftmargin=1.5em]
\item \emph{Pop.} The premise of \textsc{T-Pop} types $P$ under
$\Sigma,c$ with the open right; the operational rule chose $c$ outside
$\dom(A)$, and by Lemma~\ref{lem:tequiv} we may take the binder to be that
name. (C2) holds since $\delta(\bot,\Pop(c,q))=\mathsf{open}(q)$ and
$c\notin\dom(A)$ means $H|_c$ was empty. (C3) holds with the new right,
and $q:A$ by the premise.
\item \emph{Compensating pop.} As for pop, with the terminal fact for $d$
in $\Gamma$ matching the operational premise by (C4).
\item \emph{Bind, transform, successful verification.} Inversion gives
the right for the prior state in $\Delta$ and types the continuation with
the right for the next state. By (C3) the runtime state is the prior one,
the operational rule moves it to the next, and $\delta$ agrees, so (C2)
and (C3) hold. For transformation, $f\,b\Downarrow v$ with $f:A\to B$ and
$b:A$ gives $v:B$ by type preservation of evaluation. For verification,
$v$ keeps its type and the witness is valid by
Hypothesis~\ref{hyp:verifier}.
\item \emph{Failed verification.} Inversion of \textsc{T-Ver} types $Q$
with the right for $c$ in $\mathsf{VerificationFailed}$ and the persistent
fact $\mathsf{FailureEvidence}\langle c,B\rangle$. The operational rule
moves $c$ to $\mathsf{verificationFailed}(v,\varphi)$ and $\delta$ agrees,
so (C2) and (C3) hold, and the fact is true in $C'$, so (C4) holds.
\item \emph{Refusal.} The right is
replaced by the refusal token, the continuation is typed with the terminal
fact added to $\Gamma$, and the fact is true in $C'$, so (C4) holds with
$\Gamma'=\Gamma,\mathsf{Terminal}\langle c,\mathsf{Refused}\rangle$.
\item \emph{Collapse.} As for refusal, with the observation token and the
collapsed fact; $o=\mathsf{observe}_k(v)$ is defined since $k$ is total
on $B$.
\item \emph{Discharge and timeout.} The awaiting process holds the right
for $c$ (\textsc{T-Await}, \textsc{T-Deadline}), so $c$ is live by (C3),
and the continuation taken, $P$ or $Q$, keeps the right. The corresponding
clause of $\delta$ leaves the state unchanged, so (C2) and (C3) hold with
$\Delta'=\Delta$.
\end{itemize}

\emph{Contextual rules.} For $P\mid Q$, \textsc{T-Par} splits
$\Delta=\Delta_1\uplus\Delta_2$. The step is taken by $P$, and by the
induction hypothesis $P'$ is typed under $\Delta_1'$. The prefix that
fired either introduced a name fresh for $\Sigma$, hence not mentioned in
$\Delta_2$, or consumed a right in $\Delta_1$; Lemma~\ref{lem:recomb}
gives $\Delta_1'\uplus\Delta_2=\Delta'$, and \textsc{T-Par} retypes
$P'\mid Q$ after weakening $Q$'s typing with the new facts and names
(Lemma~\ref{lem:weak}). For $P+Q$, the premise of \textsc{T-Sum} types $P$
under the same $\Delta$, and the induction hypothesis applies. For
restriction, the induction hypothesis applies under $\Sigma,c$. For the
congruence rule, Lemma~\ref{lem:congr}.

The conditions on the other components of the configuration are carried
over unchanged or updated as described in each case.
\end{proof}

\begin{corollary}[Rights are neither duplicated nor lost]
In every configuration reachable from a well-typed one, each live
commitment has exactly one right in the linear context and each terminal
commitment has none.
\end{corollary}

\begin{proof}
(C3) is preserved by Theorem~\ref{thm:sr}.
\end{proof}

\begin{corollary}[Typed path lemma]
In a well-typed configuration, if $\Commit\langle c,A,s\rangle\in\Delta$
then $H|_c$ is a $\delta$-path from an opening record to a state of kind
$s$. It is unique in the sense that every linearization of the history's
dependence order has the same projection onto $c$.
\end{corollary}

\begin{proof}
The first sentence is Lemma~\ref{lem:canon}. For the second, records of
the same subject are pairwise dependent, so every linearization keeps
them in the same relative order.
\end{proof}

\section{Progress}

Progress in an ordinary calculus says a well-typed term is a value or can
step. Here a third possibility is legitimate: waiting for the environment.
The danger is that ``waiting'' becomes a place for errors to hide. It
cannot here, because waiting is a positive syntactic form, the
$\mathsf{await}$ prefix, which names the obligation and the principal
expected to discharge it.

A process is \emph{inert} if it is structurally congruent to a parallel
composition of $\mathbf 0$ and replications of processes with no prefix.
A configuration is \emph{waiting on} $\omega$ if every top-level prefix of
its process is an $\mathsf{await}$ prefix, with or without a deadline, and
one of them awaits $\omega$.
The \emph{top-level prefixes} of $P$ are those not beneath another
prefix; a summand's prefixes count, and a replicated body's prefixes
count.

\begin{hypothesis}[Computational totality at $C$]\label{hyp:total}
For every top-level prefix $\mathsf{transform}\;c\;\mathsf{with}\;f$ of
$C$ with $A(c)=\mathsf{bound}(q,b)$, evaluation of $f\,b$ terminates; and
for every top-level prefix $\mathsf{verify}_a\;c\;\mathsf{using}\;g$ with
$A(c)=\mathsf{transformed}(b,v)$, the verifier $g(c,v)$ returns
$\mathsf{ok}$ or $\mathsf{fail}$.
\end{hypothesis}

The hypothesis is about the particular computations a configuration is
about to perform, not about the lambda layer as a whole, which may be
Turing complete. For the fragment without fixed points it holds
automatically by strong normalization of the simply typed calculus
\cite{tait}.

\begin{theorem}[Progress]\label{thm:progress}
Let $\Sigma;\Gamma;\Delta\vdash C\;\mathsf{ok}$ with
Hypothesis~\ref{hyp:total} holding at $C$. Then exactly one of the
following holds:
\begin{enumerate}[label=(\roman*)]
\item $C$ can take a transition other than a discharge;
\item $C$ is waiting on some obligation;
\item the process of $C$ is inert, and every commitment in $\dom(A)$ is
terminal.
\end{enumerate}
\end{theorem}

\begin{proof}
If some top-level prefix is not an $\mathsf{await}$, we show it is
enabled, so (i) holds. By Lemma~\ref{lem:inv} its typing premises hold,
and by (C3) and (C4) the runtime agrees with them.
\begin{itemize}[leftmargin=1.5em]
\item A pop is always enabled: names are unbounded, so some $c$ lies
outside $\dom(A)$.
\item A compensating pop has $\mathsf{Terminal}\langle d,t\rangle$ in its
$\Gamma$, true by (C4), and its authority fact; a fresh name exists.
\item A bind, transform, verify, refuse, or collapse prefix on $c$ holds
the right for the state it requires (inversion), and by (C3) $A(c)$ is in
that state. Authority is a premise of the typing. A transform or verify
terminates by Hypothesis~\ref{hyp:total}. A collapse finds a valid witness
in the state by Lemma~\ref{lem:canon} and a defined observation since $k$
is total on the value's type.
\item A prefix under a sum, restriction, or replication fires through the
corresponding contextual rule; for replication, through
$!R\equiv R\mid\,!R$.
\end{itemize}
If every top-level prefix is an $\mathsf{await}$ and there is at least one,
(ii) holds. If there is none, the process is inert; its typing derivation
ends in instances of \textsc{T-Nil}, so $\Delta$ contains only tokens, and
by (C3) no commitment is live.

The three cases are exclusive by definition of the second and third, and
since a configuration with an enabled non-await prefix is neither waiting
nor inert.
\end{proof}

The third case is the formal statement that a well-typed program cannot
silently abandon a commitment: if nothing can happen and nothing is being
waited for, every commitment has been disposed of. A process attempting to
collapse an unverified commitment is not waiting and not inert; it is
untypable.

\section{Protocol fidelity}
\label{sec:fidelity}

Read the refusal completion table as a protocol: for each live state $s$,
the protocol offers the actions in $\mathsf{Next}(s)$, and each action
leads to the protocol of its successor state. For a single commitment this
is the session type
\[
\begin{aligned}
\mathcal S=\mathsf{open};\bigl(&\mathsf{refuse}\oplus\mathsf{bind};
(\mathsf{refuse}\oplus\mathsf{transform};(\mathsf{refuse}\\
&\oplus\mathsf{verify}^{+};(\mathsf{refuse}\oplus\mathsf{collapse})
\oplus\mathsf{verify}^{-};\mathsf{refuse}))\bigr),
\end{aligned}
\]
where $\mathsf{verify}^{+}$ and $\mathsf{verify}^{-}$ are the two outcomes
of verification, and after a negative outcome the only continuation is a
refusal.

\begin{theorem}[Protocol fidelity]\label{thm:fidelity}
Along every execution from a well-typed configuration, the sequence of
actions whose subject is $c$ is a trace of $\mathcal S$.
\end{theorem}

\begin{proof}
By Theorem~\ref{thm:sr}, each step updates the linear context by the
context step, which follows exactly one edge of the table for the subject
of the action. The first action on $c$ is its pop or compensating pop.
\end{proof}

This is fidelity for the protocol of a single commitment. Protocols that
coordinate several principals over channels, with a protocol and its dual
held by different parties, need a calculus with communication; the
correspondence with the session-type literature, including Wadler's
propositions-as-sessions reading, is taken up when that extension is made.

\section{Negative derivations}

The type system's value is measured by what it rejects.

\begin{proposition}[Rejected programs]
There is no typing derivation for any of the following, in any context in
which the named commitments' rights are as described.
\begin{enumerate}[label=(\alph*)]
\item A collapse of a commitment whose right is open, bound, or
transformed.
\item A bind of a commitment that has been refused: only a token and a
fact remain, and \textsc{T-Bind} requires a right.
\item Two uses of one right: either in parallel, since
$\Delta_1\uplus\Delta_2$ cannot contain two rights for $c$, or in sequence,
since after the first use the right has moved to a different state.
\item Verifying $c_1$ with a verifier indexed by $c_2\neq c_1$:
\textsc{T-Ver} requires $\mathsf{Verifier}\langle c_1,B\rangle$.
\item Any attributed prefix without the corresponding authority fact.
\item A process that ends holding a live right: \textsc{T-Nil} admits only
tokens.
\end{enumerate}
\end{proposition}

\begin{proof}
Each by Lemma~\ref{lem:inv}: the only rule for the prefix in question has
a premise the described context cannot satisfy.
\end{proof}

Case (c) is the static half of terminal exclusivity. Operationally, two
components racing to refuse and collapse the same commitment produce
exactly one terminal record (Chapter~\ref{ch:lts}); statically, a
well-typed program cannot contain the race, because only one component can
hold the right.

\chapter{Authority as a Modality}
\label{ch:authority}

The core calculus treats authority as a fixed policy: a set of facts
$\mathsf{may}(a,\kappa,x)$ that no rule changes. That is enough for the
safety theorems, and it makes non-amplification trivial. Real systems
transfer authority, and a transfer is exactly the kind of event that must
be recorded, authenticated, and replayable. This chapter adds delegation
as a transfer of an exclusive capability, proves that it conserves
authority rather than copying it, and then applies the resulting
distinctions to a contemporary case.

\section{Exclusive authority as a linear capability}

Two kinds of authority are distinguished. \emph{Shared} authority is a
policy fact $\mathsf{may}(a,\kappa,x)$ in $\Gamma$, as before; any number
of principals may hold it and nothing transfers it. \emph{Exclusive}
authority is a linear capability
\[
\mathsf{Cap}\langle a,c,\kappa\rangle
\]
held in $\Delta$: the right of principal $a$, and of no one else, to
perform acts of kind $\kappa$ on commitment $c$. At runtime the
configuration carries an ownership map $\mathsf{Own}$ from pairs
$(c,\kappa)$ to the single principal who currently holds the exclusive
capability. An attributed rule whose authority premise was
$\mathsf{may}(a,\kappa,c)$ now accepts either the shared fact or current
ownership, $\mathsf{Own}(c,\kappa)=a$. Exercising a capability does not
consume it; the typing rules pass it on to the continuation unchanged, so
linearity forbids duplication and discarding without forbidding use.

Exclusive capabilities enter the world only through a grant record
$\mathsf{Grant}(g,a,c,\kappa,\sigma)$ issued by a root authority $g$, and
move only through delegation.

\section{Delegation}

The prefix $\mathsf{delegate}_{a\to b}\langle c,\kappa\rangle.P$
transfers $a$'s exclusive capability over $(c,\kappa)$ to $b$. Its typing
rule consumes the delegator's capability and gives the continuation the
delegate's:
\[
\frac{\Sigma;\Gamma;\Delta,\mathsf{Cap}\langle b,c,\kappa\rangle\vdash P}
{\Sigma;\Gamma;\Delta,\mathsf{Cap}\langle a,c,\kappa\rangle\vdash
\mathsf{delegate}_{a\to b}\langle c,\kappa\rangle.P}.
\]
Its operational rule checks current ownership and an authentication by
the delegator, moves ownership, and appends an authenticated record:
{\small
\[
\frac{\mathsf{Own}(c,\kappa)=a\qquad
\mathsf{checkSig}(\sigma,a,\langle b,c,\kappa\rangle)}
{\begin{array}{c}\langle\mathsf{delegate}_{a\to b}\langle c,\kappa\rangle.P,\ldots,
\mathsf{Own},H\rangle
\xrightarrow{\mathsf{delegate}(a,b,c,\kappa,\sigma)}\\
\langle P,\ldots,\mathsf{Own}[(c,\kappa)\mapsto b],
H\cdot\mathsf{Delegate}(a,b,c,\kappa,\sigma)\rangle\end{array}}.
\]
}
Replay reconstructs $\mathsf{Own}$ from grant and delegation records and
accepts $\mathsf{Delegate}(a,b,c,\kappa,\sigma)$ only if, at that point in
the ledger, $a$ owns $(c,\kappa)$, the signature authenticates the record,
and the commitment and kind match the signed payload. Because ownership
moves, a delegation cannot be replayed a second time by its former owner:
after it, $a$ no longer owns what it would be delegating. Scope and
expiry constraints, where a policy uses them, are further conjuncts of
the same check.

\section{Results}

Write $\mathsf{Owners}_H(c,\kappa)$ for the set of principals holding the
exclusive capability for $(c,\kappa)$ after replay of $H$.

\begin{proposition}[Delegation authenticity]\label{prop:delauth}
If replay of $H$ is defined and $H$ contains
$\mathsf{Delegate}(a,b,c,\kappa,\sigma)$, then at that point in $H$ the
delegator $a$ owned $(c,\kappa)$ and $\sigma$ authenticates $a$'s transfer
to $b$.
\end{proposition}

\begin{proof}
These are the conditions under which replay accepts the record.
\end{proof}

\begin{proposition}[Conservation]\label{prop:conserve}
For every $(c,\kappa)$ that has been granted, $|\mathsf{Owners}_H(c,\kappa)|=1$
after every ledger prefix, and a delegation from $a$ to $b$ changes it by
\[
\mathsf{Owners}_{H'}(c,\kappa)=
\bigl(\mathsf{Owners}_{H}(c,\kappa)\setminus\{a\}\bigr)\cup\{b\}.
\]
In particular delegation never increases the number of holders of an
exclusive capability.
\end{proposition}

\begin{proof}
A grant sets the owner of an ungranted pair. A delegation is accepted
only from the current owner and replaces that owner by the delegate. No
other record changes $\mathsf{Own}$.
\end{proof}

\begin{proposition}[Preservation under transfer]\label{prop:delpres}
Extend well-typed configurations with the condition
$\mathsf{CapsAgree}(\mathsf{Own},\Delta)$: $\mathsf{Cap}\langle
a,c,\kappa\rangle\in\Delta$ exactly when $\mathsf{Own}(c,\kappa)=a$. Then
Theorem~\ref{thm:sr} holds for the calculus with delegation.
\end{proposition}

\begin{proof}
The new case is delegation. Inversion of its typing rule gives the
delegator's capability in $\Delta$ and types the continuation with the
delegate's; the operational rule moves ownership from $a$ to $b$, so
$\mathsf{CapsAgree}$ is maintained, and the rights, commitment map, and
history conditions are untouched except for the appended record, which
replay accepts. Every other case is as before, with the authority premise
now discharged either by a shared fact or, by $\mathsf{CapsAgree}$, by a
capability in the context. Context splitting keeps each capability in
exactly one component, by the same ownership argument as for rights
(Lemma~\ref{lem:own}).
\end{proof}

\begin{proposition}[Chain soundness]\label{prop:chain}
If an attributed record in $H$ was authorized by an exclusive capability,
there is a chain
\[
\mathsf{Grant}(g,a_0,c,\kappa,\sigma_0)\;;\;
\mathsf{Delegate}(a_0,a_1,c,\kappa,\sigma_1)\;;\;\cdots\;;\;
\mathsf{Delegate}(a_{n-1},a_n,c,\kappa,\sigma_n)
\]
of authenticated records earlier in $H$, with $a_n$ the principal of the
record.
\end{proposition}

\begin{proof}
By induction on the ledger, ownership of $(c,\kappa)$ at any point is the
end of the chain of records that set it, each accepted only from the
previous owner (Propositions~\ref{prop:delauth} and~\ref{prop:conserve}).
\end{proof}

Together these give non-amplification in the presence of delegation:
after any execution, the principals holding authority of kind $\kappa$
over $c$ are those holding the shared fact, which no rule adds, together
with the single current owner of the exclusive capability, whose
authority is traced to a grant by an authenticated chain.

\paragraph{Confinement.} The results depend on capabilities having no
other source. If the lambda layer could produce a value of capability
type, for instance a function returning
$\mathsf{Cap}\langle b,c,\kappa\rangle$ from nothing, conservation would
fail at once: a transformation could mint authority. The calculus
therefore gives capabilities no expression form. They appear only in
linear contexts, introduced by grants and moved by delegation, and no
value typing rule mentions them. That is a design constraint the
propositions rely on, and any extension that lets computation handle
capabilities as data must re-establish it.

\paragraph{What is left out.} Revocation is not in the core. A
revocation can race with the exercise of the capability it revokes, and
deciding which happened first requires a declared linearization point,
the same obligation that operational terminal exclusivity places on the
sequencer (Hypothesis~\ref{hyp:atomic}). It is a later extension.

\section[Case study: evidence is not authority]{Case study: evidence production is not disposition authority}

A recent study of proposal selection at a large neutron-scattering
facility gives a clean contemporary case for the distinctions this chapter
formalizes \cite{ding}. The authors used a language model to compare every
pair of proposals in a review cycle, aggregated the pairwise preferences
into a ranking with a Bradley--Terry model, and compared the result with
the historical human rankings. Across cycles the rank correlations ranged
from roughly $0.2$ to $0.8$, typically around $0.6$, and proposals on
which the two rankings diverged sharply were flagged for the review
committee. They also used embedding similarity between proposals as a
cheap triage signal for possible duplicates, and they are explicit that
this signal is not a validated redundancy classifier: a high-similarity
pair may be a resubmission, related work, or merely shared vocabulary, and
telling these apart needs expert judgement. Their stated position is that
the model should augment the committee rather than replace it, and that
the committee keeps final authority over flagged divergences.

Two features of the evaluation bear on authority as well as on accuracy.
Publication outcomes, used to compare the rankings' effectiveness, exist
only for accepted proposals, and acceptance was historically decided by
the human ranking; the authors note that this biases the comparison
toward the human ranking. And where either ranking had ties, the other
method's score was used as a secondary sorting key, so the two rankings
are not wholly independent in tied cases. Neither point undoes the
feasibility result. Both illustrate that evidence about an evidence
producer's reliability can itself be shaped by earlier dispositions, which
is exactly the kind of dependence a ledger makes auditable.

The study involves four roles, and the calculus can keep them apart:
\begin{itemize}[leftmargin=1.5em]
\item a \emph{signal producer}, such as the embedding-similarity score,
which flags candidates for attention and has no evidential standing;
\item an \emph{evidence producer}, such as the model's pairwise
comparisons aggregated by Bradley--Terry, which produces a ranking with
provenance that others may consult;
\item a \emph{verifier}, which checks a claim against an independent truth
condition;
\item a \emph{disposition authority}, which accepts or refuses.
\end{itemize}
In the study the model is an evidence producer. Pairwise preference is not
verification of scientific merit against an independent condition, and
nothing in the study treats it as such. A committee may choose to verify
the ranking \emph{procedure}, for instance by checking position-swap
consistency, but that is a verification of the procedure, performed under
the committee's authority, not the model verifying the proposals.

Each role corresponds to a class of records it is able to write:
\begin{center}
\small
\begin{tabular}{lll}
\toprule
Role & Authority kind & Records it can write\\
\midrule
signal producer & (annotation) & annotation-channel entries only; nothing causal\\
evidence producer & $\mathsf{discharge}$ & $\mathsf{Discharged}$ records; lifecycle-neutral\\
verifier & $\mathsf{verify}$ & $\Verify$ and $\mathsf{VerificationFailed}$ records; never terminal\\
disposition authority & $\mathsf{collapse}$, $\mathsf{refuse}$ & $\Collapse$ and explicit $\Refuse$ records\\
\bottomrule
\end{tabular}
\end{center}

\begin{proposition}[Non-derivability of authority]\label{prop:nonderiv}
In the calculus with a static policy, no transition and no typing rule
adds an authority fact, and every attributed record requires the
authority kind $\mathsf{auth}(\alpha)$ of Chapter~\ref{ch:lts}. Hence a
principal whose facts are exactly those of one role writes, in every
execution, only records in that role's row of the table. In particular
\[
\mathsf{signal}\not\Rightarrow\mathsf{evidence},\qquad
\mathsf{evidence}\not\Rightarrow\mathsf{verify},\qquad
\mathsf{verify}\not\Rightarrow\mathsf{refuse},\qquad
\mathsf{verify}\not\Rightarrow\mathsf{collapse}.
\]
\end{proposition}

\begin{proof}
The policy is a fixed set of facts consulted only by membership, and no
rule's conclusion adds to it. Theorem~\ref{thm:auth} gives the required
kind for each attributed record, and the table lists, for each kind, the
records whose required kind it is.
\end{proof}

These are non-derivability results from the rules, not empirical claims
that a model can never hold a later role. A policy may explicitly grant
verification authority to a model. What the proposition rules out is
authority arising merely because a component produced a score or a
ranking. With delegation the same holds in stronger form: a principal can
come to hold exclusive authority only at the end of an authenticated chain
of grants and delegations (Proposition~\ref{prop:chain}), each of which is
an explicit, recorded act by the previous holder.

The fourth relation depends on a design decision taken for exactly this
reason. In an earlier version of the calculus a failed verification
refused its commitment automatically, which gave every verifier a refusal
power through the negative branch even though the type system denied it
explicit refusal authority. Verification now leads to
$\mathsf{verificationFailed}$, a live state, and only a principal holding
refusal authority can end the commitment from there. Verification
classifies evidence; disposition determines continuation. The cost is one
more live state, and one more way for a commitment to wait: a failed
commitment whose disposition authority never acts stays live, and
Chapter~\ref{ch:progress} lists that as a counterexecution rather than
hiding it.

\paragraph{The encoding.} Give the model only discharge authority. Each
proposal awaits an obligation $\omega_{\mathrm{rank}}$, the model's
ranking, which the model discharges with an authenticated certificate;
the committee holds verification, collapse, and refusal:
\[
\begin{aligned}
&\mathsf{pop}\;q\;\mathsf{as}\;c.\;
\mathsf{bind}_s\;c\;\mathsf{to}\;\mathit{proposal}.\;
\mathsf{await}_m\;c\;\omega_{\mathrm{rank}}.\;
\mathsf{transform}\;c\;\mathsf{with}\;\mathit{summary}.\\
&\qquad\mathsf{verify}_h\;c\;\mathsf{using}\;\mathit{review}\;
\mathsf{then}\;\mathsf{collapse}_{h,\mathrm{alloc}}\;c.\mathbf 0\;
\mathsf{else}\;\mathbf 0,
\end{aligned}
\]
with policy $\mathsf{may}(s,\mathsf{bind},c)$ for the intake principal
$s$, $\mathsf{may}(m,\mathsf{discharge},\omega_{\mathrm{rank}})$, and
$\mathsf{may}(h,\kappa,c)$ for $\kappa\in\{\mathsf{verify},
\mathsf{collapse},\mathsf{refuse}\}$. A flagged divergence becomes a second
obligation on which the committee's own verification waits.

\begin{corollary}[Signals and evidence do not dispose]\label{prop:signals}
In every execution of this encoding, every record attributed to $m$ is a
discharge record, and no record attributed to $m$ changes any
commitment's lifecycle state.
\end{corollary}

\begin{proof}
Proposition~\ref{prop:nonderiv} for the evidence-producer row, with the
discharge clause of $\delta$ returning the incoming state.
\end{proof}

\paragraph{Provenance of the ranking.} A rank is a number, and a number
looks more determinate than the process that produced it. The discharge
certificate for $\omega_{\mathrm{rank}}$ should therefore name, in the
claim it hashes, everything the rank depends on: the comparison set, the
model and prompt versions, the ordering controls used against positional
bias, the aggregation parameters, and whether a human score was used as a
tie-breaker. The last item matters for this study in particular, since its
tie-breaking rule made the two rankings partly dependent. With those
fields in the certificate, a later reader of the ledger can tell which
ranking the committee saw and how it was made; without them, the ledger
records that a ranking was delivered but not what it was evidence of.

In a typed program the separation is static as well: any prefix in which
$m$ binds, verifies, collapses, or refuses fails to type, since each rule
requires the corresponding authority fact. The committee's action supplies
the attributable disposition. If the model's ranking were treated as the
decision itself, synthetic evidence would have become synthetic
authority, and the ledger would show exactly where: a terminal record
attributed to a principal that should not have held the authority to
write it. The general lesson is that when a system grants an automated
component authority, the question is not what the component is meant to
do but which records it is able to write.

%===================================================================
\part{The Disposition Calculus}
%===================================================================

\chapter{The Ledger Path Lemma}
\label{ch:path}

The theorems of this Part all concern what a history must contain once
the commitment map says something about a commitment. They share one
proof. Rather than repeat the same induction five times, this chapter
proves a single invariant relating the ledger to the map, and the next
three chapters read the individual results off it.

The projection $H|_c$, the subject of a record, and the lifecycle step
function $\delta$ were defined in Section~\ref{sec:records}. The whole
chapter turns on one feature of $\delta$ worth restating: it is a partial
fold, not a filter. Every clause either advances the lifecycle or, for a
discharge, checks authenticated evidence against a live state; none
passes a record through unexamined. So if $\delta^{*}(H|_c)$ is defined,
every record of $H|_c$ was consumed by exactly one step of the fold. A
history cannot fold correctly while carrying a record the fold ignored.

\section{The lemma}

Call an initial configuration $C_0=\langle P_0,S_0,A_0,O_0,H_0\rangle$
\emph{ledger consistent} when $\delta^{*}(H_0|_c)=A_0(c)$ for every
$c\in\dom(A_0)$ and $H_0|_c$ is empty for every $c\notin\dom(A_0)$. The
empty configuration, with $A_0$, $O_0$, and $H_0$ all empty, is ledger
consistent.

\begin{lemma}[Ledger path lemma]\label{lem:path}
If $C_0$ is ledger consistent and $C_0\to^{*}C=\langle P,S,A,O,H\rangle$,
then for every commitment identity $c$:
\begin{enumerate}[label=(\roman*)]
\item $c\in\dom(A)$ if and only if $H|_c$ is nonempty;
\item if $c\in\dom(A)$ then $\delta^{*}(H|_c)$ is defined and equals
$A(c)$.
\end{enumerate}
\end{lemma}

\begin{proof}
By induction on the length of the execution. The base case is ledger
consistency. For the step $C\xrightarrow{\alpha}C'$, history monotonicity
gives $H'=H\cdot\rho$ with $\rho=\record(\alpha)$, and let $s=\subj(\rho)$.

For $c\neq s$, we have
$H'|_c=H|_c$ because $\rho$ is not in the projection, and $A'(c)=A(c)$ with $c\in\dom(A')\iff c\in\dom(A)$ by
Theorem~\ref{thm:identity}. Both clauses carry over.

For $c=s$, we have $H'|_s=H|_s\cdot\rho$ and must show
$\delta(\delta^{*}(H|_s),\rho)=A'(s)$, where the inner fold is taken to be
$\bot$ when $H|_s$ is empty. We examine the primitive rule at the root of
the derivation.
\begin{itemize}[leftmargin=1.5em]
\item \emph{Pop or compensating pop.} The premise $s\notin\dom(A)$ and
clause~(i) give $H|_s$ empty, so the fold is $\bot$, and $\delta$ returns
$\mathsf{open}(q)=A'(s)$. Clause~(i) now holds since $H'|_s=\langle\rho\rangle$.
\item \emph{Bind.} The premise $A(s)=\mathsf{open}(q)$ and clause~(ii)
give $\delta^{*}(H|_s)=\mathsf{open}(q)$, and $\delta$ applied to
$\Bind(a,s,b)$ returns $\mathsf{bound}(q,b)=A'(s)$.
\item \emph{Transform.} From $A(s)=\mathsf{bound}(q,b)$ the fold yields
that state, the record carries the same $b$ because the rule reads it
from $A(s)$, and $\delta$ returns $\mathsf{transformed}(b,v)$.
\item \emph{Verification, success.} From
$A(s)=\mathsf{transformed}(b,v)$, the record carries that $v$, the premise
$\valid(w,s,v)$ satisfies the side condition, and $\delta$ returns
$\mathsf{verified}(v,w)$.
\item \emph{Verification, failure, and explicit refusal.} Both premises
place $s$ in a live state, satisfying the side condition, and $\delta$
returns $\mathsf{refused}(r)$ with the recorded $r$.
\item \emph{Discharge and timeout.} The premises give a live state and an
authentic certificate, which are the side conditions of the corresponding
clause, and $\delta$ returns the same state, as the rule leaves $A(s)$
unchanged.
\item \emph{Collapse.} From $A(s)=\mathsf{verified}(v,w)$, the record
carries that $w$, the premise $o=\mathsf{observe}_k(v)$ satisfies the
side condition, and $\delta$ returns $\mathsf{collapsed}(o)$.
\end{itemize}
In each case the rule's premises are precisely the conditions under which
$\delta$ is defined on the new record, and its conclusion is precisely
$\delta$'s value.
\end{proof}

The proof has a shape worth noticing. Nothing in it is specific to
refusal, binding, or collapse; it only checks that each operational rule
and the corresponding clause of $\delta$ say the same thing. That is the
sense in which the rest of this Part is ``nearly forced by the preceding
architecture''. The substance was spent in designing the rules so that
each premise is a check on the ledger and each conclusion is a record.

\section{Causal order}

The ancestry theorems speak of one record causally preceding another,
$e\prec_H e'$, and the spine asks that ledger order count only ``where
order represents a true dependency''. The path lemma supplies the
dependency for records of the same commitment: each record in $H|_c$ could
be appended only because the previous one had produced the state its rule
required. For records of different commitments, Chapter~\ref{ch:replay}
shows that ledger order is generally not a dependency at all, since
adjacent records of distinct subjects may be exchanged without changing
anything replay can recover. We therefore take $\prec_H$ to be the
dependence order of Definition~\ref{def:dependence}, and note that every
ancestry chain proved in this Part lies within a single projection $H|_c$,
where ledger order and $\prec_H$ agree.

\chapter{Refusal, Binding, and Evidence}
\label{ch:rbe}

Throughout this chapter and the next, $C_0$ is ledger consistent,
$C_0\to^{*}C=\langle P,S,A,O,H\rangle$, and Hypothesis~\ref{hyp:stable}
holds. Write $\mathsf{Open}(c,q)$ for either opening record, $\Pop(c,q)$
or $\Compensates(a,c,q,d)$.

\section{Refusal preservation}

Refusal must be treated as a successful terminal disposition, not as a
missing execution.

\begin{theorem}[Refusal preservation]\label{thm:refusal}
If $A(c)=\mathsf{refused}(r)$ at some configuration $C_i$ of an execution,
then at every later configuration $C_j$,
\[
A_j(c)=\mathsf{refused}(r)
\qquad\text{and}\qquad
\Refuse(a,c,r)\in H_j
\]
for the principal $a$ that refused it.
\end{theorem}

\begin{proof}
The first equation is lifecycle monotonicity, since
$\mathsf{refused}(r)$ is $\preceq_D$-maximal and is related only to
itself, together with the observation that no rule rewrites the reason of a
terminal state. For the second, the path lemma at $C_i$ gives
$\delta^{*}(H_i|_c)=\mathsf{refused}(r)$; the only clause of $\delta$
producing that value is the refusal clause, so the last record of $H_i|_c$
is $\Refuse(a,c,r)$ for some $a$; and history monotonicity carries it into
every $H_j$.
\end{proof}

The spine stated refusal preservation with an exception for a
compensating successor. In this calculus the exception is not needed:
compensation introduces a new identity and leaves the predecessor's state
untouched (Theorem~\ref{thm:identity}), so the refused commitment remains
refused forever and it is the \emph{successor} that carries the matter
forward. A later acceptance cannot erase the refusal because there is no
rule that could even address it.

\section{Binding integrity}

A binding should establish a durable relation without confusing
association with identity.

\begin{corollary}[Binding integrity]
If $A(c)=\mathsf{bound}(q,b)$, then
\[
H|_c=\langle \mathsf{Open}(c,q),\;\Bind(a,c,b)\rangle
\]
for some principal $a$. In particular the binding record is unique, it is
preceded by the opening record of the same commitment, and it binds the
value $b$ held in the state.
\end{corollary}

\begin{proof}
By the path lemma, $\delta^{*}(H|_c)=\mathsf{bound}(q,b)$. Inspecting
$\delta$, the only record sequences folding from $\bot$ to a bound state
are an opening record followed by one binding record.
\end{proof}

Commitment identities and expression values are disjoint sorts
(Chapter~\ref{ch:syntax}), so $c$ and $b$ cannot become interchangeable
because their concrete encodings coincide. The stronger clause of the
spine, that an exclusive capability is bound to at most one live
commitment, is not provable in the untyped calculus, where nothing marks a
value as exclusive. It is a consequence of linearity and is proved in
Part~IV.

\section{Witness ancestry}

\begin{corollary}[Witness ancestry]\label{cor:witness}
If $A(c)=\mathsf{verified}(v,w)$, then
\[
H|_c=\langle \mathsf{Open}(c,q),\;\Bind(a,c,b),\;
\Transform(c,f,b,v),\;\Verify(a',c,v,w)\rangle
\]
for some $q,a,b,f,a'$, and $\valid(w,c,v)$ holds.
\end{corollary}

\begin{proof}
The only record sequences folding from $\bot$ to $\mathsf{verified}(v,w)$
are an opening, a binding, a transformation, and a verification, with the
matching conditions of $\delta$ forcing the $b$ of the binding to equal
the input of the transformation and the output $v$ of the transformation
to equal the value verified. The side condition of the verification clause
is $\valid(w,c,v)$, and by stability it still holds.
\end{proof}

The four records form a chain in $\prec_H$ because they lie in one
projection. The spine's requirement that a witness ``witnesses'' the
transformation is realized by indexing: $w$ is valid for the pair $(c,v)$,
and $v$ is the output recorded in $\Transform(c,f,b,v)$. A witness produced
for another commitment, or for another value of this one, fails the
matching conditions of $\delta$ and has no rule by which it could enter
this chain. What indexing does \emph{not} rule out is a verifier $g$ that
issues valid witnesses too liberally; the calculus guarantees provenance
of evidence, not the quality of the predicate, and that limit is stated
here so that no later chapter is read as claiming otherwise.

\chapter{Collapse Soundness and Terminal Exclusivity}
\label{ch:collapse}

This is the principal local safety theorem of the untyped calculus.

\begin{theorem}[Collapse soundness]\label{thm:collapse}
If $A(c)=\mathsf{collapsed}(o)$, then
\[
H|_c=\langle \mathsf{Open}(c,q),\;\Bind(a,c,b),\;
\Transform(c,f,b,v),\;\Verify(a',c,v,w),\;
\Collapse(a'',c,k,w,o)\rangle,
\]
with $\valid(w,c,v)$ and $o=\mathsf{observe}_k(v)$, and $(c,k,o)\in O$.
\end{theorem}

\begin{proof}
By the path lemma and inspection of $\delta$, the only sequences folding
to $\mathsf{collapsed}(o)$ are a sequence folding to some
$\mathsf{verified}(v,w)$ followed by a collapse record carrying the same
$w$ and satisfying $o=\mathsf{observe}_k(v)$. Corollary~\ref{cor:witness}
applied to that prefix supplies the first four records and validity. The
observation $(c,k,o)$ was added to $O$ by the collapse rule that appended
the last record, and no rule removes observations.
\end{proof}

The theorem rules out each of the failures the spine names. There is no
collapse \emph{ex nihilo}, because the fold requires an opening record.
There is no collapse from an unrelated witness, because $\delta$ requires
the collapse record's witness to be the one held in the verified state,
which Corollary~\ref{cor:witness} ties to this commitment and value. There
is no collapse after refusal, because $\delta$ is undefined on any record
following a refusal. And there is no collapse whose ancestry has been
omitted, because the projection is a subsequence of an append-only
history and the fold must pass through every stage.

\begin{corollary}[Terminal exclusivity, history form]
\label{cor:terminal-history}
For every commitment $c$, the projection $H|_c$ contains at most one
terminal record, and if it contains one, that record is last. In
particular
\[
\Refuse(a,c,r)\in H
\quad\Longrightarrow\quad
\Collapse(a',c,k,w,o)\notin H.
\]
\end{corollary}

\begin{proof}
By the path lemma $\delta^{*}(H|_c)$ is defined. Since $\delta$ is
undefined on every record whose incoming state is terminal, no record
follows a terminal record in a sequence on which $\delta^{*}$ is defined.
\end{proof}

Reconsideration therefore never produces two terminal records for one
identity. It produces a chain of distinct identities, which is the subject
of the next chapter.

\begin{remark}[Which collapse rule]
The spine stated collapse soundness with a single condition
$\operatorname{denotes}(v,o)$. Here the condition is $o=\mathsf{observe}_k(v)$
for the declared rule $k$, recorded in the ledger. This is where the
disposition calculus meets the kernel's rule-indexed collapse, and it
bears directly on the third comparative claim: two commitments with equal
verified values may collapse under different rules to different
observations, and the ledger says which rule produced which. Whether this
is more than a relabelling of a family of observational equivalences is
the question of Chapter~\ref{ch:spectrum}.
\end{remark}

\chapter{Compensation}
\label{ch:compensation}

Compensation must add history rather than revise it. It is the formal
boundary between historical reversibility, which the calculus forbids, and
material repair, which it permits.

\begin{theorem}[Compensating continuation]\label{thm:compensation}
If $C\xrightarrow{\mathsf{compensate}(a,c',q,c)}C'$, then
\begin{enumerate}[label=(\roman*)]
\item $c'\neq c$;
\item $H'|_c=H|_c$ and $A'(c)=A(c)$, so the terminal record of $c$ is
present in $H'$ and unchanged;
\item $\Compensates(a,c',q,c)\in H'$ and $H\sqsubseteq H'$;
\item $\mathsf{may}(a,\mathsf{compensate},c)$ held in $C$.
\end{enumerate}
\end{theorem}

\begin{proof}
The rule's premises give $c'\notin\dom(A)$ and $\mathsf{terminal}(A(c))$,
so $c\in\dom(A)$ and (i) follows. Part (ii) is Theorem~\ref{thm:identity}
together with the fact that the new record's subject is $c'$; by
Corollary~\ref{cor:terminal-history} and the path lemma the last record of
$H|_c$ is terminal. Part (iii) is history monotonicity and (iv) is
Theorem~\ref{thm:auth}.
\end{proof}

The spine required that ``no projection may entail that the original
disposition never occurred''. That clause can be given a precise form.

\begin{definition}[Supersession]
Write $c\leadsto c'$ when $\Compensates(a,c',q,c)\in H$. The
\emph{current successor} of $c$ is the last element of the longest chain
$c\leadsto c_1\leadsto\cdots\leadsto c_n$ in $H$, if such chains are
unique, and the set of chain endpoints otherwise.
\end{definition}

\begin{proposition}[Supersession is a view, not a revision]
\label{prop:supersession}
The relation $\leadsto$ is acyclic. For every commitment $c$ and every
extension $H\sqsubseteq H'$ produced by execution,
$\delta^{*}(H'|_c)=\delta^{*}(H|_c)$ whenever $c$ was terminal in $H$.
\end{proposition}

\begin{proof}
If $c\leadsto c'$ then the opening record of $c'$ appears later in $H$ than
the opening record of $c$, since $c$ was already in the domain when $c'$
was introduced. Opening positions strictly increase along any chain, so no
chain returns to its start. For the second statement, a terminal $c$
admits no further record in its projection by
Corollary~\ref{cor:terminal-history}, so $H'|_c=H|_c$.
\end{proof}

Ancestry that runs through compensation crosses projections: the record
\[
\Compensates(a,c',q,c)
\]
lies in $H|_{c'}$ while referring to $c$. The
path lemma covers each projection separately, so the crossing needs its
own statement. Write $c\triangleleft_H c'$ when
$\Compensates(a,c',q,c)\in H$.

\begin{lemma}[Compensation ancestry]\label{lem:companc}
Let $H$ be the history of an execution from a consistent configuration.
If $c\triangleleft_H c'$, then $H|_c$ ends in a terminal record, and that
record precedes the compensation record both in $H$ and in the dependence
order $\prec_H$. Consequently the ancestry of any commitment is a finite
alternating sequence
\[
\pi_{c_0}\;\triangleleft\;\pi_{c_1}\;\triangleleft\;\cdots\;\triangleleft\;\pi_{c_n},
\]
where each $\pi_{c_i}$ is the lifecycle path $H|_{c_i}$, each $\pi_{c_i}$
with $i<n$ ends in a terminal record, and every record is
$\prec_H$-after every record to its left.
\end{lemma}

\begin{proof}
At the configuration where the compensation record was appended, the
rule's premise made $A(c)$ terminal, so by the path lemma the last record
of $H|_c$ at that point was terminal, and it was already in the history.
The compensation record refers to $c$, so it is dependent on every record
of $H|_c$ (Definition~\ref{def:dependence}), and hence $\prec_H$-after
them. Records within one projection are $\prec_H$-ordered because they
share a subject. Finiteness and the absence of repeats follow from
acyclicity of $\triangleleft_H$ (Proposition~\ref{prop:supersession}).
\end{proof}

A projection that reports only current successors is legitimate; it is a
choice of what to display. What no execution can produce is a history from
which the predecessor's disposition has become unrecoverable, because that
disposition is a function of $H|_c$ and $H|_c$ never changes again. The
distinction between the two is exactly the distinction between a
$\Collapse$ rule and an edit.

Chains need not be unique. Two principals, each authorized, may
compensate the same refused commitment, producing two successors. The
calculus records both and does not choose; whether a policy should forbid
this, or a later collapse should arbitrate between successors, is listed
in Appendix~\ref{app:open}.

%===================================================================
\part{Replay, Progress, and Concurrency}
%===================================================================

\chapter{Replay}
\label{ch:replay}

\section{Replay as a fold}

Replay reconstructs the disposition and observation state from the
ledger. Its state is a pair $X=\langle A,O\rangle$ and it applies one
record at a time:
\[
\mathsf{apply}(\langle A,O\rangle,\rho)=
\begin{cases}
\langle A[s\mapsto\delta(A(s),\rho)],\;O\cup\{(s,k,o)\}\rangle
  & \text{if }\rho=\Collapse(a,s,k,w,o),\\
\langle A[s\mapsto\delta(A(s),\rho)],\;O\rangle & \text{otherwise},
\end{cases}
\]
where $s=\subj(\rho)$ and $A(s)$ is read as $\bot$ when
$s\notin\dom(A)$. For a compensation record $\Compensates(a,s,q,d)$,
$\mathsf{apply}$ is additionally defined only when $A(d)$ is terminal,
matching the operational premise. $\mathsf{apply}$ is undefined whenever
$\delta$ is, so replay of a ledger that no execution could have produced
fails rather than silently repairing it. Then
\[
\replay_{X_0}(H)=\text{the left fold of }\mathsf{apply}\text{ over }H
\text{ from }X_0.
\]

\begin{theorem}[Replay correspondence]\label{thm:replay}
Let $C_0$ be ledger consistent with
$\replay_{\langle\varnothing,\varnothing\rangle}(H_0)=\langle A_0,O_0\rangle$.
If $C_0\to^{*}\langle P,S,A,O,H\rangle$, then
\[
\replay_{\langle\varnothing,\varnothing\rangle}(H)=\langle A,O\rangle.
\]
\end{theorem}

\begin{proof}
By induction on the execution. For a step appending $\rho$ with subject
$s$, $\mathsf{apply}$ updates the map at $s$ by $\delta$, which by the
path lemma agrees with the operational update, and leaves every other key
unchanged, which agrees with Theorem~\ref{thm:identity}. For a
compensation record the additional condition on $A(d)$ is the operational
premise. The observation component changes only on a collapse record, by
the same triple the collapse rule adds to $O$.
\end{proof}

The store $S$ is absent from the statement, and not by oversight. In the
rules of Chapter~\ref{ch:lts} no transition changes $S$; the component is
carried for the extensions in which transformations read and write mutable
resources, receive external input, or behave nondeterministically. In
those extensions $S$ is reconstructible exactly when every state-affecting
input is itself recorded, and not otherwise. The present theorem claims
only the part that holds unconditionally.

Replay also serves as a checker. Because $\mathsf{apply}$ is partial,
$\replay(H)$ is defined only on histories whose every projection is a legal
lifecycle path. By the path lemma every history produced by execution
passes this check, so a ledger received from elsewhere can be audited by
replaying it: success certifies that each commitment's records form an
admissible path, with valid witnesses, matching data, and at most one
terminal record, provided Hypothesis~\ref{hyp:stable} holds at the auditor.

\section{Consistency as an invariant}

The replay theorem is stated for executions from a ledger-consistent
start. It is cleaner to state its content as an invariant on single
configurations, since that is what a running or restored system needs.

\begin{definition}[Consistency]
A configuration $C=\langle P,S,A,O,H\rangle$ is \emph{consistent},
$\mathsf{Consistent}(C)$, when $\replay(H)=\langle A,O\rangle$.
\end{definition}

\begin{proposition}[Consistency is established and preserved]
\label{prop:consistent}
The empty configuration is consistent. If $\mathsf{Consistent}(C)$ and
$C\xrightarrow{\alpha}C'$, then $\mathsf{Consistent}(C')$.
\end{proposition}

\begin{proof}
The first statement is immediate. For the second, let $\rho$ be the
appended record with subject $s$. By consistency $A(s)$ is the replayed
state of $s$, the rule's premises are exactly the conditions under which
$\delta(A(s),\rho)$ is defined, and its conclusion writes that value. The
map is unchanged elsewhere (Theorem~\ref{thm:identity}), and the
observation set changes only on collapse, by the triple replay adds.
\end{proof}

Theorem~\ref{thm:replay} is the iterate of this proposition. A system that
starts empty is consistent forever. A system resumed from stored state is
admitted only through a checked initialization, which combines the
consistency test with the completeness anchor of the next section:
\[
\frac{\replay(H)=\langle A,O\rangle\qquad
\mathsf{ValidAnchor}(H,\chi)}
{\mathsf{init}(A,O,H,\chi)\;\mathsf{ok}}.
\]
One invariant then covers empty systems, restored systems, and embedded
long-running systems alike, and the ledger path lemma applies to all of
them.

\section{Completeness and anchors}

Replay as a checker certifies that each commitment's records form an
admissible path. It cannot certify that the ledger is complete.

\begin{theorem}[Prefix indistinguishability]\label{thm:prefix}
If $H$ is the history of an execution from a ledger-consistent
configuration, then every prefix $H'\sqsubseteq H$ is itself the history
of an execution from the same configuration, and $\replay(H')$ is defined.
Consequently no predicate on histories alone separates complete ledgers
from truncated ones.
\end{theorem}

\begin{proof}
Stop the execution after its first $|H'|$ transitions. By history
monotonicity its history at that point is $H'$, and by
Theorem~\ref{thm:replay} its replay is defined. A predicate that accepted
$H$ as complete and rejected $H'$ as truncated would reject the complete
history of this shorter execution.
\end{proof}

Truncation is only the simplest case. The laboratory found interior
deletions that replay cannot detect either (Chapter~\ref{ch:worlds}), and
the general statement is this.

\begin{theorem}[Subhistory indistinguishability]\label{thm:subhistory}
There are histories $H$ produced by execution and proper subsequences
$H'$ of $H$, not prefixes of it, such that $\replay(H)$ and $\replay(H')$
are both defined and $H'$ is itself the history of an execution. Hence
replay validity certifies the local admissibility of the records
presented, not their completeness relative to what actually occurred.
\end{theorem}

\begin{proof}
Take $H=\Pop(c,q)\,;\Bind(a,c,b)\,;\Transform(c,f,b,v)\,;\Refuse(a,c,r)$
and delete the transformation. The remainder
$\Pop(c,q)\,;\Bind(a,c,b)\,;\Refuse(a,c,r)$ is the history of the
execution that refuses from the bound state, which the explicit refusal
rule permits. Both replay.
\end{proof}

Detection therefore needs information that is not in the history. The
usual object is a hash chain with a signed checkpoint. Fix a hash function
$h$ and let
\[
h_0=\mathsf{init},\qquad h_i=h(h_{i-1}\parallel e_i),
\]
where $e_i$ is the $i$-th ledger entry with its event identity. A
\emph{checkpoint} is a triple $\mathsf{Anchor}(n,h_n,\sigma)$ in which
$\sigma$ is the sequencer's signature on $(n,h_n)$.

\begin{proposition}[Relative completeness]\label{prop:anchor}
Suppose an auditor has retained an authentic checkpoint
$\mathsf{Anchor}(n,h_n,\sigma)$ independently of the ledger, the
signature scheme is unforgeable, and $h$ is collision resistant. If a
presented ledger $H$ has hash chain value $h_n$ at position $n$, then its
first $n$ entries match the anchored sequence in length, order, and
content, except with negligible probability. A presented ledger shorter
than $n$, or with a different chain value at $n$, is detected, whether
the discrepancy is a truncation, an interior deletion, an insertion, a
reordering, or an altered record.
\end{proposition}

\begin{proof}
The first two cases are direct checks. For the third, two different
sequences of $n$ entries hashing to the same $h_n$ yield, at the first
position where the chains diverge, a collision in $h$. The argument is
the standard one for hash chains and is used here without further
elaboration; its strength is exactly that of the cryptographic
assumptions.
\end{proof}

The completeness obtained is relative: relative to the anchor, and to
trust in whoever issued it or in the transparency mechanism that holds
it. The hypothesis that carries the weight is \emph{independent
retention}.
A sequencer that can replace both the ledger and the checkpoint can
publish a shorter history with a fresh, internally consistent checkpoint,
and nothing internal to the pair reveals it. Deployments close that gap by
placing checkpoints where the sequencer cannot rewrite them: in an
external append-only transparency log of the kind used for web
certificates \cite{rfc6962},
or distributed among independent witnesses. The lesson fits the calculus
closely. A history can vouch for the admissibility of what it contains; a
claim that nothing is missing needs a witness outside the history whose
completeness is in question.

\section{Duplicate delivery}

When records travel between replicas they may arrive more than once. Pair
each record with an \emph{event identity} $i$ assigned by the sequencer
when the record is appended, so that a ledger is a sequence of pairs
$(i,\rho)$. Extend the replay state with the set $I$ of identities already
applied, and let
\[
\mathsf{apply}(\langle A,O,I\rangle,(i,\rho))=
\begin{cases}
\langle A,O,I\rangle & \text{if }i\in I,\\
\langle \mathsf{apply}(\langle A,O\rangle,\rho),\;I\cup\{i\}\rangle
& \text{otherwise.}
\end{cases}
\]

\begin{hypothesis}[Identity coherence]\label{hyp:coherence}
Any two ledger entries with the same event identity carry the same record.
\end{hypothesis}

\begin{theorem}[Idempotent disposition]
For every state $X$ and entry $e$,
$\mathsf{apply}(\mathsf{apply}(X,e),e)=\mathsf{apply}(X,e)$. Under
Hypothesis~\ref{hyp:coherence}, if $\mathsf{dedup}(H)$ keeps the first
occurrence of each identity, then
\[
\replay(H)=\replay(\mathsf{dedup}(H))
\]
on the $\langle A,O\rangle$ components, both sides being defined or
undefined together.
\end{theorem}

\begin{proof}
The first statement is immediate, since after the first application $i\in
I$. For the second, induct on $H$. An entry whose identity is new is
applied on both sides. An entry whose identity was seen earlier is a no-op
on the left and absent on the right; by coherence it is the same record
as the earlier occurrence, so dropping it cannot change which occurrence
was the one applied.
\end{proof}

The mechanism is deliberately simple; any system with an applied-set
achieves idempotence. The content of the theorem is in its hypothesis.
Exactly-once \emph{disposition} is obtained from at-least-once
\emph{delivery} precisely when the sequencer never reuses an identity for a
different record, and that is an obligation on the arbiter, not a property
the replay procedure can establish for itself.

\section{Which reorderings replay ignores}

Duplicate tolerance concerns the same record arriving twice. A different
question is whether distinct records may arrive in a different order.

\begin{definition}[Independence and dependence order]
\label{def:dependence}
Records $\rho$ and $\sigma$ are \emph{independent}, written $\rho\,I\,\sigma$,
when $\subj(\rho)\neq\subj(\sigma)$ and neither is a compensation record
whose predecessor is the other's subject. Two histories are \emph{trace
equivalent}, $H\sim_I H'$, when one can be obtained from the other by
finitely many exchanges of adjacent independent records. The
\emph{dependence order} $\prec_H$ is the transitive closure of the relation
``$\rho$ precedes $\sigma$ in $H$ and $\rho$, $\sigma$ are not
independent''.
\end{definition}

\begin{proposition}[Replay respects commutation]\label{prop:commute}
If $H\sim_I H'$ then $\replay(H)$ is defined if and only if $\replay(H')$
is, and when defined they are equal.
\end{proposition}

\begin{proof}
It suffices to treat one exchange, $H=H_1\cdot\rho\cdot\sigma\cdot H_2$ and
$H'=H_1\cdot\sigma\cdot\rho\cdot H_2$ with $\rho\,I\,\sigma$. From the
state reached after $H_1$, applying $\rho$ reads and writes the map at
$\subj(\rho)$, and if $\rho$ is a compensation record, also reads it at the
predecessor. By independence, $\sigma$ writes neither of these keys, and
symmetrically. So each application is defined, and produces the same
update, whether it comes first or second. The observation components are
unions of sets and commute. The remainder $H_2$ is then applied to equal
states.
\end{proof}

\begin{corollary}
The disposition and observation state recoverable from a ledger is a
function of its Mazurkiewicz trace $[H]_{\sim_I}$. Every ancestry chain of
Chapters~\ref{ch:rbe} and \ref{ch:collapse} is a $\prec_H$-chain.
\end{corollary}

\begin{proof}
The first sentence is Proposition~\ref{prop:commute}. For the second,
consecutive records of one projection share a subject and are therefore
dependent.
\end{proof}

This result has two consequences for the rest of the book. It makes
precise the spine's phrase ``ledger order where order represents a true
dependency'': the true dependencies are those of $\prec_H$, and the total
order imposed by the sequencer carries more information than replay can
use. And it bears on the first comparative claim, since a trace with an
independence relation determined by subjects is very close to the
configuration structure of an event structure. Chapter~\ref{ch:eventstructures} takes that
comparison up directly, and Chapter~\ref{ch:sequencer} asks what the sequencer's extra
information is for, given that replay does not need it.

\chapter{Petri Nets, Traces, and Event Structures}
\label{ch:eventstructures}
\plan{Occurrence nets, Mazurkiewicz traces, and prime event structures with
causality and conflict. Starting from Proposition~\ref{prop:commute},
constructs an event structure from a ledger: events are records,
causality is $\prec_H$, and conflict is the relation between two terminal
records for one subject, which Corollary~\ref{cor:terminal-history}
already shows can never both occur. Tests the first comparative claim and
isolates the residue: whether refusal is exactly conflict, and what
role the reasons and witnesses play, which an event structure does not
carry.}

\chapter{Logical Clocks, CRDTs, and the Single Sequencer}
\label{ch:sequencer}
\plan{Relates happened-before and its linear extensions to the arbiter's
sequence. Given that replay depends only on the trace, asks what the total
order is needed for: event identity assignment
(Hypothesis~\ref{hyp:coherence}), freshness of commitment identities
across replicas, and the arbitration of competing compensations. Asks which
disposition programs admit a conflict-free replicated implementation.}

\chapter{Eventual Disposition}
\label{ch:progress}

Progress (Theorem~\ref{thm:progress}) is a statement about one
configuration: it is never silently stuck. Eventual disposition is a
statement about whole executions: every commitment is eventually refused
or collapsed. The second does not follow from the first. An execution can
make progress forever while neglecting a particular commitment, a
transformation can run forever, a verifier can fail to decide, and an
outside party can simply never answer. The theorem of this chapter holds
only under hypotheses that exclude each of these, and they are stated as
named predicates on executions so that it cannot be read as stronger than
it is.

\section{Hypotheses on executions}

An \emph{execution} $\xi$ is a finite or infinite sequence of transitions
$C_0\to C_1\to\cdots$; if finite, it is \emph{maximal} when its last
configuration has no transition, discharges included. A \emph{redex} is an
occurrence of a top-level prefix, tracked from one configuration to the
next while the steps taken are elsewhere.
\begin{itemize}[leftmargin=1.5em]
\item $\mathsf{Fair}(\xi)$: no redex is enabled at every configuration
from some point on without ever being taken. This is weak fairness, a
condition on the scheduler; the calculus itself does not schedule.
\item $\mathsf{AllTransformsTerminate}(\xi)$: every transformation redex
on a bound commitment in $\xi$ evaluates to a value.
\item $\mathsf{AllVerifiersDecide}(\xi)$: every verification redex on a
transformed commitment in $\xi$ returns $\mathsf{ok}$ or an authenticated
$\mathsf{fail}$.
\item $\mathsf{AllObligationsResponsive}(\xi)$: every $\mathsf{await}$
redex in $\xi$ is eventually settled: discharged, or, for the deadline
form, discharged or certified as timed out.
\end{itemize}
The second and third are quantified over the computations the execution
actually invokes, not over the lambda layer as a whole. The lambda layer
may be Turing complete; the theorem asks only that the particular
transformations and verifiers this execution calls on terminate.

\section{The theorem}

\begin{theorem}[Eventual disposition]\label{thm:eventual}
Let $\Sigma;\Gamma;\Delta\vdash C_0\;\mathsf{ok}$ and let $\xi$ be an
execution from $C_0$ that is maximal or infinite and satisfies
$\mathsf{Fair}$, $\mathsf{AllTransformsTerminate}$,
$\mathsf{AllVerifiersDecide}$, and $\mathsf{AllObligationsResponsive}$.
Then every commitment that is live at some configuration of $\xi$ is
terminal at some later configuration of $\xi$.
\end{theorem}

\begin{proof}
Let $c$ be live at $C_i$ and suppose, for contradiction, that it is never
terminal afterward. By Theorem~\ref{thm:sr}, every $C_j$ with $j\ge i$ is
well typed and its linear context contains exactly one right for $c$.

\emph{Ownership.} Flatten the process of $C_j$ into its top-level
parallel components, looking through restriction. By \textsc{T-Par} and
Lemma~\ref{lem:own}, exactly one component's context contains the right
for $c$; call it the \emph{owner} $O_j$. It is not a replication, since
\textsc{T-Rep} requires an empty context. By \textsc{T-Nil} it is not
$\mathbf 0$. So $O_j$ is a prefix or a sum of prefixes.

\emph{Persistence.} If a prefix of $O_j$ is enabled at $C_j$, it remains
enabled until $O_j$ acts. Its premises concern the states of commitments
whose rights $O_j$ holds, which no other component can change, since
changing a state requires holding its right; authority facts, which are
fixed; terminal facts, which are monotone; freshness, which is always
satisfiable; and the results of evaluation and verification, which are
deterministic. If the prefix is an $\mathsf{await}$, responsiveness
supplies its settlement, by discharge or timeout.

\emph{Enabledness.} By Theorem~\ref{thm:progress} and its proof, every
non-await prefix of a well-typed configuration is enabled, given that the
computation it invokes terminates, which the second and third hypotheses
provide.

\emph{Descent.} Let $\|P\|$ count the prefix occurrences of $P$ that are
not beneath a replication. Each step taken by the owner consumes one
prefix, and possibly discards the other summands of a sum, and the right
for $c$ passes to a subterm of what remains. Hence the size of the owner
of $c$ strictly decreases with each owner step, while steps of other
components leave the owner unchanged.

If $\xi$ is infinite, fairness and persistence force the owner to act
infinitely often, so its size would decrease without bound, which is
impossible. If $\xi$ is finite and maximal, its last configuration has no
transition; by responsiveness it is not waiting, since a pending
settlement would be a transition; so by Theorem~\ref{thm:progress} it is
inert and every commitment is terminal. Either way, $c$ is eventually
terminal.
\end{proof}

The proof uses the type system at every step. Ownership is linearity.
Persistence depends on ownership: no other component can disable the
owner's prefix. Descent depends on \textsc{T-Rep}: rights are never
replicated, so the owner is a finite term. And the finite case depends on
\textsc{T-Nil}: a program cannot end holding a live right.

\section{What fails without each hypothesis}

Each hypothesis is needed, and each has a smallest counterexecution: a
well-typed configuration and an execution satisfying every other
hypothesis in which some commitment stays live forever. In each case $c$
has just been opened by a principal $a$ with the authority the program
uses.
\begin{enumerate}[leftmargin=1.5em]
\item \emph{Unfair scheduling.}
\[
!\,\mathsf{pop}\;q'\;\mathsf{as}\;d.\,\mathsf{refuse}_a\;d\;\mathsf{because}\;r.\mathbf 0
\;\mid\;\mathsf{refuse}_a\;c\;\mathsf{because}\;r.\mathbf 0 .
\]
A scheduler that always chooses the replicated component takes a step
forever and never the refusal of $c$, although it is enabled throughout.
Every commitment the replication creates is disposed of at once; $c$
never moves.
\item \emph{Divergent transformation.} With $A(c)=\mathsf{bound}(q,b)$ and
$f$ a fixed point that diverges on $b$,
\[
\mathsf{transform}\;c\;\mathsf{with}\;f.\,\mathsf{refuse}_a\;c\;\mathsf{because}\;r.\mathbf 0 .
\]
The only redex never completes its evaluation premise, and $c$ stays
bound.
\item \emph{Nondeciding verifier.} With $A(c)=\mathsf{transformed}(b,v)$
and a verifier $g$ that on $(c,v)$ neither returns a witness nor issues a
certificate,
\[
\mathsf{verify}_a\;c\;\mathsf{using}\;g\;\mathsf{then}\;
\mathsf{refuse}_a\;c\;\mathsf{because}\;r.\mathbf 0\;\mathsf{else}\;\mathbf 0 .
\]
Stability (Hypothesis~\ref{hyp:stable}) is not enough; the verifier must
decide.
\item \emph{Permanently unanswered obligation.}
\[
\mathsf{await}_e\;c\;\omega.\,\mathsf{refuse}_a\;c\;\mathsf{because}\;r.\mathbf 0,
\]
where $e$ never discharges $\omega$. The configuration is well typed and
waiting, not stuck in the sense of Theorem~\ref{thm:progress}, and the
ledger correctly shows nothing. No other component can refuse $c$, since
the awaiting process holds the right.
\item \emph{An environment that neither discharges nor times out.}
\[
\mathsf{await}_e\;c\;\omega.\,P\;\mathsf{within}_b\;t\;\mathsf{else}\;
\mathsf{refuse}_a\;c\;\mathsf{because}\;\mathsf{expired}.\mathbf 0,
\]
where $e$ never discharges and the clock authority $b$ never certifies
expiry. Declaring a deadline does not by itself end the wait; the
deadline is only as responsive as the authority that attests to it.
\item \emph{An unattended verification failure.} With
$A(c)=\mathsf{transformed}(b,v)$ and a verifier that returns a certified
negative result,
\[
\mathsf{verify}_a\;c\;\mathsf{using}\;g\;\mathsf{then}\;P\;\mathsf{else}\;
\mathsf{await}_h\;c\;\omega_{\mathrm{review}}.\,
\mathsf{refuse}_h\;c\;\mathsf{because}\;r.\mathbf 0,
\]
where the disposition authority $h$ never completes its review. The
commitment sits in $\mathsf{verificationFailed}$ forever. This is waiting,
not a type error, and it exposes the institutional dependency the
failed-verification state was introduced to make visible: a negative
result that no authorized party acts on.
\end{enumerate}
The last example does not need a hypothesis of its own. In a well-typed
program the continuation after a failed verification holds the right, and
\textsc{T-Nil} forbids it to stop holding it, so it must either refuse
directly, in which case fairness makes the refusal happen, or wait on an
obligation, in which case responsiveness does. What some presentations
would state separately as ``failed-verification obligations receive an
authorized response'' is, here, the instance of
$\mathsf{AllObligationsResponsive}$ for the obligation the failed branch
awaits.
A sixth example shows that typing is needed as well: the untyped process
$\mathsf{pop}\;q\;\mathsf{as}\;c.\mathbf 0$ satisfies all four hypotheses
and leaves $c$ open forever. \textsc{T-Nil} rejects it.

The last two examples are where the liveness question meets institutions.
A commitment waiting on a party that never answers is not failing in any
sense the calculus can detect; the ledger shows it waiting, correctly. A
deadline turns silence into a record, $\mathsf{TimedOut}$, that says the
obligation went unanswered within the declared interval. It does not turn
silence into refusal. What follows a timeout, whether further waiting,
compensation, or refusal, is decided by an authorized rule in the
$\mathsf{else}$ branch, and choosing that rule is a policy the calculus
can express but does not make.

%===================================================================

\chapter{Bisimulation and the Spectrum as Collapse Rules}
\label{ch:spectrum}
\plan{Recasts bisimilarity, and then trace, failures, testing, and ready
equivalences, as collapse rules over histories, and compares the result
with van Glabbeek's spectrum. Settles the third comparative claim, using
the rule-indexed collapse records of Chapter~\ref{ch:lts} as the test
case.}

\chapter{Checking the Layer Translation}
\label{ch:layers-back}
\plan{With forward simulation and backward adequacy proved at the level
of ledgers in Chapter~\ref{ch:layers}, this chapter checks the translation
against Gorla's encodability criteria, where compositionality and name
invariance are expected to hold and divergence reflection turns on the
reading of transformation. It then attacks the transparent-reading
conjecture: that the kernel's encoding of application agrees with
evaluation in the lambda layer. A failure there would not undo the ledger
theorems; it would locate precisely the computational content that the
kernel does not itself supply.}

\chapter{Typed Disposition Safety}
\label{ch:typedsafety}

All the pieces are in place. This chapter states the book's central safety
theorem and proves it by a short induction, as the proof spine promised,
because each clause has already been established by a local result.

\begin{theorem}[Typed disposition safety]\label{thm:tds}
Let $\Sigma_0;\Gamma_0;\Delta_0\vdash C_0\;\mathsf{ok}$, and let
$C_0\to C_1\to\cdots\to C_n$ be any finite execution. Then there are
$\Sigma_n$, $\Gamma_n$, $\Delta_n$ such that:
\begin{enumerate}[label=(\arabic*)]
\item $\Sigma_n;\Gamma_n;\Delta_n\vdash C_n\;\mathsf{ok}$;
\item every live commitment of $C_n$ has exactly one right in $\Delta_n$,
and every terminal commitment has none;
\item every attributed transition of the execution was authorized;
\item every verified state and every $\Verify$ record carries a witness
valid for its own commitment and value;
\item every collapsed commitment has a complete ledger path through
opening, binding, transformation, verification, and collapse;
\item every refusal appended during the execution is present in $H_n$;
\item no commitment has two terminal records in $H_n$;
\item $\replay(H_n)=\langle A_n,O_n\rangle$;
\item the process of $C_n$ contains no well-typed continuation acting on a
commitment whose right its context does not hold, and at most one
component holds each right.
\end{enumerate}
\end{theorem}

\begin{proof}
By induction on $n$, the case $n=0$ being the hypothesis. For the step,
Theorem~\ref{thm:sr} carries (1) from $C_n$ to $C_{n+1}$. The remaining
clauses are properties of any well-typed, consistent configuration and
of the execution leading to it:
\begin{itemize}[leftmargin=1.5em]
\item (2) is condition (C3) of (1), that is, the corollary to
Theorem~\ref{thm:sr};
\item (3) is Theorem~\ref{thm:auth}, applied to each step;
\item (4) is Corollary~\ref{cor:witness} with Lemma~\ref{lem:canon}, and
non-transferability under Hypothesis~\ref{hyp:index};
\item (5) is Theorem~\ref{thm:collapse};
\item (6) is Theorem~\ref{thm:refusal} with history monotonicity;
\item (7) is Corollary~\ref{cor:terminal-history};
\item (8) is condition (C2) of (1), maintained by
Proposition~\ref{prop:consistent};
\item (9) is Lemma~\ref{lem:inv}, since each prefix acting on $c$ requires
$c$'s right in its context, together with Lemma~\ref{lem:own}, since two
components cannot both hold it.
\end{itemize}
The untyped results used here were proved for executions from a
consistent configuration, which (1) guarantees.
\end{proof}

Two things are deliberately absent from the theorem. It says nothing about
whether commitments are eventually disposed of; that is
Theorem~\ref{thm:eventual}, and it needs fairness, terminating
computations, deciding verifiers, and responsive obligations, none of
which a safety theorem should assume. And it says nothing about whether
the ledger is complete; that is Proposition~\ref{prop:anchor}, and it
needs a checkpoint retained outside the history.

The three validation levels of the Preface can now be seen together.
The laboratory's bounded propositions check the untyped clauses (5), (6),
and (7) and the replay clause (8) on every ledger of a finite class, and
found the specification errors the book corrected along the way. The
untyped proofs establish (3) through (8) for every execution of every
program. The typed metatheory adds (1), (2), and (9), which the untyped
calculus cannot even state: that rights are neither duplicated nor
abandoned, and that no well-typed program contains a competitor for a
commitment it does not own.

\chapter{A Map of Positions}
\plan{One table collecting every comparison: calculus, direction of
encoding, criteria met, outcome, and grade of assertion. Answers the three
comparative claims in their final form.}

%===================================================================
\appendix
%===================================================================

\chapter{Proof Spine}

The proof dependency is: unique abstract parsing and alpha-invariance;
fresh commitment introduction and lifecycle preservation; append-only
history; value and process substitution; subject reduction; authority
preservation and binding integrity; witness ancestry; collapse soundness
and terminal exclusivity; replay correspondence; compensation
conservation; and finally typed disposition safety. Liveness is stated
separately because it requires fairness, responsive external dependencies,
terminating transformations, and decidable verification. Safety holds for
every finite execution without those assumptions.

\begin{center}
\small
\begin{tabular}{lll}
\toprule
Result & Status & Location\\
\midrule
Unique tokenization & proved & Lem.~\ref{lem:lex}\\
Unique parsing & proved (LL(1)) & Thm.~\ref{thm:parse}\\
Printer adequacy & proved & Thm.~\ref{thm:adequacy}\\
Equivariance & proved & Ch.~\ref{ch:syntax}\\
Commitment identity & proved & Thm.~\ref{thm:identity}\\
Lifecycle and history monotonicity & proved & Ch.~\ref{ch:lts}\\
Authorized transition & proved (static policy) & Thm.~\ref{thm:auth}\\
Source authenticity of verification failures & proved & Prop.~\ref{prop:vf}\\
Ledger acceptance of verification failures & proved (given certificates) & Prop.~\ref{prop:vf-ledger}\\
Diamond for independent steps & proved & Prop.~\ref{prop:diamond}\\
Ledger path lemma & proved & Lem.~\ref{lem:path}\\
Refusal preservation & proved & Thm.~\ref{thm:refusal}\\
Binding integrity & proved (exclusivity: typed) & Ch.~\ref{ch:rbe}\\
Witness ancestry & proved & Cor.~\ref{cor:witness}\\
Collapse soundness & proved & Thm.~\ref{thm:collapse}\\
Terminal exclusivity & proved & Cor.~\ref{cor:terminal-history}\\
Compensating continuation & proved & Thm.~\ref{thm:compensation}\\
Compensation ancestry & proved & Lem.~\ref{lem:companc}\\
Consistency invariant & proved & Prop.~\ref{prop:consistent}\\
Replay correspondence & proved & Thm.~\ref{thm:replay}\\
Prefix indistinguishability & proved & Thm.~\ref{thm:prefix}\\
Subhistory indistinguishability & proved & Thm.~\ref{thm:subhistory}\\
Relative completeness & proved (conditional) & Prop.~\ref{prop:anchor}\\
Idempotent disposition & proved (under coherence) & Ch.~\ref{ch:replay}\\
Replay respects commutation & proved & Prop.~\ref{prop:commute}\\
Causal shuffle & proved & Thm.~\ref{thm:shuffle}\\
Modal terminal exclusivity & proved & Prop.~\ref{prop:modal-terminal}\\
No cycles, one re-entry & proved & Prop.~\ref{prop:reentry}\\
Kernel translation: forward simulation & proved (ledger level) & Thm.~\ref{thm:fwd}\\
Kernel translation: backward adequacy & proved (ledger, universal process) & Thm.~\ref{thm:bwd}\\
Kernel translation: transparent reading & conjecture & Ch.~\ref{ch:layers-back}\\
Value substitution, linear and unrestricted & proved & Ch.~\ref{ch:types}\\
Canonical forms for rights & proved & Lem.~\ref{lem:canon}\\
Preservation & proved & Thm.~\ref{thm:sr}\\
Progress with typed waiting & proved & Thm.~\ref{thm:progress}\\
Protocol fidelity (single commitment) & proved & Thm.~\ref{thm:fidelity}\\
Typed disposition safety & proved & Thm.~\ref{thm:tds}\\
Eventual disposition & proved (four hypotheses) & Thm.~\ref{thm:eventual}\\
Non-derivability of authority & proved (static policy) & Prop.~\ref{prop:nonderiv}\\
Delegation authenticity & proved & Prop.~\ref{prop:delauth}\\
Conservation of exclusive capabilities & proved & Prop.~\ref{prop:conserve}\\
Preservation under transfer & proved & Prop.~\ref{prop:delpres}\\
Delegation chain soundness & proved & Prop.~\ref{prop:chain}\\
\bottomrule
\end{tabular}
\end{center}

\chapter{Complete Syntax, Typing, and Transition Rules}
\plan{The full grammar of both layers, all typing rules, all labelled
transition rules, and the replay clause paired with each.}

\chapter{Encodability Criteria, Stated Formally}
\plan{Formal statements of the criteria of Chapter~\ref{ch:encodings},
with notation fixed for the whole book.}

\chapter{A Worked Encoding: A CCS Fragment into the Calculus}
\label{app:worked}
\plan{A full translation of a small CCS fragment with each criterion
checked explicitly.}

\chapter{Open Problems}
\label{app:open}
\begin{enumerate}[leftmargin=1.5em]
\item \emph{The transparent reading.} The layer translation is proved
under the opaque reading, where evaluation and verification results are
guard data. Whether the kernel's encoding of application agrees with
evaluation in the lambda layer, so that the kernel itself performs the
computation, is open (Chapter~\ref{ch:layers-back}).
\item \emph{Verifier versioning.} The theorems of Part~V assume
Hypothesis~\ref{hyp:stable}. With versioned verifiers, the version must be
recorded and witness ancestry restated relative to it; whether an old
verification remains evidence under a new version is a policy question the
calculus should make explicit rather than decide.
\item \emph{Competing compensations.} Two authorized principals may
compensate one terminal commitment. Should the calculus forbid this,
record both and require a later arbitration, or leave it to a collapse
rule?
\item \emph{Recording the road not taken.} The choice rule leaves no
record of the discarded summand. A variant that records it may or may not
be more expressive (Chapter~\ref{ch:choice}).
\item \emph{Restriction and visibility.} Restriction scopes naming but not
visibility. Whether a scoped ledger is desirable, and what it would do to
replay, is open (Chapter~\ref{ch:pi}).
\item \emph{Mutable store.} Replay correspondence excludes $S$. The exact
recording obligations under which $S$ becomes reconstructible remain to
be stated.
\item \emph{Denied attempts.} An unauthorized attempt has no transition and
leaves no record. Recording $\mathsf{DeniedAttempt}$ on the annotation
channel would make attempted overreach auditable without making it causal;
who records it, and whether an untrusted component can flood that channel,
is open.
\item \emph{Anchoring the ledger's end.} Replay cannot detect removal of a
ledger's final records (Chapter~\ref{ch:worlds}). The sequencer must
publish an anchor; which anchor suffices under which failure model is
open.
\item \emph{The purpose of total order.} Replay depends only on the trace
(Proposition~\ref{prop:commute}). Which guarantees, if any, require the
sequencer's total order beyond identity assignment and freshness
(Chapter~\ref{ch:sequencer})?
\end{enumerate}

%===================================================================
\backmatter
\begin{thebibliography}{99}
\bibitem{agha} G.~Agha. \emph{Actors: A Model of Concurrent Computation in Distributed Systems}. MIT Press, 1986.
\bibitem{brookes} S.~D.~Brookes, C.~A.~R.~Hoare, and A.~W.~Roscoe. A theory of communicating sequential processes. \emph{Journal of the ACM} 31(3), 1984.
\bibitem{church} A.~Church. \emph{The Calculi of Lambda-Conversion}. Princeton University Press, 1941.
\bibitem{cousot} P.~Cousot and R.~Cousot. Abstract interpretation: a unified lattice model for static analysis of programs. \emph{POPL}, 1977.
\bibitem{girard} J.-Y.~Girard. Linear logic. \emph{Theoretical Computer Science} 50, 1987.
\bibitem{gorla} D.~Gorla. Towards a unified approach to encodability and separation results for process calculi. \emph{Information and Computation} 208(9), 2010.
\bibitem{hewitt} C.~Hewitt, P.~Bishop, and R.~Steiger. A universal modular ACTOR formalism for artificial intelligence. \emph{IJCAI}, 1973.
\bibitem{hoare} C.~A.~R.~Hoare. \emph{Communicating Sequential Processes}. Prentice Hall, 1985.
\bibitem{honda} K.~Honda. Types for dyadic interaction. \emph{CONCUR}, 1993.
\bibitem{lamport} L.~Lamport. Time, clocks, and the ordering of events in a distributed system. \emph{Communications of the ACM} 21(7), 1978.
\bibitem{mazurkiewicz} A.~Mazurkiewicz. Concurrent program schemes and their interpretations. DAIMI Report PB-78, Aarhus University, 1977.
\bibitem{milner} R.~Milner. \emph{A Calculus of Communicating Systems}. LNCS 92, Springer, 1980.
\bibitem{mpw} R.~Milner, J.~Parrow, and D.~Walker. A calculus of mobile processes, I and II. \emph{Information and Computation} 100(1), 1992.
\bibitem{npw} M.~Nielsen, G.~Plotkin, and G.~Winskel. Petri nets, event structures and domains, part I. \emph{Theoretical Computer Science} 13, 1981.
\bibitem{palamidessi} C.~Palamidessi. Comparing the expressive power of the synchronous and asynchronous $\pi$-calculi. \emph{Mathematical Structures in Computer Science} 13(5), 2003.
\bibitem{park} D.~Park. Concurrency and automata on infinite sequences. \emph{Theoretical Computer Science}, LNCS 104, 1981.
\bibitem{petri} C.~A.~Petri. \emph{Kommunikation mit Automaten}. Dissertation, University of Bonn, 1962.
\bibitem{plotkin} G.~D.~Plotkin. A structural approach to operational semantics. DAIMI FN-19, Aarhus University, 1981.
\bibitem{nasrollahi} F.~S.~Fatemi~Nasrollahi, E.~Newby, and R.~Albert. Modularity-based approach identifies small sets of control nodes in synthetic and biological Boolean networks. \emph{Scientific Reports} 16:29905, 2026. doi:10.1038/s41598-026-48763-1.
\bibitem{ding} L.~Ding, J.~Thomson, J.~Taylor, and C.~Do. LLMs can assist with proposal selection at large user facilities. \emph{Scientific Reports} 16:29879, 2026. doi:10.1038/s41598-026-60226-1.
\bibitem{rfc6962} B.~Laurie, A.~Langley, and E.~Kasper. Certificate Transparency. RFC 6962, 2013.
\bibitem{tait} W.~W.~Tait. Intensional interpretations of functionals of finite type I. \emph{Journal of Symbolic Logic} 32(2), 1967.
\bibitem{pierce} B.~C.~Pierce. \emph{Types and Programming Languages}. MIT Press, 2002.
\bibitem{shapiro} M.~Shapiro, N.~Pregui\c{c}a, C.~Baquero, and M.~Zawirski. Conflict-free replicated data types. \emph{SSS}, 2011.
\bibitem{vanglabbeek} R.~J.~van Glabbeek. The linear time--branching time spectrum I. In \emph{Handbook of Process Algebra}, Elsevier, 2001.
\bibitem{wadler} P.~Wadler. Propositions as sessions. \emph{ICFP}, 2012.
\bibitem{winskel} G.~Winskel. Event structures. In \emph{Advances in Petri Nets}, LNCS 255, Springer, 1987.
\bibitem{flyxion-spec} Flyxion. Spherepop specification. Unpublished manuscript, 2025--2026.
\bibitem{flyxion-binding} Flyxion. The algebra of binding. Unpublished manuscript, 2026.
\bibitem{flyxion-witness} Flyxion. The witness plane: admissibility, verification, and commitment in event systems. Unpublished manuscript, 2026.
\end{thebibliography}

\end{document}
